%============================================================
%  Bd. 08 - Bijektive Funktionen
%============================================================

\documentclass[main.tex]{subfiles}

\ifSubfilesClassLoaded{
  \BandSetupStandalone{B08}
}{
}

\title{Bd. 08 - Bijektive Funktionen}
\author{Martin Kunze}
\date{}
\setcounter{file}{8}

\hypersetup{
  pdftitle={Bd. 08 - Bijektive Funktionen},
  pdfauthor={Martin Kunze}
}

\begin{document}

\title{Bd. 08 - Bijektive Funktionen}
\author{Martin Kunze}
\date{}
\hypersetup{pageanchor=false}
\maketitle
\hypersetup{pageanchor=true}
\setcounter{tocdepth}{3}
\tableofcontents
\setcounter{file}{8}
\renewcommand{\FormulaBandID}{8}

\chapter{Einleitung}

Dieser Band verbindet die in Band~06 und Band~07 aufgebauten Begriffe der
Injektivität und Surjektivität zur Bijektivität. Der Aufbau beginnt wie in
Band~03 mit der zusammenfassenden Begriffsdefinition und entwickelt danach
die charakteristischen Konstruktionen: Identität, Umkehrfunktion,
Komposition und Erweiterung.

\chapter{Axiomatischer Aufbau der bijektiven Funktionen}

\section{Begriff und Grundsätze}

\FormulaDefDeltaK[Begriff der bijektiven Funktion]{F\colon A\bij B}{Bijektive Funktion}{
  \DeltaRow{Mengen}{A\dsep B\dsep x\dsep y}
  \DeltaPrem{Funktionen}{F\colon A\to B}[\FormulaRefAuto{Funktion}]
%
  \DeltaRow{\textbf{Axiome}}{}
  %
  % — Injektivität —
  \DeltaRow{Injektivität}
           {F\colon A\inj B\dsep x,y\in A\dsep F(x)=F(y)\vdash x=y}
           [\FormulaRefAuto{F\colon A\inj B\dsep x\in A\dsep y\in A\dsep F(x)=F(y)\vdash x=y}]
  %
  % — Surjekvität —
  \DeltaRow{Surjektivität}
           {y\in B\vdash \exists x\in A\; F(x)=y}
           [\FormulaRefAuto{y\in B\vdash \exists x\in A\; F(x)=y}]
  \DeltaRow{\textbf{Neue Symbole}}{}
  \DeltaPrem{Bijektive Funktionen}{ }
}

\FormulaThmDeltaK[Subbeweiseinsetzung für bijektive Funktionen]{%
  [F\colon A\to B]\vdots P\dsep F\colon A\bij B\vdash P%
}{BijectiveFunctionSubproofElimination}{%
  \DeltaRow{Mengen}{A\dsep B}
  \DeltaRow{Aussagevariablen}{P}
  \DeltaRow{Funktionen}{F}
}
\begin{tabproof}
  \proofstep{1}{F\colon A\bij B}{\rA}
  \proofstep{1}{F\colon A\to B}{%
    \FormulaRefAuto{Bijektive Funktion}[def]{1}}
  \proofstep{1}{P}{[F\colon A\to B]\vdots P(2)}
\end{tabproof}

\FormulaThmDelta{F\colon A\bij B\dsep x\in A\vdash F(x)\in B}{
\DeltaRow{Mengen}{A\dsep B}
}
\begin{tabproof}
  \proofstep{1}{F\colon A\bij B}{\rA}
  \proofstep{2}{x\in A}{\rA}
  \proofstep{1}{F\colon A\to B}{\FormulaRefAuto{Bijektive Funktion}{1}}
  \proofstep{1,2}{F(x)\in B}{%
    \FormulaRefAuto{F\colon A\to B\dsep x\in A\vdash F(x)\in B}{3,2}}
\end{tabproof}

\FormulaThmDelta[Elemente einer Teilmenge liegen im Bild der Teilmenge unter bijektiven Funktionen]{%
  C\subseteq A\dsep F\colon A\bij B\dsep x\in C\vdash F(x)\in F[C]
}{
\DeltaRow{Mengen}{A\dsep B\dsep C\dsep x}
\DeltaRow{Funktionen}{F}
}
\begin{tabproof}
  \proofstep{1}{C\subseteq A}{\rA}
  \proofstep{2}{F\colon A\bij B}{\rA}
  \proofstep{3}{x\in C}{\rA}
  \proofstep{2}{F \colon A \to B}{%
      \FormulaRefAuto{Bijektive Funktion}{2}}
  \proofstep{1,2,3}{F(x)\in F[C]}{%
    \FormulaRefAuto{C\subseteq A\dsep F\colon A\to B\dsep x\in C\vdash F(x)\in F[C]}{%
      1,4,3}}
\end{tabproof}

\FormulaThmDelta{F\colon A\inj B\dsep F\colon A\sur B \vdash F\colon A\bij B}{
\DeltaRow{Mengen}{A\dsep B}
\DeltaPrem{Funktionen}{F\colon A\to B}
}
\begin{tabproofsplit}
\proofpart{Injektivität}
\proofstep{1}{F\colon A\inj B}{\rA}
    \proofstep{2}{F\colon A\sur B}{\rA}
    \proofstep{1}{\forall x,y\in A\,(F(x)=F(y)\rightarrow x=y)}{%
      \FormulaRefAuto{F\colon A\inj B\dsep x\in A\dsep y\in A\dsep F(x)=F(y)\vdash x=y}{1}}

\closeproofpart

\proofpart{Surjektivität}
  \setcounter{proofstepnr}{3}% Fortlaufende Schritte dieses Beweises.
\proofstep{1}{\forall y\in B\,\exists x\in A\,(F(x)=y)}{\FormulaRefAuto{y\in B\vdash \exists x\in A\; F(x)=y}{2}}

\closeproofpart

\proofpart{Bijektivität}
  \setcounter{proofstepnr}{4}% Fortlaufende Schritte dieses Beweises.
\proofstep{1,2}{F\colon A\bij B}{\FormulaRefAuto{Bijektive Funktion}{3,4}}

\closeproofpart
\end{tabproofsplit}

\FormulaThmDelta{F\colon A\bij B \vdash F\colon A\inj B}{
\DeltaRow{Mengen}{A\dsep B}
}
\begin{tabproof}
    \proofstep{1}{F\colon A\bij B}{\rA}
    \proofstep{1}{F\colon A\to B}{\FormulaRefAuto{Bijektive Funktion}{1}}
    \proofstep{1}{\forall x,y\in A\, (F(x)=F(y)\rightarrow x=y)}{\FormulaRefAuto{Bijektive Funktion}{1}}
    \proofstep{1}{F\colon A\inj B}{\FormulaRefAuto{Injektive Funktion}{2,3}}
\end{tabproof}

\FormulaThmDelta{F\colon A\bij B \vdash F\colon A\sur B}{
\DeltaRow{Mengen}{A\dsep B}
}
\begin{tabproof}
    \proofstep{1}{F\colon A\bij B}{\rA}
    \proofstep{1}{F\colon A\to B}{\FormulaRefAuto{Bijektive Funktion}{1}}
    \proofstep{1}{\forall y\in B\exists x\in A\; F(x)=y}{\FormulaRefAuto{Bijektive Funktion}{1}}
    \proofstep{1}{F\colon A\sur B}{\FormulaRefAuto{Surjektive Funktion}{2,3}}
\end{tabproof}

\FormulaThmDelta{F\colon A\inj B\land F\colon A\sur B \eqvdash F\colon A\bij B }{
\DeltaRow{Mengen}{A\dsep B}
}
\begin{tabproofsplit}
    \proofpart{\(\vdash\)}
    \proofstep{1}{F\colon A\inj B\land F\colon A\sur B}{\rA}
    \proofstep{1}{F\colon A\inj B}{\rAEa{1}}
    \proofstep{1}{F\colon A\sur B}{\rAEb{1}}
    \proofstep{1}{F\colon A\bij B}{\FormulaRefAuto{F\colon A\inj B\dsep F\colon A\sur B \vdash F\colon A\bij B}{2,3}}
    \closeproofpart
    \proofpart{\(\dashv\)}
    \proofstep{1}{F\colon A\bij B}{\rA}
    \proofstep{1}{F\colon A\inj B}{\FormulaRefAuto{F\colon A\bij B \vdash F\colon A\inj B}{1}}
    \proofstep{1}{F\colon A\sur B}{\FormulaRefAuto{F\colon A\bij B \vdash F\colon A\sur B}{1}}
    \proofstep{1}{F\colon A\inj B\land F\colon A\sur B}{\rAI{2,3}}
    \closeproofpart
\end{tabproofsplit}

\FormulaThmDelta[Eindeutigkeit des Urbildes bei Bijektivität]{%
F\colon A\bij B \dsep y\in B \vdash \exists! x\in A\; F(x)=y
}{
  \DeltaRow{Mengen}{A \dsep B \dsep F\dsep x \dsep y \dsep z}
}
\begin{tabproofsplit}
\proofpart{Existenz}
% Annahmen
  \proofstep{1}{F\colon A\bij B}{\rA}
  \proofstep{2}{y\in B}{\rA}

  % Existenz eines Urbildes aus Surjektivität
  \proofstep{1,2}{\exists x\in A\,F(x)=y}{%
    \FormulaRefAuto{y\in B\vdash \exists x\in A\; F(x)=y}{1,2}}

  % Uniqueness-Teil: angenommen zwei Urbilder x,z

\closeproofpart

\proofpart{Eindeutigkeit}
  \setcounter{proofstepnr}{3}% Fortlaufende Schritte dieses Beweises.
\proofstep{4}{x\in A \land F(x)=y}{\rA}
  \proofstep{5}{z\in A \land F(z)=y}{\rA}

  % Aus den Konjunktionen die Einzelteile holen
  \proofstep{4}{x\in A}{\rAEa{4}}
  \proofstep{4}{F(x)=y}{\rAEb{4}}
  \proofstep{5}{z\in A}{\rAEa{5}}
  \proofstep{5}{F(z)=y}{\rAEb{5}}

  % F(x)=y und F(z)=y ⇒ F(x)=F(z) (über Gleichheitseigenschaften)
  \proofstep{5}{y=F(z)}{\FormulaRefAuto{a=b \vdash b=a}{9}}
  \proofstep{4,5}{F(x)=F(z)}{\FormulaRefAuto{a = b,\, b = c \vdash a = c}{7,10}}

  % Injektivität verwenden: aus x,z∈A und F(x)=F(z) folgt x=z
  \proofstep{1,4,5}{x=z}{%
    \FormulaRefAuto{F\colon A\inj B\dsep x\in A\dsep y\in A\dsep F(x)=F(y)\vdash x=y}{1,6,8,11}}

  % ∃!-Einführung: aus Existenz (3) und Eindeutigkeit (Annahmen 4,5; Schluss 12)

\closeproofpart

\proofpart{Eindeutige Existenz}
  \setcounter{proofstepnr}{12}% Fortlaufende Schritte dieses Beweises.
\proofstep{1,2}{\exists! x\in A\,F(x)=y}{\UEI{3,4,5,12}}

\closeproofpart
\end{tabproofsplit}

\subsection{Urbildmengen unter bijektiven Funktionen}

\FormulaThmDelta[Urbilder liegen im Definitionsbereich bijektiver Funktionen]{%
  C\subseteq B\dsep F\colon A\bij B \vdash F^{-1}[C]\subseteq A
}{
  \DeltaRow{Mengen}{A\dsep B\dsep C}
  \DeltaRow{Funktionen}{F}
}
\begin{tabproof}
  \proofstep{1}{C\subseteq B}{\rA}
  \proofstep{2}{F\colon A\bij B}{\rA}
  \proofstep{2}{F\colon A\to B}{\FormulaRefAuto{Bijektive Funktion}{2}}
  \proofstep{1,2}{F^{-1}[C]\subseteq A}{%
    \FormulaRefAuto{F\colon A\to B\dsep C\subseteq B \vdash F^{-1}[C]\subseteq A}{3,1}}
\end{tabproof}


\section{Bild einer bijektiven Funktion}

\FormulaThmDelta[Bild einer bijektiven Funktion ist die Zielmenge]{%
  F\colon A\bij B \vdash F[A]=B
}{%
  \DeltaRow{Mengen}{A\dsep B\dsep x\dsep y}
  \DeltaRow{Funktionen}{F}
}
\begin{tabproofsplitwide}
  \proofpartwideind[Teilmengenrichtung \(F[A]\subseteq B\)]{%
    F\colon A\bij B \vdash F[A]\subseteq B%
  }[%
    F\colon A\bij B \vdash F[A]\subseteq B%
  ]
    \proofstepwidestar[1]{F\colon A\bij B}{\rA}
    \proofstepwidestar[1]{F\colon A\to B}{\FormulaRefAuto{Bijektive Funktion}{1}}
    \proofstepwidestar[1]{F[A]\subseteq B}{%
      \FormulaRefAuto{F\colon A\to B \vdash F[A]\subseteq B}{2}}
  \closeproofpartwideind

  \proofpartwideind[Teilmengenrichtung \(B\subseteq F[A]\)]{%
    F\colon A\bij B \vdash B\subseteq F[A]%
  }[%
    F\colon A\bij B \vdash B\subseteq F[A]%
  ]
    \proofstepwidestar[1]{F\colon A\bij B}{\rA}
    \proofstepwidestar[1]{F\colon A\sur B}{\FormulaRefAuto{Bijektive Funktion}{1}}
    \proofstepwidestar[1]{F\colon A\to B}{\FormulaRefAuto{Bijektive Funktion}{1}}

    \proofstepwidestar[4]{y\in B}{\rA}
    \proofstepwidestar[1,4]{\exists x\in A\,F(x)=y}{%
      \FormulaRefAuto{y\in B\vdash \exists x\in A\; F(x)=y}{4}}

    \proofstepwidestar[6]{x\in A\land F(x)=y}{\rA}
    \proofstepwidestar[6]{x\in A}{\rAEa{6}}
    \proofstepwidestar[6]{F(x)=y}{\rAEb{6}}
    \proofstepwidestar[6]{y=F(x)}{\FormulaRefAuto{a=b \vdash b=a}{8}}
    \proofstepwidestar[6]{x\in A\land y=F(x)}{\rAI{7,9}}
    \proofstepwidestar[6]{\exists x\in A\,y=F(x)}{\rEI{10}}
    \proofstepwidestar[1,4,6]{y\in F[A]}{%
      \FormulaRefAuto{F\colon A\to B\dsep \exists x\in A\,y=F(x)\vdash y\in F[A]}{3,11}}
    \proofstepwidestar[1,4]{y\in F[A]}{\rEE{5,6,12}}

    \proofstepwidestar[1]{\forall y\,\bigl(y\in B\rightarrow y\in F[A]\bigr)}{\rUI{\rRI{4,13}}}
    \proofstepwidestar[1]{B\subseteq F[A]}{%
      \FormulaRefAuto{A\subseteq B \coloneq \forall x\,(x\in A \rightarrow x\in B)}{14}}
  \closeproofpartwideind

  \proofpartwide{Zusammenfassung}
    \proofstepwidestar[1]{F\colon A\bij B}{\rA}
    \proofstepwidestar[1]{F[A]\subseteq B}{%
      \FormulaRefAuto{F\colon A\bij B \vdash F[A]\subseteq B}{1}}
    \proofstepwidestar[1]{B\subseteq F[A]}{%
      \FormulaRefAuto{F\colon A\bij B \vdash B\subseteq F[A]}{1}}
    \proofstepwidestar[1]{F[A]=B}{%
      \FormulaRefAuto{A\subseteq B\dsep B\subseteq A \vdash A=B}{2,3}}
  \closeproofpartwide
\end{tabproofsplitwide}

\chapter{Beispiele und Konstruktionen}

\section{Identitätsfunktion}

\subsection{Definition der Identitätsfunktion}

\FormulaDefDeltaK[Identitätsfunktion auf \(A\)]{%
  \begin{aligned}[t]
    &\Id_A\colon A\to A,\\[-2pt]
    &\forall x\in A\,\bigl(\Id_A(x)\coloneqq x\bigr)
  \end{aligned}%
}{IdentityFunctionDef}{%
\DeltaRow{Mengen}{A}
\DeltaRow{Funktionen}{\Id_A}
}
\begin{tabproof}
  \proofstep{1}{x\in A}{\rA}
\end{tabproof}

\FormulaThmDeltaK[Identitätsfunktion]
{\Id_A\colon A\to A}{IdA}
{
\DeltaRow{Mengen}{A}
}
\begin{tabproof}
    \proofstep{}{\Id_A\colon A\to A}{%
      \FormulaRefAuto{IdentityFunctionDef}[def]}
\end{tabproof}


\FormulaThmDelta
{\TotRel{\Id_A,A,A}}
{
\DeltaRow{Mengen}{A}
}
\begin{tabproof}
    \proofstep{}{\Id_A\colon A\to A}{\FormulaRefAuto{IdA}}
    \proofstep{}{\TotRel{\Id_A,A,A}}{\FormulaRefAuto{F\colon A\to B\vdash\TotRel{F,A,B}}{1}}
\end{tabproof}


\FormulaThmDelta[Identitätsfunktion auf \(A\)]
{x\in A\vdash \Id_A(x)=x}
{
\DeltaRow{Mengen}{A\dsep x}
}
\begin{tabproof}
    \proofstep{1}{x\in A}{\rA}
    \proofstep{}{\forall u\in A\,\bigl(\Id_A(u)=u\bigr)}{%
      \FormulaRefAuto{IdentityFunctionDef}[def]}
    \proofstep{1}{\Id_A(x)=x}{%
      \rRE{\rUE{2},1}}
\end{tabproof}

\FormulaThmDelta{%
  x\in A\vdash x=\Id_A(x)%
}{
  \DeltaRow{Mengen}{A\dsep x}
}
\begin{tabproof}
  \proofstep{1}{x\in A}{\rA}
  \proofstep{1}{\Id_A(x)=x}{%
    \FormulaRefAuto{x\in A\vdash \Id_A(x)=x}}
  \proofstep{1}{x=\Id_A(x)}{%
    \FormulaRefAuto{a=b\vdash b=a}}
\end{tabproof}

\subsection{Die Bijektionseigenschaften}

\paragraph{Injektivität}
\FormulaThmDelta[Injektivität von \(\Id_A\)]{%
x\in A \dsep y\in A \dsep \Id_A(x)=\Id_A(y) \vdash x=y
}{
  \DeltaRow{Mengen}{x \dsep y \dsep A}
}
\begin{tabproof}
  \proofstep{1}{\Id_A(x)=\Id_A(y)}{\rA}
  \proofstep{2}{x\in A}{\rA}
  \proofstep{3}{y\in A}{\rA}
  \proofstep{2}{\Id_A(x)=x}{\FormulaRefAuto{x\in A\vdash \Id_A(x)=x}{2}}
  \proofstep{3}{\Id_A(y)=y}{\FormulaRefAuto{x\in A\vdash \Id_A(x)=x}{3}}
  \proofstep{1,2,3}{x=y}{\FormulaRefAuto{a = b\dsep a = c\dsep b=d\vdash c = d}{1,4,5}}
\end{tabproof}

\paragraph{Surjektivität}
\FormulaThmDelta[Surjektivität von \(\Id_A\)]{%
x\in A \vdash \exists y\in A\; \Id_A(y)=x
}{
  \DeltaRow{Mengen}{x \dsep y \dsep A}
}
\begin{tabproof}
  \proofstep{1}{x\in A}{\rA}
  \proofstep{1}{\Id_A(x)=x}{\FormulaRefAuto{x\in A\vdash \Id_A(x)=x}{1}}
  \proofstep{1}{\exists y\in A\; \Id_A(y)=x}{\rEI{\rAI{1,2}}}
\end{tabproof}

\paragraph{Bijektivität}
\FormulaThmDelta[\(\Id_A\) als bijektive Funktion]{%
\Id_A\colon A\bij A
}{
  \DeltaRow{Mengen}{A}
}
\begin{tabproofsplit}
\proofpart{Funktionstyp}
\proofstep{}{\Id_A\colon A\to A}{\FormulaRefAuto{IdA}}

\closeproofpart

\proofpart{Injektivität}
  \setcounter{proofstepnr}{1}% Fortlaufende Schritte dieses Beweises.
\proofstep{}{\forall x,y\in A\,(\Id_A(x)=\Id_A(y) \rightarrow x=y)}{\FormulaRefAuto{x\in A \dsep y\in A \dsep \Id_A(x)=\Id_A(y) \vdash x=y}}

\closeproofpart

\proofpart{Surjektivität}
  \setcounter{proofstepnr}{2}% Fortlaufende Schritte dieses Beweises.
\proofstep{}{\forall x\in A\exists y\in A\,\Id_A(y)=x}{\FormulaRefAuto{x\in A \vdash \exists y\in A\; \Id_A(y)=x}}

\closeproofpart

\proofpart{Bijektivität}
  \setcounter{proofstepnr}{3}% Fortlaufende Schritte dieses Beweises.
\proofstep{}{\Id_A\colon A\bij A}{\FormulaRefAuto{Bijektive Funktion}{1,2,3}}

\closeproofpart
\end{tabproofsplit}


\section{Einschränkungen und Korestriktionen}

\FormulaThmDelta[Einschränkung auf ein Urbild als bijektive Funktion]{%
  F\colon A\bij B \dsep T\subseteq B \vdash F\restriction_{F^{-1}[T]}\colon F^{-1}[T]\bij T
}{
  \DeltaRow{Mengen}{A\dsep B\dsep T}
  \DeltaRow{Funktionen}{F}
}
\begin{tabproofsplit}
\proofpart{Injektivität der Einschränkung}
\proofstep{1}{F\colon A\bij B}{\rA}
  \proofstep{2}{T\subseteq B}{\rA}
  \proofstep{1}{F\colon A\inj B}{\FormulaRefAuto{Bijektive Funktion}{1}}
  \proofstep{1}{F\colon A\sur B}{\FormulaRefAuto{Bijektive Funktion}{1}}
  \proofstep{2,3}{F\restriction_{F^{-1}[T]}\colon F^{-1}[T]\inj T}{%
    \FormulaRefAuto{F\colon A\inj B\dsep T\subseteq B \vdash F\restriction_{F^{-1}[T]}\colon F^{-1}[T]\inj T}{3,2}}

\closeproofpart

\proofpart{Surjektivität der Einschränkung}
  \setcounter{proofstepnr}{5}% Fortlaufende Schritte dieses Beweises.
\proofstep{2,4}{F\restriction_{F^{-1}[T]}\colon F^{-1}[T]\sur T}{%
    \FormulaRefAuto{F\colon A\sur B \dsep T\subseteq B \vdash F\restriction_{F^{-1}[T]}\colon F^{-1}[T]\sur T}{4,2}}

\closeproofpart

\proofpart{Bijektivität}
  \setcounter{proofstepnr}{6}% Fortlaufende Schritte dieses Beweises.
\proofstep{1,2}{F\restriction_{F^{-1}[T]}\colon F^{-1}[T]\bij T}{%
    \FormulaRefAuto{Bijektive Funktion}{5,6}}

\closeproofpart
\end{tabproofsplit}


Der vorstehende Satz ist exakt: Eine Einschränkung einer Bijektion ist genau
dann bijektiv auf den gewählten Zielbereich, wenn ihr Definitionsbereich
punktweise dessen Urbild beschreibt.

\FormulaThmDeltaK[Punktweises Kriterium für bijektive Einschränkungen]{%
  \begin{aligned}[t]
    &F\colon A\bij B\dsep S\subseteq A\dsep T\subseteq B\vdash{}\\[-2pt]
    &\Bigl(
      F\restriction_S\colon S\bij T
      \leftrightarrow
      \forall x\in A\,\bigl(x\in S\leftrightarrow F(x)\in T\bigr)
    \Bigr)
  \end{aligned}%
}{BijectiveRestrictionIffPointwiseCompatibility}{%
  \DeltaRow{Mengen}{A\dsep B\dsep S\dsep T\dsep x}
  \DeltaRow{Funktionen}{F}
}
\begin{tabproofsplitwide}
  \proofpartwideindR[Bildrichtung \(x\in S\Rightarrow F(x)\in T\)]{%
    \begin{aligned}[t]
      &F\colon A\bij B\dsep S\subseteq A\dsep T\subseteq B\dsep
        F\restriction_S\colon S\bij T\dsep x\in A\dsep x\in S\\[-2pt]
      &\vdash F(x)\in T
    \end{aligned}%
  }{%
    F\colon A\bij B\dsep S\subseteq A\dsep T\subseteq B\dsep
    F\restriction_S\colon S\bij T\dsep x\in A\dsep x\in S
    \vdash F(x)\in T
  }[BijectiveRestrictionSubsetMembershipImpliesTargetMembership]
    \proofstepwidestar[1]{F\colon A\bij B}{\rA}
    \proofstepwidestar[2]{S\subseteq A}{\rA}
    \proofstepwidestar[3]{T\subseteq B}{\rA}
    \proofstepwidestar[4]{F\restriction_S\colon S\bij T}{\rA}
    \proofstepwidestar[5]{x\in A}{\rA}
    \proofstepwidestar[6]{x\in S}{\rA}
    \proofstepwidestar[4,6]{(F\restriction_S)(x)\in T}{%
      \FormulaRefAuto{F\colon A\bij B\dsep x\in A\vdash F(x)\in B}[thm]{4,6}}
    \proofstepwidestarbreakkeep[2,6]{(F\restriction_S)(x)=F(x)}{%
      \FormulaRefAuto{C\subseteq A\dsep x\in C\vdash
        F\restriction_C(x)=F(x)}[thm]{2,6}}
    \proofstepwidestar[2,4,6]{F(x)\in T}{\rIE{8,7}}
  \closeproofpartwideindR

  \proofpartwideindR[Urbildrichtung \(F(x)\in T\Rightarrow x\in S\)]{%
    \begin{aligned}[t]
      &F\colon A\bij B\dsep S\subseteq A\dsep T\subseteq B\dsep
        F\restriction_S\colon S\bij T\dsep x\in A\dsep F(x)\in T\\[-2pt]
      &\vdash x\in S
    \end{aligned}%
  }{%
    F\colon A\bij B\dsep S\subseteq A\dsep T\subseteq B\dsep
    F\restriction_S\colon S\bij T\dsep x\in A\dsep F(x)\in T
    \vdash x\in S
  }[BijectiveRestrictionTargetMembershipImpliesSubsetMembership]
    \proofstepwidestar[1]{F\colon A\bij B}{\rA}
    \proofstepwidestar[2]{S\subseteq A}{\rA}
    \proofstepwidestar[3]{T\subseteq B}{\rA}
    \proofstepwidestar[4]{F\restriction_S\colon S\bij T}{\rA}
    \proofstepwidestar[5]{x\in A}{\rA}
    \proofstepwidestar[6]{F(x)\in T}{\rA}
    \proofstepwidestar[4]{(F\restriction_S)[S]=T}{%
      \FormulaRefAuto{F\colon A\bij B \vdash F[A]=B}[thm]{4}}
    \proofstepwidestar[1]{F\colon A\to B}{%
      \FormulaRefAuto{Bijektive Funktion}[def]{1}}
    \proofstepwidestarbreakkeep[1,2]{(F\restriction_S)[S]=F[S]}{%
      \FormulaRefAuto{C\subseteq A \dsep F\colon A\to B
        \vdash (F\restriction_C)[C]=F[C]}[thm]{2,8}}
    \proofstepwidestar[1,2,4]{T=F[S]}{%
      \FormulaRefAuto{a=b\dsep a=c\vdash b=c}[thm]{7,9}}
    \proofstepwidestar[1,2,4,6]{F(x)\in F[S]}{\rIE{10,6}}
    \proofstepwidestar[1]{F\colon A\inj B}{%
      \FormulaRefAuto{F\colon A\bij B\vdash F\colon A\inj B}[thm]{1}}
    \proofstepwidestar[2]{S\in\powerset(A)}{%
      \FormulaRefAuto{x\in\powerset(A)\eqvdash x\subseteq A}[thm]{2}}
    \proofstepwidestar[1,2,4,5,6]{x\in S}{%
      \FormulaRefAuto{F\colon A\inj B \dsep Y\in\powerset(A)
        \dsep x\in A\dsep F(x)\in F[Y] \vdash x\in Y}[thm]{%
        12,13,5,11}}
  \closeproofpartwideindR

  \pagebreak[3]
  \proofpartwideindR[Zusammenführung zur punktweisen Äquivalenz]{%
    \begin{aligned}[t]
      &F\colon A\bij B\dsep S\subseteq A\dsep T\subseteq B\dsep
        F\restriction_S\colon S\bij T\\[-2pt]
      &\vdash
        \forall x\in A\,\bigl(x\in S\leftrightarrow F(x)\in T\bigr)
    \end{aligned}%
  }{%
    F\colon A\bij B\dsep S\subseteq A\dsep T\subseteq B\dsep
    F\restriction_S\colon S\bij T
    \vdash
    \forall x\in A\,\bigl(x\in S\leftrightarrow F(x)\in T\bigr)
  }[BijectiveRestrictionImpliesPointwiseCompatibility]
    \proofstepwidestar[1]{F\colon A\bij B}{\rA}
    \proofstepwidestar[2]{S\subseteq A}{\rA}
    \proofstepwidestar[3]{T\subseteq B}{\rA}
    \proofstepwidestar[4]{F\restriction_S\colon S\bij T}{\rA}
    \proofstepwidestar[5]{x\in A}{\rA}
    \proofstepwidestar[6]{x\in S}{\rA}
    \proofstepwidestar[1,2,3,4,5,6]{F(x)\in T}{%
      \FormulaRefAuto{BijectiveRestrictionSubsetMembershipImpliesTargetMembership}[thm]{%
        1,2,3,4,5,6}}
    \proofstepwidestar[1,2,3,4,5]{x\in S\rightarrow F(x)\in T}{%
      \rRI{6,7}}
    \proofstepwidestar[9]{F(x)\in T}{\rA}
    \proofstepwidestar[1,2,3,4,5,9]{x\in S}{%
      \FormulaRefAuto{BijectiveRestrictionTargetMembershipImpliesSubsetMembership}[thm]{%
        1,2,3,4,5,9}}
    \proofstepwidestar[1,2,3,4,5]{F(x)\in T\rightarrow x\in S}{%
      \rRI{9,10}}
    \proofstepwidestar[1,2,3,4,5]{x\in S\leftrightarrow F(x)\in T}{%
      \rLRI{8,11}}
    \proofstepwidestar[1,2,3,4]{%
      \forall x\in A\,\bigl(x\in S\leftrightarrow F(x)\in T\bigr)}{%
      \rUI{\rRI{5,12}}}
  \closeproofpartwideindR

  \proofpartwideindR[Rückrichtung zur eingeschränkten Bijektivität]{%
    \begin{aligned}[t]
      &F\colon A\bij B\dsep S\subseteq A\dsep T\subseteq B\dsep
        \forall x\in A\,\bigl(x\in S\leftrightarrow F(x)\in T\bigr)\\[-2pt]
      &\vdash F\restriction_S\colon S\bij T
    \end{aligned}%
  }{%
    F\colon A\bij B\dsep S\subseteq A\dsep T\subseteq B\dsep
    \forall x\in A\,\bigl(x\in S\leftrightarrow F(x)\in T\bigr)
    \vdash F\restriction_S\colon S\bij T
  }[PointwiseCompatibilityImpliesBijectiveRestriction]
    \proofstepwidestar[1]{F\colon A\bij B}{\rA}
    \proofstepwidestar[2]{S\subseteq A}{\rA}
    \proofstepwidestar[3]{T\subseteq B}{\rA}
    \proofstepwidestar[4]{%
      \forall x\in A\,\bigl(x\in S\leftrightarrow F(x)\in T\bigr)}{\rA}
    \proofstepwidestar[1]{F\colon A\to B}{%
      \FormulaRefAuto{Bijektive Funktion}[def]{1}}
    \proofstepwidestar[6]{x\in A}{\rA}
    \proofstepwidestar[4,6]{x\in S\leftrightarrow F(x)\in T}{%
      \rRE{\rUE{4},6}}
    \proofstepwidestar[4,6]{F(x)\in T\leftrightarrow x\in S}{%
      \FormulaRefAuto{P\leftrightarrow Q\vdash Q\leftrightarrow P}[thm]{7}}
    \proofstepwidestar[4]{%
      \forall x\in A\,\bigl(F(x)\in T\leftrightarrow x\in S\bigr)}{%
      \rUI{\rRI{6,8}}}
    \proofstepwidestar[1,2,3,4]{F^{-1}[T]=S}{%
      \FormulaRefAuto{PreimageEqFromPointwiseIff}[thm]{5,3,2,9}}
    \proofstepwidestar[1,3]{%
      F\restriction_{F^{-1}[T]}\colon F^{-1}[T]\bij T}{%
      \FormulaRefAuto{F\colon A\bij B\dsep T\subseteq B
        \vdash F\restriction_{F^{-1}[T]}\colon F^{-1}[T]\bij T}[thm]{1,3}}
    \proofstepwidestar[1,2,3,4]{F\restriction_S\colon S\bij T}{%
      \rIE{10,11}}
  \closeproofpartwideindR

  \pagebreak[4]
  \proofpartwide{Schluss}
    \proofstepwidestar[1]{F\colon A\bij B}{\rA}
    \proofstepwidestar[2]{S\subseteq A}{\rA}
    \proofstepwidestar[3]{T\subseteq B}{\rA}
    \proofstepwidestar[4]{F\restriction_S\colon S\bij T}{\rA}
    \proofstepwidestar[1,2,3,4]{%
      \forall x\in A\,\bigl(x\in S\leftrightarrow F(x)\in T\bigr)}{%
      \FormulaRefAuto{BijectiveRestrictionImpliesPointwiseCompatibility}[thm]{%
        1,2,3,4}}
    \proofstepwidestarbreakkeep[1,2,3]{%
      F\restriction_S\colon S\bij T
      \rightarrow
      \forall x\in A\,\bigl(x\in S\leftrightarrow F(x)\in T\bigr)}{%
      \rRI{4,5}}
    \proofstepwidestar[7]{%
      \forall x\in A\,\bigl(x\in S\leftrightarrow F(x)\in T\bigr)}{\rA}
    \proofstepwidestar[1,2,3,7]{F\restriction_S\colon S\bij T}{%
      \FormulaRefAuto{PointwiseCompatibilityImpliesBijectiveRestriction}[thm]{%
        1,2,3,7}}
    \proofstepwidestarbreakkeep[1,2,3]{%
      \forall x\in A\,\bigl(x\in S\leftrightarrow F(x)\in T\bigr)
      \rightarrow F\restriction_S\colon S\bij T}{\rRI{7,8}}
    \proofstepwidestarbreakkeep[1,2,3]{%
      F\restriction_S\colon S\bij T
      \leftrightarrow
      \forall x\in A\,\bigl(x\in S\leftrightarrow F(x)\in T\bigr)}{%
      \rLRI{6,9}}
  \closeproofpartwide
\end{tabproofsplitwide}


\FormulaThmDelta[Korestriktion auf das Bild einer injektiven Funktion ist bijektiv]{%
  F\colon A\inj B \vdash F\colon A\bij F[A]
}{%
  \DeltaRow{Mengen}{A\dsep B\dsep x\dsep z}
  \DeltaRow{Funktionen}{F}
}
\begin{tabproof}
  \proofstep{1}{F\colon A\inj B}{\rA}
  \proofstep{1}{F\colon A\to B}{\FormulaRefAuto{Injektive Funktion}{1}}
  \proofstep{3}{x\in A}{\rA}
  \proofstep{3}{F(x)=F(x)}{\rII}
  \proofstep{1,3}{\exists z\in A\,F(x)=F(z)}{\rEI{\rAI{3,4}}}
  \proofstep{1,3}{F(x)\in F[A]}{%
    \FormulaRefAuto{F\colon A\to B\dsep \exists x\in A\,y=F(x)\vdash y\in F[A]}{2,5}}
  \proofstep{1}{\forall x\in A\,F(x)\in F[A]}{\rUI{\rRI{3,6}}}
  \proofstep{1}{F\colon A\inj F[A]}{%
    \FormulaRefAuto{F\colon A\inj B\dsep \forall x\in A\,F(x)\in C \vdash F\colon A\inj C}{1,7}}
  \proofstep{1}{F\colon A\sur F[A]}{%
    \FormulaRefAuto{F\colon A\to B \vdash F\colon A\sur F[A]}{2}}
  \proofstep{1}{F\colon A\bij F[A]}{%
    \FormulaRefAuto{F\colon A\inj B\dsep F\colon A\sur B \vdash F\colon A\bij B}{8,9}}
\end{tabproof}

\section{Umkehrfunktion}

\subsection{Definition der Umkehrfunktion}

\FormulaDefDeltaK[Umkehrfunktion von \(F\)]{%
  \begin{aligned}[t]
    &F^{-1}\colon B\to A,\\[-2pt]
    &\forall y\in B\,
      \Bigl(F^{-1}(y)\coloneqq
        \iota x\bigl(x\in A\land F(x)=y\bigr)\Bigr)
  \end{aligned}%
}{InverseFunctionDef}{%
  \DeltaRow{Mengen}{A\dsep B}
  \DeltaPrem{Funktionen}{F\colon A\bij B}
  \DeltaRow{Funktionen}{F^{-1}}
}
\begin{tabproof}
  \proofstep{1}{F\colon A\bij B}{\rA}
  \proofstep{2}{y\in B}{\rA}
  \proofstep{1,2}{%
    \begin{aligned}[t]
      &\iota x\bigl(x\in A\land F(x)=y\bigr)\in A\land{}\\[-2pt]
      &F\!\left(\iota x\bigl(x\in A\land F(x)=y\bigr)\right)=y
    \end{aligned}%
  }{%
    \FormulaRefAuto{F\colon A\bij B\dsep y\in B
      \vdash \exists! x\in A\,F(x)=y}[thm]{1,2}}
  \proofstep{1,2}{\iota x\bigl(x\in A\land F(x)=y\bigr)\in A}{%
    \rAEa{3}}
\end{tabproof}

Der innere \(\iota\)-Term bezeichnet gemäß der Iota-Konvention den durch
den Eindeutigkeitssatz bestimmten Urbildzeugen; nur die äußere
Funktions-Iota-Konstruktion entfällt.

\FormulaThmDeltaR[Umkehrfunktion von \(F\)]
{F^{-1}\colon B\to A}{F\colon A\bij B\vdash F^{-1}\colon B\to A}
{
  \DeltaRow{Mengen}{A\dsep B}
  \DeltaPrem{Funktionen}{F\colon A\bij B}
}
\begin{tabproof}
  \proofstep{}{F^{-1}\colon B\to A}{%
    \FormulaRefAuto{InverseFunctionDef}[def]}
\end{tabproof}

\FormulaThmDelta
{\TotRel{F^{-1},B,A}}
{
  \DeltaRow{Mengen}{A\dsep B}
  \DeltaPrem{Funktionen}{F\colon A\bij B}
}
\begin{tabproof}
    \proofstep{}{F^{-1}\colon B\to A}{\FormulaRefAuto{F\colon A\bij B\vdash F^{-1}\colon B\to A}}
    \proofstep{1}{\TotRel{F^{-1},B,A}}{\FormulaRefAuto{F\colon A\to B\vdash\TotRel{F,A,B}}{1}}
\end{tabproof}

\FormulaThmDelta[Wert der Umkehrfunktion von \(F\)]{%
  y\in B\vdash
  F^{-1}(y)=\iota x\bigl(x\in A\land F(x)=y\bigr)%
}
{
  \DeltaRow{Mengen}{A\dsep B\dsep x\dsep y}
  \DeltaPrem{Funktionen}{F\colon A\bij B}
}
\begin{tabproof}
  \proofstep{1}{y\in B}{\rA}
  \proofstep{}{\forall z\in B\,
    \Bigl(F^{-1}(z)=\iota u\bigl(u\in A\land F(u)=z\bigr)\Bigr)}{%
    \FormulaRefAuto{InverseFunctionDef}[def]}
  \proofstep{1}{F^{-1}(y)=\iota x\bigl(x\in A\land F(x)=y\bigr)}{%
    \rRE{\rUE{2},1}}
\end{tabproof}

\FormulaThmDelta[Typisierung der Umkehrfunktion]
{y\in B \vdash F^{-1}(y)\in A}
{
  \DeltaRow{Mengen}{A\dsep B\dsep x\dsep y}
  \DeltaPrem{Funktionen}{F\colon A\bij B}
}
\begin{tabproof}
  \proofstep{1}{y\in B}{\rA}
  \proofstep{}{F^{-1}\colon B\to A}{%
    \FormulaRefAuto{F\colon A\bij B\vdash F^{-1}\colon B\to A}[thm]}
  \proofstep{1}{F^{-1}(y)\in A}{%
    \FormulaRefAuto{F\colon A\to B\dsep x\in A\vdash F(x)\in B}[thm]{2,1}}
\end{tabproof}

\subsection{Grundgleichungen der Umkehrfunktion}

\FormulaThmDelta[Links-Umkehrung]
{x\in A \vdash F^{-1}(F(x)) = x}
{
  \DeltaRow{Mengen}{A\dsep B\dsep x}
  \DeltaPrem{Funktionen}{F\colon A\bij B}
}
\begin{tabproof}
  \proofstep{1}{x\in A}{\rA}
  \proofstep{1}{F(x)\in B}{%
    \FormulaRefAuto{F\colon A\to B\dsep x\in A\vdash F(x)\in B}{1}}
  \proofstep{}{\forall z\in B\,
    \Bigl(F^{-1}(z)=\iota u\bigl(u\in A\land F(u)=z\bigr)\Bigr)}{%
    \FormulaRefAuto{InverseFunctionDef}[def]}
  \proofstep{1}{%
    F^{-1}(F(x))=
    \iota u\bigl(u\in A\land F(u)=F(x)\bigr)}{%
    \rRE{\rUE{3},2}}
  \proofstep{1}{%
    \begin{aligned}[t]
      &\iota u\bigl(u\in A\land F(u)=F(x)\bigr)\in A\land{}\\[-2pt]
      &F\!\left(\iota u\bigl(u\in A\land F(u)=F(x)\bigr)\right)=F(x)
    \end{aligned}%
  }{%
    \FormulaRefAuto{F\colon A\bij B\dsep y\in B
      \vdash \exists! x\in A\,F(x)=y}[thm]{2}}
  \proofstep{1}{%
    \iota u\bigl(u\in A\land F(u)=F(x)\bigr)\in A}{\rAEa{5}}
  \proofstep{1}{%
    F\!\left(\iota u\bigl(u\in A\land F(u)=F(x)\bigr)\right)=F(x)}{%
    \rAEb{5}}
  \proofstep{1}{F^{-1}(F(x))\in A}{\rIE{4,6}}
  \proofstep{1}{F(F^{-1}(F(x)))=F(x)}{\rIE{4,7}}
  \proofstep{}{F\colon A\inj B}{%
      \FormulaRefAuto{F\colon A\bij B \vdash F\colon A\inj B}}
  \proofstep{1}{F^{-1}(F(x))=x}{%
    \FormulaRefAuto{F\colon A\inj B\dsep x\in A\dsep y\in A\dsep F(x)=F(y)\vdash x=y}{%
      10,8,1,9}}
\end{tabproof}

\FormulaThmDeltaR[Rechts-Umkehrung]
{y\in B \vdash F(F^{-1}(y)) = y}{F\colon A\bij B \dsep y\in B \vdash F(F^{-1}(y)) = y}
{
  \DeltaRow{Mengen}{A\dsep B\dsep y}
  \DeltaPrem{Funktionen}{F\colon A\bij B}
}
\begin{tabproof}
  \proofstep{1}{y\in B}{\rA}
  \proofstep{}{\forall z\in B\,
    \Bigl(F^{-1}(z)=\iota u\bigl(u\in A\land F(u)=z\bigr)\Bigr)}{%
    \FormulaRefAuto{InverseFunctionDef}[def]}
  \proofstep{1}{%
    F^{-1}(y)=\iota u\bigl(u\in A\land F(u)=y\bigr)}{%
    \rRE{\rUE{2},1}}
  \proofstep{1}{%
    \begin{aligned}[t]
      &\iota u\bigl(u\in A\land F(u)=y\bigr)\in A\land{}\\[-2pt]
      &F\!\left(\iota u\bigl(u\in A\land F(u)=y\bigr)\right)=y
    \end{aligned}%
  }{%
    \FormulaRefAuto{F\colon A\bij B\dsep y\in B
      \vdash \exists! x\in A\,F(x)=y}[thm]{1}}
  \proofstep{1}{%
    F\!\left(\iota u\bigl(u\in A\land F(u)=y\bigr)\right)=y}{\rAEb{4}}
  \proofstep{1}{F(F^{-1}(y))=y}{\rIE{3,5}}
\end{tabproof}

\subsection{Mengenbezogene Grundgleichungen der Umkehrfunktion}

\FormulaThmDelta[Bild des Urbilds einer Teilmenge unter einer Bijektion]{%
  Y\subseteq B \dsep F\colon A\bij B \vdash F[F^{-1}[Y]]=Y
}{%
  \DeltaRow{Mengen}{A\dsep B\dsep Y}
  \DeltaRow{Funktionen}{F}
}
\begin{tabproof}
  \proofstep{1}{Y\subseteq B}{\rA}
  \proofstep{2}{F\colon A\bij B}{\rA}
  \proofstep{1,2}{F^{-1}[Y]\subseteq A}{%
    \FormulaRefAuto{C\subseteq B\dsep F\colon A\bij B \vdash F^{-1}[C]\subseteq A}{1,2}}
  \proofstep{2}{F\colon A\to B}{\FormulaRefAuto{Bijektive Funktion}{2}}
  \proofstep{1,2}{F\restriction_{F^{-1}[Y]}\colon F^{-1}[Y]\bij Y}{%
    \FormulaRefAuto{F\colon A\bij B \dsep T\subseteq B \vdash F\restriction_{F^{-1}[T]}\colon F^{-1}[T]\bij T}{2,1}}
  \proofstep{1,2}{(F\restriction_{F^{-1}[Y]})[F^{-1}[Y]]=Y}{%
    \FormulaRefAuto{F\colon A\bij B \vdash F[A]=B}{5}}
  \proofstep{1,2}{(F\restriction_{F^{-1}[Y]})[F^{-1}[Y]]=F[F^{-1}[Y]]}{%
    \FormulaRefAuto{C\subseteq A \dsep F\colon A\to B \vdash (F\restriction_C)[C]=F[C]}{3,4}}
  \proofstep{1,2}{F[F^{-1}[Y]]=(F\restriction_{F^{-1}[Y]})[F^{-1}[Y]]}{%
    \FormulaRefAuto{a=b\vdash b=a}{7}}
  \proofstep{1,2}{F[F^{-1}[Y]]=Y}{%
    \FormulaRefAuto{a=b,\, b=c \vdash a=c}{8,6}}
\end{tabproof}

\subsection{Die Bijektionseigenschaften der Umkehrfunktion}

\FormulaThmDelta[Injektivität von \(F^{-1}\)]%
{x\in B\dsep y\in B\dsep F^{-1}(x)=F^{-1}(y)\vdash x=y}%
{
  \DeltaRow{Mengen}{x\dsep y\dsep A\dsep B}
  \DeltaPrem{Bijektive Funktionen}{F\colon A\bij B}
}
\begin{tabproof}
  \proofstep{1}{x\in B}{\rA}
  \proofstep{2}{y\in B}{\rA}
  \proofstep{3}{F^{-1}(x)=F^{-1}(y)}{\rA}
  \proofstep{}{F^{-1}\colon B\to A}{\FormulaRefAuto{F\colon A\bij B\vdash F^{-1}\colon B\to A}}
  \proofstep{}{F\colon A\to B}{\FormulaRefAuto{Bijektive Funktion}}
  \proofstep{1}{F^{-1}(x)\in A}{\FormulaRefAuto{F\colon A\to B\dsep x\in A\vdash F(x)\in B}{1,4}}
  \proofstep{2}{F^{-1}(y)\in A}{\FormulaRefAuto{F\colon A\to B\dsep x\in A\vdash F(x)\in B}{2,4}}
  \proofstep{1}{F(F^{-1}(x))=x}{%
    \FormulaRefAuto{F\colon A\bij B \dsep y\in B \vdash F(F^{-1}(y)) = y}{1}}
  \proofstep{2}{F(F^{-1}(y))=y}{%
    \FormulaRefAuto{F\colon A\bij B \dsep y\in B \vdash F(F^{-1}(y)) = y}{2}}
  \proofstep{1,2,3}{F(F^{-1}(x))=F(F^{-1}(y))}{%
    \FormulaRefAuto{F\colon A\to B\dsep x\in A\dsep y\in A\dsep x=y\vdash F(x)=F(y)}{5,6,7,3}}
  \proofstep{1,2,3}{x=y}{%
    \FormulaRefAuto{a = b\dsep a = c\dsep b=d\vdash c = d}{10,8,9}}
\end{tabproof}

\FormulaThmDelta[Surjektivität von \(F^{-1}\)]%
{x\in A\vdash \exists y\in B\; F^{-1}(y)=x}%
{
  \DeltaRow{Mengen}{x\dsep y\dsep A\dsep B}
  \DeltaPrem{Bijektive Funktionen}{F\colon A\bij B}
}
\begin{tabproof}
  \proofstep{1}{x\in A}{\rA}
  % Typisierung von F:
  \proofstep{1}{F(x)\in B}{%
    \FormulaRefAuto{F\colon A\to B\dsep x\in A\vdash F(x)\in B}{1}}
  % Links-Umkehrung: F^{-1}(F(x))=x
  \proofstep{1}{F^{-1}(F(x))=x}{%
    \FormulaRefAuto{x\in A \vdash F^{-1}(F(x)) = x}{1}}
  % Zeugenbildung y \coloneq F(x)
  \proofstep{1}{\exists y\in B\; F^{-1}(y)=x}{%
    \rEI{\rAI{2,3}}}
\end{tabproof}

\FormulaThmDeltaR[\(F^{-1}\) als bijektive Funktion]%
{F^{-1}\colon B\bij A}{F\colon A\bij B\vdash F^{-1}\colon B\bij A}%
{
  \DeltaRow{Mengen}{A\dsep B}
  \DeltaPrem{Bijektive Funktionen}{F\colon A\bij B}
}
\begin{tabproofsplit}
\proofpart{Funktionstyp}
\proofstep{}{F^{-1}\colon B\to A}{\FormulaRefAuto{F\colon A\bij B\vdash F^{-1}\colon B\to A}}

\closeproofpart

\proofpart{Injektivität}
  \setcounter{proofstepnr}{1}% Fortlaufende Schritte dieses Beweises.
\proofstep{}{\forall x,y\in B\, (F^{-1}(x)=F^{-1}(y)\rightarrow x=y)}{%
    \FormulaRefAuto{x\in B\dsep y\in B\dsep F^{-1}(x)=F^{-1}(y)\vdash x=y}}

\closeproofpart

\proofpart{Surjektivität}
  \setcounter{proofstepnr}{2}% Fortlaufende Schritte dieses Beweises.
\proofstep{}{\forall x\in A\exists y\in B\, F^{-1}(y)=x}{%
    \FormulaRefAuto{x\in A\vdash \exists y\in B\; F^{-1}(y)=x}}

\closeproofpart

\proofpart{Bijektivität}
  \setcounter{proofstepnr}{3}% Fortlaufende Schritte dieses Beweises.
\proofstep{}{F^{-1}\colon B\bij A}{%
    \FormulaRefAuto{Bijektive Funktion}{1,2,3}}

\closeproofpart
\end{tabproofsplit}


\FormulaThmDelta[Eigeninversität der Identität]{%
\Id_A^{-1}=\Id_A
}{
  \DeltaRow{Mengen}{A}
}
\begin{tabproof}
  \proofstep{}{%
    \Id_A\colon A\bij A}{%
    \FormulaRefAuto{\Id_A\colon A\bij A}}

  \proofstep{2}{%
    x\in A}{%
    \rA}

  \proofstep{2}{%
    \Id_A^{-1}(\Id_A(x)) = x}{%
    \FormulaRefAuto{x\in A \vdash F^{-1}(F(x)) = x}{1,2}}

  \proofstep{2}{%
    \Id_A(x)=x}{%
    \FormulaRefAuto{x\in A\vdash \Id_A(x)=x}{2}}

  \proofstep{2}{%
    \Id_A^{-1}(x)=x}{%
    \rIE{4,3}}

  \proofstep{2}{%
    \Id_A(x)=\Id_A(x)}{%
    \rII}

  \proofstep{2}{%
    x=\Id_A(x)}{%
    \rIE{4,6}}

  \proofstep{2}{%
    \Id_A^{-1}(x)=\Id_A(x)}{%
    \rIE{7,5}}

  \proofstep{}{%
    x\in A \rightarrow \Id_A^{-1}(x)=\Id_A(x)}{%
    \rRI{2,8}}

  \proofstep{}{%
    \forall x\in A\, \Id_A^{-1}(x)=\Id_A(x)}{%
    \rUI{9}}

  \proofstep{}{\Id_A^{-1}\colon A\to A}{%
    \FormulaRefAuto{F\colon A\bij B\vdash F^{-1}\colon B\to A}[thm]{1}}
  \proofstep{}{\Id_A\colon A\to A}{\FormulaRefAuto{IdA}[thm]}
  \proofstep{}{\Id_A^{-1}=\Id_A}{%
    \FormulaRefAuto{FunctionExtensionality}[thm]{11,12,10}}
\end{tabproof}


\subsection{Weitere Eigenschaften}

\FormulaThmDelta{%
\forall y\in B\,\bigl(F(G(y))=y\bigr) \vdash G = F^{-1}
}{
  \DeltaRow{Mengen}{x \dsep y \dsep A \dsep B}
  \DeltaPrem{Bijektive Funktionen}{F\colon A\bij B}
  \DeltaPrem{Funktionen}{G\colon B\to A}
}
\begin{tabproof}
  \proofstep{1}{\forall y\in B\,\bigl(F(G(y))=y\bigr)}{\rA}
  \proofstep{2}{y\in B}{\rA}
  \proofstep{1,2}{F(G(y))=y}{\rRE{\rUE{1},2}}
  \proofstep{1,2}{F(F^{-1}(y))=y}{\FormulaRefAuto{F\colon A\bij B \dsep y\in B \vdash F(F^{-1}(y)) = y}{2}}
  \proofstep{2}{G(y)\in A}{\FormulaRefAuto{F\colon A\to B\dsep x\in A\vdash F(x)\in B}{2}}
  \proofstep{2}{F^{-1}(y)\in A}{\FormulaRefAuto{F\colon A\to B\dsep x\in A\vdash F(x)\in B}{2}}
  \proofstep{}{F\colon A\inj B}{%
      \FormulaRefAuto{F\colon A\bij B \vdash F\colon A\inj B}}
  \proofstep{1,2}{G(y)=F^{-1}(y)}{%
    \FormulaRefAuto{F\colon A\inj B\dsep x\in A\dsep y\in A\dsep F(x)=F(y)\vdash x=y}{%
      7,5,6,3,4}}
  \proofstep{1}{\forall y\in B\,\Bigl(G(y)=F^{-1}(y)\Bigr)}{\rUI{\rRI{2,8}}}
  \proofstep{}{G\colon B\to A}{\rA}
  \proofstep{}{F^{-1}\colon B\to A}{%
    \FormulaRefAuto{F\colon A\bij B\vdash F^{-1}\colon B\to A}[thm]}
  \proofstep{1}{G=F^{-1}}{%
    \FormulaRefAuto{FunctionExtensionality}[thm]{10,11,9}}
\end{tabproof}

\FormulaThmDelta{%
\forall x\in A\,\bigl(F^{-1}(G(x))=x\bigr) \vdash G = F
}{
  \DeltaRow{Mengen}{x \dsep A \dsep B}
  \DeltaPrem{Bijektive Funktionen}{F\colon A\bij B}
  \DeltaPrem{Funktionen}{G\colon A\to B}
}
\begin{tabproof}
  \proofstep{1}{\forall x\in A\,\bigl(F^{-1}(G(x))=x\bigr)}{\rA}
  \proofstep{2}{x\in A}{\rA}
  \proofstep{1,2}{F^{-1}(G(x))=x}{\rRE{\rUE{1},2}}
  \proofstep{2}{F^{-1}\bigl(F(x)\bigr)=x}{\FormulaRefAuto{x\in A \vdash F^{-1}(F(x)) = x}{2}}
  \proofstep{2}{G(x)\in B}{\FormulaRefAuto{F\colon A\to B\dsep x\in A\vdash F(x)\in B}{2}}
  \proofstep{2}{F(x)\in B}{\FormulaRefAuto{F\colon A\to B\dsep x\in A\vdash F(x)\in B}{2}}
  \proofstep{1,2}{F^{-1}(G(x))=F^{-1}(F(x))}{\FormulaRefAuto{a=b, c=b\vdash a=c}{3,4}}
  \proofstep{1,2}{G(x)=F(x)}{\FormulaRefAuto{x\in B\dsep y\in B\dsep F^{-1}(x)=F^{-1}(y)\vdash x=y}{5,6,7}}
  \proofstep{1}{\forall x\in A\,\Bigl(G(x)=F(x)\Bigr)}{\rUI{\rRI{2,8}}}
  \proofstep{}{G\colon A\to B}{\rA}
  \proofstep{}{F\colon A\to B}{\FormulaRefAuto{Bijektive Funktion}}
  \proofstep{1}{G=F}{\FormulaRefAuto{FunctionExtensionality}[thm]{10,11,9}}
\end{tabproof}


\section{Kompositionen und Inversen}

\subsection{Erhalt der Bijektivität}

\FormulaThmDeltaK[Die Komposition als bijektive Funktion]{%
F\colon A\bij B \dsep G\colon B\bij C\vdash G\circ F\colon A\bij C
}{BijectiveFunctionComposition}{
  \DeltaRow{Mengen}{A\dsep B\dsep C}
  \DeltaPrem{Funktionen}{F\colon A\to B \dsep G\colon B\to C}
}
\begin{tabproofsplit}
\proofpart{Injektivität der Komposition}
\proofstep{1}{F\colon A\bij B}{\rA}
    \proofstep{2}{G\colon B\bij C}{\rA}
    \proofstep{1}{F\colon A\sur B}{\FormulaRefAuto{F\colon A\bij B\vdash F\colon A\sur B}{1}}
    \proofstep{1}{F\colon A\inj B}{\FormulaRefAuto{F\colon A\bij B\vdash F\colon A\inj B}{1}}
    \proofstep{2}{G\colon B\sur C}{\FormulaRefAuto{F\colon A\bij B\vdash F\colon A\sur B}{2}}
    \proofstep{2}{G\colon B\inj C}{\FormulaRefAuto{F\colon A\bij B\vdash F\colon A\inj B}{2}}
    \proofstep{1,2}{ G\circ F\colon A\inj C}{\FormulaRefAuto{F\colon A\inj B \dsep G\colon B\inj C\vdash G\circ F\colon A\inj C}{4,6}}

\closeproofpart

\proofpart{Surjektivität der Komposition}
  \setcounter{proofstepnr}{7}% Fortlaufende Schritte dieses Beweises.
\proofstep{1,2}{ G\circ F\colon A\sur C}{\FormulaRefAuto{G\circ F\colon A\sur C}{3,5}}

\closeproofpart

\proofpart{Bijektivität}
  \setcounter{proofstepnr}{8}% Fortlaufende Schritte dieses Beweises.
\proofstep{1,2}{ G\circ F\colon A\bij C}{\FormulaRefAuto{F\colon A\inj B\dsep F\colon A\sur B\vdash F\colon A\bij B}{7,8}}

\closeproofpart
\end{tabproofsplit}



\FormulaThmDelta[Linksneutralität]{%
\Id_B \circ F = F
}{
  \DeltaRow{Mengen}{A \dsep B}
  \DeltaPrem{Funktionen}{F\colon A\to B}
}
\begin{tabproofwide}
  \proofstepwidestar[1]{x\in A}{\rA}
  \proofstepwidestar[1]{F(x)\in B}{\FormulaRefAuto{F\colon A\to B\dsep x\in A\vdash F(x)\in B}{1}}

  \proofstepwide[1]{(\Id_B\circ F)(x)}{=}{\Id_B(F(x))}{\FormulaRefAuto{x\in A \vdash (G\circ F)(x)=G(F(x))}{1}}
  \proofstepwide[1]{}{=}{F(x)}{\FormulaRefAuto{x\in A \vdash \Id_A(x)=x}{3}}
  \proofstepwide[1]{(\Id_B\circ F)(x)}{=}{F(x)}{\rChain{3,4}}

  \proofstepwide[]{\forall x\in A\,\bigl((\Id_B\circ F)(x)}{=}{F(x)\bigr)}{\rUI{\rRI{1,5}}}
  \proofstepwidestar[]{F\colon A\to B}{\rA}
  \proofstepwidestar[]{\Id_B\colon B\to B}{\FormulaRefAuto{IdA}[thm]}
  \proofstepwidestar[]{\Id_B\circ F\colon A\to B}{%
    \FormulaRefAuto{CompositionDef}[def]{7,8}}
  \proofstepwide[]{\Id_B\circ F}{=}{F}{%
    \FormulaRefAuto{FunctionExtensionality}[thm]{9,7,6}}
\end{tabproofwide}

% — Rechtsneutralität —
\FormulaThmDelta[Rechtsneutralität]{%
F \circ \Id_A = F
}{
  \DeltaRow{Mengen}{A \dsep B}
  \DeltaPrem{Funktionen}{F\colon A\to B}
}
\begin{tabproofwide}
  \proofstepwidestar[1]{x\in A}{\rA}

  % Fixpunkteigenschaft der Identität auf A
  \proofstepwidestar[1]{\Id_A(x)=x}{\FormulaRefAuto{x\in A \vdash \Id_A(x)=x}{1}}

  % Komposition auswerten
  \proofstepwide[1]{(F\circ \Id_A)(x)}{=}{F(\Id_A(x))}{\FormulaRefAuto{x\in A \vdash (G\circ F)(x)=G(F(x))}{1}}

  % Ersetze \Id_A(x) durch x
  \proofstepwide[1]{}{=}{F(x)}{\rIE{2,3}}

  \proofstepwide[1]{(F\circ \Id_A)(x)}{=}{F(x)}{\rChain{3,4}}

  % Verallgemeinern und Extensionalität
  \proofstepwide[]{\forall x\in A\,\bigl((F\circ \Id_A)(x)}{=}{F(x)\bigr)}{\rUI{\rRI{1,5}}}
  \proofstepwidestar[]{F\colon A\to B}{\rA}
  \proofstepwidestar[]{\Id_A\colon A\to A}{\FormulaRefAuto{IdA}[thm]}
  \proofstepwidestar[]{F\circ \Id_A\colon A\to B}{%
    \FormulaRefAuto{CompositionDef}[def]{8,7}}
  \proofstepwide[]{F\circ \Id_A}{=}{F}{%
    \FormulaRefAuto{FunctionExtensionality}[thm]{9,7,6}}
\end{tabproofwide}

\FormulaThmDelta{%
F\circ F^{-1} = \Id_B
}{
  \DeltaRow{Mengen}{A \dsep B}
  \DeltaPrem{Bijektive Funktionen}{F\colon A\bij B}
}
\begin{tabproof}
    \proofstep{1}{y\in B}{\rA}
    \proofstep{1}{F(F^{-1}(y))=y}{\FormulaRefAuto{F\colon A\bij B\dsep y\in B \vdash F\bigl(F^{-1}(y)\bigr)=y}{1}}
    \proofstep{1}{(F\circ F^{-1})(y)=F(F^{-1}(y))}{\FormulaRefAuto{x\in A \vdash (G\circ F)(x)=G(F(x))}{2}}
    \proofstep{1}{(F\circ F^{-1})(y)=y}{\rIE{2,3}}
    \proofstep{1}{\Id_B(y)=y}{\FormulaRefAuto{x\in A \vdash \Id_A(x)=x}{1}}
    \proofstep{1}{(F\circ F^{-1})(y)=\Id_B(y)}{\rIE{4,5}}
    \proofstep{1}{\forall y\in B\,(F\circ F^{-1})(y)=\Id_B(y)}{\rUI{\rRI{1,6}}}
    \proofstep{}{F^{-1}\colon B\to A}{%
      \FormulaRefAuto{F\colon A\bij B\vdash F^{-1}\colon B\to A}[thm]}
    \proofstep{}{F\colon A\to B}{\FormulaRefAuto{Bijektive Funktion}}
    \proofstep{}{F\circ F^{-1}\colon B\to B}{%
      \FormulaRefAuto{CompositionDef}[def]{8,9}}
    \proofstep{}{\Id_B\colon B\to B}{\FormulaRefAuto{IdA}[thm]}
    \proofstep{1}{F\circ F^{-1}=\Id_B}{%
      \FormulaRefAuto{FunctionExtensionality}[thm]{10,11,7}}
\end{tabproof}

\FormulaThmDelta{%
F^{-1}\circ F = \Id_A
}{
  \DeltaRow{Mengen}{x \dsep A \dsep B}
  \DeltaPrem{Bijektive Funktionen}{F\colon A\bij B}
}
\begin{tabproof}
  \proofstep{1}{x\in A}{\rA}
  \proofstep{1}{F^{-1}(F(x))=x}{\FormulaRefAuto{x\in A \vdash F^{-1}(F(x))=x}{1}}
  \proofstep{1}{(F^{-1}\circ F)(x)=F^{-1}(F(x))}{\FormulaRefAuto{x\in A \vdash (G\circ F)(x)=G(F(x))}{1}}
  \proofstep{1}{(F^{-1}\circ F)(x)=x}{\rIE{2,3}}
  \proofstep{1}{\Id_A(x)=x}{\FormulaRefAuto{x\in A \vdash \Id_A(x)=x}{1}}
  \proofstep{1}{(F^{-1}\circ F)(x)=\Id_A(x)}{\rIE{4,5}}
  \proofstep{ }{\forall x\in A\,\bigl((F^{-1}\circ F)(x)=\Id_A(x)\bigr)}{\rUI{\rRI{1,6}}}
  \proofstep{}{F\colon A\to B}{\FormulaRefAuto{Bijektive Funktion}}
  \proofstep{}{F^{-1}\colon B\to A}{%
    \FormulaRefAuto{F\colon A\bij B\vdash F^{-1}\colon B\to A}[thm]}
  \proofstep{}{F^{-1}\circ F\colon A\to A}{%
    \FormulaRefAuto{CompositionDef}[def]{8,9}}
  \proofstep{}{\Id_A\colon A\to A}{\FormulaRefAuto{IdA}[thm]}
  \proofstep{ }{F^{-1}\circ F=\Id_A}{%
    \FormulaRefAuto{FunctionExtensionality}[thm]{10,11,7}}
\end{tabproof}

% ------------------------------------------------------------
% Hilfslemma: die breite Kettenrechnung auslagern
% ------------------------------------------------------------
\FormulaThmDelta{%
  x\in A \vdash (G^{-1}\circ F^{-1})\bigl((F\circ G)(x)\bigr)=x%
}{
  \DeltaRow{Mengen}{x \dsep A \dsep B \dsep C}
  \DeltaPrem{Bijektive Funktionen}{G\colon A\bij B \dsep F\colon B\bij C}
}
\begin{tabproofwide}
  \proofstepwidestar[1]{x\in A}{\rA}
  \proofstepwidestar[1]{G(x)\in B}{\FormulaRefAuto{F\colon A\to B\dsep x\in A\vdash F(x)\in B}{1}}
  \proofstepwidestar[1]{F(G(x))\in C}{\FormulaRefAuto{F\colon A\to B\dsep x\in A\vdash F(x)\in B}{2}}

  \proofstepwidestar[1]{(F\circ G)(x)=F(G(x))}{\FormulaRefAuto{x\in A \vdash (G\circ F)(x)=G(F(x))}{1}}

  \proofstepwidestar[1]{F^{-1}\bigl(F(G(x))\bigr)=G(x)}{\FormulaRefAuto{x\in A \vdash F^{-1}\bigl(F(x)\bigr)=x}{2}}
  \proofstepwidestar[1]{G^{-1}\bigl(G(x)\bigr)=x}{\FormulaRefAuto{x\in A \vdash F^{-1}\bigl(F(x)\bigr)=x}{1}}

  \proofstepwide[1]{%
    \begin{aligned}[t]
      (G^{-1}\circ F^{-1})\\
      \bigl((F\circ G)(x)\bigr)
    \end{aligned}
  }{=}{%
    \begin{aligned}[t]
      (G^{-1}\circ F^{-1})\\
      \bigl(F(G(x))\bigr)
    \end{aligned}
  }{\rIE{4,\rII}}

  \proofstepwide[1]{}{=}{G^{-1}\Bigl(F^{-1}\bigl(F(G(x))\bigr)\Bigr)}{\FormulaRefAuto{x\in A \vdash (G\circ F)(x)=G(F(x))}{3}}
  \proofstepwide[1]{}{=}{G^{-1}\bigl(G(x)\bigr)}{\rIE{5,8}}
  \proofstepwide[1]{}{=}{x}{\rIE{6,9}}

  \proofstepwidestar[1]{(G^{-1}\circ F^{-1})\bigl((F\circ G)(x)\bigr)=x}{\rChain{7,10}}
\end{tabproofwide}


% ------------------------------------------------------------
% Inversenformel für Kompositionen (kurzer Beweis: Block ersetzt)
% ------------------------------------------------------------
\FormulaThmDelta[Inverse der Komposition]{%
  G^{-1}\circ F^{-1} = (F\circ G)^{-1}%
}{
  \DeltaRow{Mengen}{x \dsep A \dsep B \dsep C}
  \DeltaPrem{Bijektive Funktionen}{G\colon A\bij B \dsep F\colon B\bij C}
}
\begin{tabproof}
  \proofstep{1}{x\in A}{\rA}

  \proofstep{1}{(G^{-1}\circ F^{-1})\bigl((F\circ G)(x)\bigr)=x}{%
    \FormulaRefAuto{x\in A \vdash (G^{-1}\circ F^{-1})\bigl((F\circ G)(x)\bigr)=x}{1}%
  }

  \proofstep{}{%
    \forall x\in A\,\Bigl((G^{-1}\circ F^{-1})\bigl((F\circ G)(x)\bigr)=x\Bigr)%
  }{%
    \rUI{\rRI{1,2}}%
  }

  \proofstep{}{%
    G^{-1}\circ F^{-1} = (F\circ G)^{-1}%
  }{%
    \FormulaRefAuto{\forall y\in B\,\bigl(F(G(y))=y\bigr) \vdash G = F^{-1}}{3}%
  }
\end{tabproof}




\subsection{Eigenschaften der Komponenten}

\FormulaThmDeltaR{%
x\in A\dsep y\in A\dsep F(x)=F(y)\vdash x=y
}{G\circ F\colon A\to C bijektiv\dsep x\in A\dsep y\in A\dsep F(x)=F(y)\vdash x=y}{
  \DeltaRow{Mengen}{A\dsep B\dsep C}
  \DeltaPrem{Funktionen}{F\colon A\to B\dsep G\colon B\to C}
  \DeltaPrem{Bijektive Funktionen}{G\circ F\colon A\bij C }
}
\begin{tabproof}
    \proofstep{1}{x\in A}{\rA}
    \proofstep{2}{y\in A}{\rA}
    \proofstep{3}{F(x)=F(y)}{\rA}
    \proofstep{}{G\colon B\to C}{\rA}
    \proofstep{1}{F(x)\in B}{\FormulaRefAuto{F\colon A\to B\dsep x\in A\vdash F(x)\in B}{1}}
    \proofstep{2}{F(y)\in B}{\FormulaRefAuto{F\colon A\to B\dsep x\in A\vdash F(x)\in B}{2}}
    \proofstep{1,2,3}{G(F(x))=G(F(y))}{\FormulaRefAuto{F\colon A\to B\dsep x\in A\dsep y\in A\dsep x=y\vdash F(x)=F(y)}{4,5,6,3}}
    \proofstep{1,2,3}{(G\circ F)(x)=(G\circ F)(y)}{\FormulaRefAuto{x\in A\dsep y\in A\dsep G(F(x))=G(F(y))\vdash (G\circ F)(x)=(G\circ F)(y)}{7}}
    \proofstep{1,2,3}{x=y}{\FormulaRefAuto{Injektivität}{1,2,8}}
\end{tabproof}

\FormulaThmDeltaR{%
y\in C\vdash \exists x\in B\; G(x)=y
}{G\circ F\colon A\to C bijektiv\dsep y\in C\vdash \exists x\in B\; G(x)=y}{
  \DeltaRow{Mengen}{A\dsep B\dsep C}
  \DeltaPrem{Funktionen}{F\colon A\to B\dsep G\colon B\to C}
  \DeltaPrem{Bijektive Funktionen}{G\circ F\colon A\bij C }
}
\begin{tabproof}
    \proofstep{1}{y\in C}{\rA}
    \proofstep{1}{\exists x\in A\, (G\circ F)(x)=y}{\FormulaRefAuto{y\in B\vdash \exists x\in A\; F(x)=y}{1}}
    \proofstep{3}{x\in A\land (G\circ F)(x)=y}{\rA}
    \proofstep{3}{x\in A}{\rAEa{3}}
    \proofstep{3}{(G\circ F)(x)=y}{\rAEb{3}}
    \proofstep{3}{(G\circ F)(x)=G(F(x))}{\FormulaRefAuto{x\in A \vdash (G\circ F)(x)=G(F(x))}{4}}
    \proofstep{3}{G(F(x))=y}{\FormulaRefAuto{a=b,\, a=c \vdash b=c}{6,5}}
    \proofstep{3}{F(x)\in B}{\FormulaRefAuto{F\colon A\to B\dsep x\in A\vdash F(x)\in B}{4}}
    \proofstep{3}{F(x)\in B\land G(F(x))=y}{\rAI{8,7}}
    \proofstep{3}{\exists x\in B\, G(x)=y}{\rEI{9}}
    \proofstep{1}{\exists x\in B\, G(x)=y}{\rEE{2,3,10}}
\end{tabproof}

\FormulaThmDeltaR{%
x\in B\dsep y\in B\dsep G(x)=G(y)\vdash x=y
}{G\circ F\colon A\to C bijektiv\dsep F\colon A\to B surjektiv\dsep x\in B\dsep y\in B\dsep G(x)=G(y)\vdash x=y}{
  \DeltaRow{Mengen}{A\dsep B\dsep C}
  \DeltaPrem{Funktionen}{G\colon B\to C}
  \DeltaPrem{Surjektive Funktionen}{F\colon A\sur B}
  \DeltaPrem{Bijektive Funktionen}{G\circ F\colon A\bij C }
}
\begin{tabproof}
    \proofstep{1}{x\in B}{\rA}
    \proofstep{2}{y\in B}{\rA}
    \proofstep{3}{G(x)=G(y)}{\rA}
    \proofstep{1}{\exists u\in A\,F(u)=x}{\FormulaRefAuto{y\in B\vdash \exists x\in A\; F(x)=y}{1}}
    \proofstep{2}{\exists v\in A\,F(v)=y}{\FormulaRefAuto{y\in B\vdash \exists x\in A\; F(x)=y}{2}}
    \proofstep{6}{u\in A\land F(u)=x}{\rA}
    \proofstep{6}{u\in A}{\rAEa{6}}
    \proofstep{6}{F(u)=x}{\rAEb{6}}
    \proofstep{9}{v\in A\land F(v)=y}{\rA}
    \proofstep{9}{v\in A}{\rAEa{9}}
    \proofstep{9}{F(v)=y}{\rAEb{9}}
    \proofstep{3,6,9}{G(F(u))=G(F(v))}{\FormulaRefAuto{a=b\dsep c=d \dsep P(b,d)\vdash P(a,c)}{8,11,3}}
    \proofstep{3,6,9}{(G\circ F)(u)=(G\circ F)(v)}{\FormulaRefAuto{x\in A\dsep y\in A\dsep G(F(x))=G(F(y))\vdash (G\circ F)(x)=(G\circ F)(y)}{7,10,12}}
    \proofstep{3,6,9}{u=v}{\FormulaRefAuto{x\in A \dsep y\in A \dsep (G\circ F)(x)=(G\circ F)(y) \vdash x=y}{7,10,13}}
    \proofstep{}{F\colon A\to B}{\FormulaRefAuto{Surjektive Funktion}}
    \proofstep{3,6,9}{F(u)=F(v)}{\FormulaRefAuto{F\colon A\to B\dsep x\in A\dsep y\in A\dsep x=y\vdash F(x)=F(y)}{15,7,10,14}}
    \proofstep{3,6,9}{x=y}{\rIE{11,\rIE{8,16}}}
    \proofstep{1,3,9}{x=y}{\rEE{4,6,17}}
    \proofstep{1,2,3}{x=y}{\rEE{5,9,18}}
\end{tabproof}

\FormulaThmDelta{%
 G\circ F\colon A\bij C\vdash F\colon A\inj B
}{
  \DeltaRow{Mengen}{A \dsep B \dsep C}
  \DeltaPrem{Funktionen}{F\colon A\to B \dsep G\colon B\to C}
}
\begin{tabproof}
    \proofstep{1}{ G\circ F\colon A\bij C}{\rA}
    \proofstep{1}{ \forall x,y\in A\, (F(x)=F(y)\rightarrow x=y)}{\FormulaRefAuto{G\circ F\colon A\to C bijektiv\dsep x\in A\dsep y\in A\dsep F(x)=F(y)\vdash x=y}{1}}
    \proofstep{1}{F\colon A\inj B}{\FormulaRefAuto{Injektive Funktion}{2}}
\end{tabproof}

\FormulaThmDelta{%
 G\circ F\colon A\bij C\vdash G\colon B\sur C
}{
  \DeltaRow{Mengen}{A \dsep B \dsep C}
  \DeltaPrem{Funktionen}{F\colon A\to B \dsep G\colon B\to C}
}
\begin{tabproof}
    \proofstep{1}{ G\circ F\colon A\bij C}{\rA}
    \proofstep{1}{\forall y\in C\exists x\in B\; G(x)=y}{\FormulaRefAuto{G\circ F\colon A\to C bijektiv\dsep y\in C\vdash \exists x\in B\; G(x)=y}{1}}
    \proofstep{1}{G\colon B\sur C}{\FormulaRefAuto{Surjektive Funktion}{2}}
\end{tabproof}


\FormulaThmDelta{%
 F\colon A\sur B\dsep G\circ F\colon A\bij C\vdash G\colon B\bij C
}{
  \DeltaRow{Mengen}{A \dsep B \dsep C}
  \DeltaPrem{Funktionen}{F\colon A\to B \dsep G\colon B\to C}
}
\begin{tabproof}
    \proofstep{1}{ F\colon A\sur B}{\rA}
    \proofstep{2}{ G\circ F\colon A\bij C}{\rA}
    \proofstep{2}{G\colon B\sur C}{\FormulaRefAuto{ G\circ F\colon A\bij C\vdash G\colon B\sur C}{2}}
    \proofstep{1,2}{\forall x,y\in B\, (G(x)=G(y)\rightarrow x=y)}{\FormulaRefAuto{G\circ F\colon A\to C bijektiv\dsep F\colon A\to B surjektiv\dsep x\in B\dsep y\in B\dsep G(x)=G(y)\vdash x=y}{1,2}}
    \proofstep{1,2}{G\colon B\inj C}{\FormulaRefAuto{Injektive Funktion}{4}}
    \proofstep{1,2}{G\colon B\bij C}{\FormulaRefAuto{F\colon A\inj B\dsep F\colon A\sur B\vdash F\colon A\bij B}{5,3}}
\end{tabproof}

\FormulaThmDelta{%
F\circ G = \Id_A\dsep G\circ F =\Id_B \vdash F\colon B\bij A
}{
  \DeltaRow{Mengen}{A}
  \DeltaPrem{Funktionen}{G\colon A\to B \dsep F\colon B\to A}
}
\begin{tabproofsplit}
\proofpart{Surjektivität}
\proofstep{1}{ F\circ G = \Id_A}{\rA}
    \proofstep{2}{ G\circ F =\Id_B}{\rA}
    \proofstep{}{\Id_A\colon A\bij A}{\FormulaRefAuto{\Id_A\colon A\bij A}}
    \proofstep{}{\Id_B\colon B\bij B}{\FormulaRefAuto{\Id_A\colon A\bij A}}
    \proofstep{1}{F\circ G\colon A\bij A}{\rIE{1,3}}
    \proofstep{2}{G\circ F\colon B\bij B}{\rIE{2,4}}
    \proofstep{1}{F\colon B\sur A}{\FormulaRefAuto{G\circ F\colon A\bij C\vdash G\colon B\sur C}{5}}

\closeproofpart

\proofpart{Injektivität}
  \setcounter{proofstepnr}{7}% Fortlaufende Schritte dieses Beweises.
\proofstep{2}{F\colon B\inj A}{\FormulaRefAuto{G\circ F\colon A\bij C\vdash F\colon A\inj B}{6}}

\closeproofpart

\proofpart{Bijektivität}
  \setcounter{proofstepnr}{8}% Fortlaufende Schritte dieses Beweises.
\proofstep{1,2}{F\colon B\bij A}{\FormulaRefAuto{F\colon A\inj B\dsep F\colon A\sur B\vdash F\colon A\bij B}{8,7}}

\closeproofpart
\end{tabproofsplit}

\FormulaThmDelta{%
F\circ G = \Id_A\dsep G\circ F =\Id_B \vdash G\colon A\bij B
}{
  \DeltaRow{Mengen}{A}
  \DeltaPrem{Funktionen}{G\colon A\to B \dsep F\colon B\to A}
}
\begin{tabproofsplit}
\proofpart{Surjektivität}
\proofstep{1}{ F\circ G = \Id_A}{\rA}
    \proofstep{2}{ G\circ F =\Id_B}{\rA}
    \proofstep{}{\Id_A\colon A\bij A}{\FormulaRefAuto{\Id_A\colon A\bij A}}
    \proofstep{}{\Id_B\colon B\bij B}{\FormulaRefAuto{\Id_A\colon A\bij A}}
    \proofstep{1}{F\circ G\colon A\bij A}{\rIE{1,3}}
    \proofstep{2}{G\circ F\colon B\bij B}{\rIE{2,4}}
    \proofstep{2}{G\colon A\sur B}{\FormulaRefAuto{G\circ F\colon A\bij C\vdash G\colon B\sur C}{6}}

\closeproofpart

\proofpart{Injektivität}
  \setcounter{proofstepnr}{7}% Fortlaufende Schritte dieses Beweises.
\proofstep{1}{G\colon A\inj B}{\FormulaRefAuto{G\circ F\colon A\bij C\vdash F\colon A\inj B}{5}}

\closeproofpart

\proofpart{Bijektivität}
  \setcounter{proofstepnr}{8}% Fortlaufende Schritte dieses Beweises.
\proofstep{1,2}{G\colon A\bij B}{\FormulaRefAuto{F\colon A\inj B\dsep F\colon A\sur B\vdash F\colon A\bij B}{8,7}}

\closeproofpart
\end{tabproofsplit}

\FormulaThmDelta{%
F\circ G = \Id_A\dsep G\circ F =\Id_B \vdash G=F^{-1}
}{
  \DeltaRow{Mengen}{A}
  \DeltaPrem{Funktionen}{G\colon A\to B \dsep F\colon B\to A}
}
\begin{tabproofwide}
    \proofstepwidestar[1]{ F\circ G = \Id_A}{\rA}
    \proofstepwidestar[2]{ G\circ F =\Id_B}{\rA}
    \proofstepwidestar[1,2]{F\colon B\bij A}{\FormulaRefAuto{F\circ G = \Id_A\dsep G\circ F =\Id_B \vdash F\colon B\bij A}{1,2}}
    \proofstepwidestar[]{G\colon A\to B}{\rA}
    \proofstepwidestar[1,2]{F\colon B\to A}{\FormulaRefAuto{Bijektive Funktion}{3}}
    \proofstepwidestar[1,2]{F^{-1}\colon A\to B}{%
      \FormulaRefAuto{F\colon A\bij B\vdash F^{-1}\colon B\to A}[thm]{3}}
    \proofstepwide[]{G}{=}{G\circ \Id_A}{\FormulaRefAuto{F \circ \Id_A = F}}
    \proofstepwidestar[1,2]{F\circ F^{-1}=\Id_A}{%
      \FormulaRefAuto{F\circ F^{-1} = \Id_B}{3}}
    \proofstepwide[1,2]{G\circ \Id_A}{=}{G\circ (F\circ F^{-1})}{\rIE{8,7}}
    \proofstepwide[1,2]{}{=}{(G\circ F)\circ F^{-1}}{%
      \FormulaRefAuto{F\colon A\to B\dsep G\colon B\to C\dsep H\colon C\to D
        \vdash H\circ (G\circ F) = (H\circ G)\circ F}[thm]{6,5,4}}
    \proofstepwide[1,2]{}{=}{\Id_B\circ F^{-1}}{\rIE{2,10}}
    \proofstepwide[1,2]{}{=}{F^{-1}}{\FormulaRefAuto{\Id_B\circ F = F}}
    \proofstepwide[1,2]{G}{=}{F^{-1}}{\rChain{7,12}}
\end{tabproofwide}

\section{Induzierte Potenzmengenabbildung}

\FormulaThmDelta[Injektivität der induzierten Potenzmengenabbildung]{%
  F\colon A\bij B \dsep X\in\powerset(A)\dsep Y\in\powerset(A)\dsep
  \PowMap{F}(X)=\PowMap{F}(Y) \vdash X=Y
}{%
  \DeltaRow{Mengen}{A\dsep B\dsep X\dsep Y}
  \DeltaRow{Funktionen}{F\dsep \PowMap{F}}
}
\begin{tabproofsplitwide}
  \proofpartwideindR[Hilfsschritt \(F[X]=F[Y]\)]{%
    F\colon A\bij B \dsep X\in\powerset(A)\dsep Y\in\powerset(A)\dsep
    \PowMap{F}(X)=\PowMap{F}(Y) \vdash F[X]=F[Y]
  }{%
    F\colon A\bij B \dsep X\in\powerset(A)\dsep Y\in\powerset(A)\dsep
    \PowMap{F}(X)=\PowMap{F}(Y) \vdash F[X]=F[Y]
  }
    \proofstepwidestar[1]{F\colon A\bij B}{\rA}
    \proofstepwidestar[2]{X\in\powerset(A)}{\rA}
    \proofstepwidestar[3]{Y\in\powerset(A)}{\rA}
    \proofstepwidestar[4]{\PowMap{F}(X)=\PowMap{F}(Y)}{\rA}
    \proofstepwidestar[1]{F\colon A\to B}{\FormulaRefAuto{Bijektive Funktion}{1}}
    \proofstepwidestar[1,2]{\PowMap{F}(X)=F[X]}{%
      \FormulaRefAuto{F\colon A\to B \dsep X\in\powerset(A)\vdash \PowMap{F}(X)=F[X]}{5,2}}
    \proofstepwidestar[1,3]{\PowMap{F}(Y)=F[Y]}{%
      \FormulaRefAuto{F\colon A\to B \dsep X\in\powerset(A)\vdash \PowMap{F}(X)=F[X]}{5,3}}
    \proofstepwidestar[1,2]{F[X]=\PowMap{F}(X)}{\FormulaRefAuto{a=b\vdash b=a}{6}}
    \proofstepwidestar[1,3,4]{\PowMap{F}(X)=F[Y]}{\rIE{4,7}}
    \proofstepwidestar[1,2,3,4]{F[X]=F[Y]}{\rIE{8,9}}
  \closeproofpartwideindR

  \proofpartwide{Schluss}
    \proofstepwidestar[1]{F\colon A\bij B}{\rA}
    \proofstepwidestar[2]{X\in\powerset(A)}{\rA}
    \proofstepwidestar[3]{Y\in\powerset(A)}{\rA}
    \proofstepwidestar[4]{\PowMap{F}(X)=\PowMap{F}(Y)}{\rA}
    \proofstepwidestar[1]{F\colon A\inj B}{\FormulaRefAuto{F\colon A\bij B \vdash F\colon A\inj B}{1}}
    \proofstepwidestar[1,2,3,4]{F[X]=F[Y]}{%
      \FormulaRefAuto{F\colon A\bij B \dsep X\in\powerset(A)\dsep Y\in\powerset(A)\dsep \PowMap{F}(X)=\PowMap{F}(Y) \vdash F[X]=F[Y]}{1,2,3,4}}
    \proofstepwidestar[1,2,3,4]{X=Y}{%
      \FormulaRefAuto{F\colon A\inj B \dsep X\in\powerset(A)\dsep Y\in\powerset(A)\dsep F[X]=F[Y] \vdash X=Y}{5,2,3,6}}
  \closeproofpartwide
\end{tabproofsplitwide}

\FormulaThmDelta[Surjektivität der induzierten Potenzmengenabbildung]{%
  F\colon A\bij B \dsep Y\in\powerset(B)\vdash \exists X\in\powerset(A)\,\PowMap{F}(X)=Y
}{%
  \DeltaRow{Mengen}{A\dsep B\dsep X\dsep Y}
  \DeltaRow{Funktionen}{F\dsep \PowMap{F}}
}
\begin{tabproof}
  \proofstep{1}{F\colon A\bij B}{\rA}
  \proofstep{2}{Y\in\powerset(B)}{\rA}
  \proofstep{2}{Y\subseteq B}{\FormulaRefAuto{x \in \powerset(A)\eqvdash x \subseteq A}{2}}
  \proofstep{1,2}{F^{-1}[Y]\subseteq A}{%
    \FormulaRefAuto{C\subseteq B\dsep F\colon A\bij B \vdash F^{-1}[C]\subseteq A}{3,1}}
  \proofstep{1,2}{F^{-1}[Y]\in\powerset(A)}{%
    \FormulaRefAuto{x \in \powerset(A)\eqvdash x \subseteq A}{4}}
  \proofstep{1}{F\colon A\to B}{\FormulaRefAuto{Bijektive Funktion}{1}}
  \proofstep{1,2}{F[F^{-1}[Y]]=Y}{%
    \FormulaRefAuto{Y\subseteq B \dsep F\colon A\bij B \vdash F[F^{-1}[Y]]=Y}{3,1}}
  \proofstep{1,2}{\PowMap{F}(F^{-1}[Y])=F[F^{-1}[Y]]}{%
    \FormulaRefAuto{F\colon A\to B \dsep X\in\powerset(A)\vdash \PowMap{F}(X)=F[X]}{6,5}}
  \proofstep{1,2}{\PowMap{F}(F^{-1}[Y])=Y}{%
    \FormulaRefAuto{a=b,\, b=c \vdash a=c}{8,7}}
  \proofstep{1,2}{\exists X\in\powerset(A)\,\PowMap{F}(X)=Y}{\rEI{\rAI{5,9}}}
\end{tabproof}

\FormulaThmDelta[Induzierte Potenzmengenabbildung einer Bijektion]{%
  F\colon A\bij B \vdash \PowMap{F}\colon \powerset(A)\bij \powerset(B)
}{%
  \DeltaRow{Mengen}{A\dsep B\dsep X\dsep Y}
  \DeltaRow{Funktionen}{F\dsep \PowMap{F}}
}
\begin{tabproofsplit}
\proofpart{Injektivität}
\proofstep{1}{F\colon A\bij B}{\rA}
  \proofstep{1}{F\colon A\to B}{\FormulaRefAuto{Bijektive Funktion}{1}}
  \proofstep{1}{\PowMap{F}\colon \powerset(A)\to \powerset(B)}{%
    \FormulaRefAuto{F\colon A\to B \vdash \PowMap{F}\colon \powerset(A)\to \powerset(B)}{2}}
  \proofstep{1}{\forall X,Y\in\powerset(A)\,(\PowMap{F}(X)=\PowMap{F}(Y)\rightarrow X=Y)}{%
      \FormulaRefAuto{F\colon A\bij B \dsep X\in\powerset(A)\dsep Y\in\powerset(A)\dsep \PowMap{F}(X)=\PowMap{F}(Y) \vdash X=Y}{1}}
  \proofstep{1}{\PowMap{F}\colon \powerset(A)\inj \powerset(B)}{%
    \FormulaRefAuto{Injektive Funktion}{3,4}}

\closeproofpart

\proofpart{Surjektivität}
  \setcounter{proofstepnr}{5}% Fortlaufende Schritte dieses Beweises.
\proofstep{1}{\forall Y\in\powerset(B)\,\exists X\in\powerset(A)\,(\PowMap{F}(X)=Y)}{%
      \FormulaRefAuto{F\colon A\bij B \dsep Y\in\powerset(B)\vdash \exists X\in\powerset(A)\,\PowMap{F}(X)=Y}{1}}
\proofstep{1}{\PowMap{F}\colon \powerset(A)\sur \powerset(B)}{%
    \FormulaRefAuto{Surjektive Funktion}{3,6}}

\closeproofpart

\proofpart{Bijektivität}
  \setcounter{proofstepnr}{7}% Fortlaufende Schritte dieses Beweises.
\proofstep{1}{\PowMap{F}\colon \powerset(A)\bij \powerset(B)}{%
    \FormulaRefAuto{F\colon A\inj B\dsep F\colon A\sur B \vdash F\colon A\bij B}{5,7}}

\closeproofpart
\end{tabproofsplit}

\subsection{Rückwirkung der induzierten Potenzmengenabbildung}

Die folgende Rückrichtung betrifft nicht irgendeine Bijektion zwischen zwei
Potenzmengen.  Entscheidend ist, dass die Abbildung von einer bereits
gegebenen Grundabbildung \(F\) induziert wird.  Einermengen machen dann die
Injektivität und Surjektivität von \(F\) in der Potenzmengenabbildung
sichtbar.

\FormulaThmDeltaK[Injektivität wird von der induzierten
Potenzmengenabbildung reflektiert]{%
  \begin{aligned}[t]
    &F\colon A\to B\dsep
      \PowMap{F}\colon\powerset(A)\inj\powerset(B)\dsep{}\\[-2pt]
    &x,y\in A\dsep F(x)=F(y)\vdash x=y
  \end{aligned}%
}{PowerMapInjectivityPointwiseReflection}{%
  \DeltaRow{Mengen}{A\dsep B\dsep x\dsep y}
  \DeltaRow{Funktionen}{F\dsep\PowMap{F}}
}
\begin{tabproofwide}
  \proofstepwidestar[1]{F\colon A\to B}{\rA}
  \proofstepwidestar[2]{%
    \PowMap{F}\colon\powerset(A)\inj\powerset(B)}{\rA}
  \proofstepwidestar[3]{x\in A}{\rA}
  \proofstepwidestar[4]{y\in A}{\rA}
  \proofstepwidestar[5]{F(x)=F(y)}{\rA}
  \proofstepwidestar[3]{\{ x \} \subseteq A}{%
      \FormulaRefAuto{a\in A\vdash\{a\}\subseteq A}{3}}
  \proofstepwidestar[3]{\{x\}\in\powerset(A)}{%
    \FormulaRefAuto{x\in\powerset(A)\eqvdash x\subseteq A}{%
      6}}
  \proofstepwidestar[4]{\{ y \} \subseteq A}{%
      \FormulaRefAuto{a\in A\vdash\{a\}\subseteq A}{4}}
  \proofstepwidestar[4]{\{y\}\in\powerset(A)}{%
    \FormulaRefAuto{x\in\powerset(A)\eqvdash x\subseteq A}{%
      8}}
  \proofstepwidestar[1,3]{F [ \{ x \} ] = \{ F ( x ) \}}{%
      \FormulaRefAuto{x\in A\dsep F\colon A\to B
        \vdash F[\{x\}]=\{F(x)\}}{3,1}}
  \proofstepwidestar[1,3]{\PowMap F ( \{ x \} ) = F [ \{ x \} ]}{%
      \FormulaRefAuto{F\colon A\to B\dsep X\in\powerset(A)
        \vdash\PowMap{F}(X)=F[X]}{1,7}}
  \proofstepwidestar[1,3]{\PowMap{F}(\{x\})=\{F(x)\}}{%
    \FormulaRefAuto{a=b,\, b=c \vdash a=c}{%
      11,
      10}}
  \proofstepwidestar[1,4]{F [ \{ y \} ] = \{ F ( y ) \}}{%
      \FormulaRefAuto{x\in A\dsep F\colon A\to B
        \vdash F[\{x\}]=\{F(x)\}}{4,1}}
  \proofstepwidestar[1,4]{\PowMap F ( \{ y \} ) = F [ \{ y \} ]}{%
      \FormulaRefAuto{F\colon A\to B\dsep X\in\powerset(A)
        \vdash\PowMap{F}(X)=F[X]}{1,9}}
  \proofstepwidestar[1,4]{\PowMap{F}(\{y\})=\{F(y)\}}{%
    \FormulaRefAuto{a=b,\, b=c \vdash a=c}{%
      14,
      13}}
  \proofstepwidestar[5]{\{F(x)\}=\{F(y)\}}{%
    \FormulaRefAuto{\{a\}=\{b\}\eqvdash a=b}{5}}
  \proofstepwidestar[1,3,4,5]{%
    \PowMap{F}(\{x\})=\{F(y)\}}{%
    \FormulaRefAuto{a=b\dsep b=c\vdash a=c}{12,16}}
  \proofstepwidestar[1,4]{\{ F ( y ) \} = \PowMap F ( \{ y \} )}{%
      \FormulaRefAuto{a=b\vdash b=a}{15}}
  \proofstepwidestar[1,3,4,5]{%
    \PowMap{F}(\{x\})=\PowMap{F}(\{y\})}{%
    \FormulaRefAuto{a=b\dsep b=c\vdash a=c}{%
      17,18}}
  \proofstepwidestar[1,2,3,4,5]{\{x\}=\{y\}}{%
    \FormulaRefAuto{F\colon A\inj B\dsep x\in A\dsep y\in A\dsep
      F(x)=F(y)\vdash x=y}{2,7,9,19}}
  \proofstepwidestar[1,2,3,4,5]{x=y}{%
    \FormulaRefAuto{\{a\}=\{b\}\eqvdash a=b}{20}}
\end{tabproofwide}

\FormulaThmDeltaK[Injektivität der induzierten Potenzmengenabbildung
erzwingt Injektivität]{%
  F\colon A\to B\dsep
  \PowMap{F}\colon\powerset(A)\inj\powerset(B)
  \vdash F\colon A\inj B%
}{PowerMapInjectiveReflectsInjectivity}{%
  \DeltaRow{Mengen}{A\dsep B\dsep x\dsep y}
  \DeltaRow{Funktionen}{F\dsep\PowMap{F}}
}
\begin{tabproof}
  \proofstep{1}{F\colon A\to B}{\rA}
  \proofstep{2}{\PowMap{F}\colon\powerset(A)\inj\powerset(B)}{\rA}
  \proofstep{1,2}{%
    \forall x,y\in A\,\bigl(F(x)=F(y)\rightarrow x=y\bigr)}{%
    \FormulaRefAuto{PowerMapInjectivityPointwiseReflection}[thm]{1,2}}
  \proofstep{1,2}{F\colon A\inj B}{%
    \FormulaRefAuto{Injektive Funktion}[def]{1,3}}
\end{tabproof}

\FormulaThmDeltaK[Surjektivität wird von der induzierten
Potenzmengenabbildung reflektiert]{%
  \begin{aligned}[t]
    &F\colon A\to B\dsep
      \PowMap{F}\colon\powerset(A)\sur\powerset(B)\dsep y\in B\vdash{}\\[-2pt]
    &\exists x\in A\,F(x)=y
  \end{aligned}%
}{PowerMapSurjectivityPointwiseReflection}{%
  \DeltaRow{Mengen}{A\dsep B\dsep X\dsep x\dsep y}
  \DeltaRow{Funktionen}{F\dsep\PowMap{F}}
}
\begin{tabproofwide}
  \proofstepwidestar[1]{F\colon A\to B}{\rA}
  \proofstepwidestar[2]{%
    \PowMap{F}\colon\powerset(A)\sur\powerset(B)}{\rA}
  \proofstepwidestar[3]{y\in B}{\rA}
  \proofstepwidestar[3]{\{ y \} \subseteq B}{%
      \FormulaRefAuto{a\in A\vdash\{a\}\subseteq A}{3}}
  \proofstepwidestar[3]{\{y\}\in\powerset(B)}{%
    \FormulaRefAuto{x\in\powerset(A)\eqvdash x\subseteq A}{%
      4}}
  \proofstepwidestar[2,3]{%
    \exists X\in\powerset(A)\,\PowMap{F}(X)=\{y\}}{%
    \FormulaRefAuto{y\in B\vdash\exists x\in A\;F(x)=y}{2,5}}
  \proofstepwidestar[7]{%
    X\in\powerset(A)\land\PowMap{F}(X)=\{y\}}{\rA}
  \proofstepwidestar[7]{X\in\powerset(A)}{\rAEa{7}}
  \proofstepwidestar[7]{\PowMap{F}(X)=\{y\}}{\rAEb{7}}
  \proofstepwidestar[1,7]{\PowMap{F}(X)=F[X]}{%
    \FormulaRefAuto{F\colon A\to B\dsep X\in\powerset(A)
      \vdash\PowMap{F}(X)=F[X]}{1,8}}
  \proofstepwidestar[1,7]{F[X]=\{y\}}{%
    \FormulaRefAuto{a=b\dsep a=c\vdash b=c}{10,9}}
  \proofstepwidestar[]{y\in\{y\}}{%
    \FormulaRefAuto{x\in\{a\}\eqvdash x=a}{\rII}}
  \proofstepwidestar[1,7]{y\in F[X]}{\rIE{11,12}}
  \proofstepwidestar[7]{X\subseteq A}{%
    \FormulaRefAuto{x\in\powerset(A)\eqvdash x\subseteq A}{8}}
  \proofstepwidestar[1,7]{\exists x\in X\,y=F(x)}{%
    \FormulaRefAuto{C\subseteq A\dsep F\colon A\to B\dsep
      y\in F[C]\vdash\exists x\in C\,y=F(x)}{14,1,13}}
  \proofstepwidestar[16]{x\in X\land y=F(x)}{\rA}
  \proofstepwidestar[7,16]{x\in A}{%
    \FormulaRefAuto{A\subseteq B\dsep x\in A\vdash x\in B}{14,\rAEa{16}}}
  \proofstepwidestar[16]{F(x)=y}{%
    \FormulaRefAuto{a=b\vdash b=a}{\rAEb{16}}}
  \proofstepwidestar[7,16]{\exists x\in A\,F(x)=y}{%
    \rEI{\rAI{17,18}}}
  \proofstepwidestar[1,7]{\exists x\in A\,F(x)=y}{%
    \rEE{15,16,19}}
  \proofstepwidestar[1,2,3]{\exists x\in A\,F(x)=y}{%
    \rEE{6,7,20}}
\end{tabproofwide}

\FormulaThmDeltaK[Surjektivität der induzierten Potenzmengenabbildung
erzwingt Surjektivität]{%
  F\colon A\to B\dsep
  \PowMap{F}\colon\powerset(A)\sur\powerset(B)
  \vdash F\colon A\sur B%
}{PowerMapSurjectiveReflectsSurjectivity}{%
  \DeltaRow{Mengen}{A\dsep B\dsep x\dsep y}
  \DeltaRow{Funktionen}{F\dsep\PowMap{F}}
}
\begin{tabproof}
  \proofstep{1}{F\colon A\to B}{\rA}
  \proofstep{2}{\PowMap{F}\colon\powerset(A)\sur\powerset(B)}{\rA}
  \proofstep{1,2}{\forall y\in B\,\exists x\in A\,F(x)=y}{%
    \FormulaRefAuto{PowerMapSurjectivityPointwiseReflection}[thm]{1,2}}
  \proofstep{1,2}{F\colon A\sur B}{%
    \FormulaRefAuto{Surjektive Funktion}[def]{1,3}}
\end{tabproof}

\FormulaThmDeltaK[Bijektivität der induzierten Potenzmengenabbildung
erzwingt Bijektivität]{%
  F\colon A\to B\dsep
  \PowMap{F}\colon\powerset(A)\bij\powerset(B)
  \vdash F\colon A\bij B%
}{PowerMapBijectiveReflectsBijectivity}{%
  \DeltaRow{Mengen}{A\dsep B}
  \DeltaRow{Funktionen}{F\dsep\PowMap{F}}
}
\begin{tabproofsplit}
\proofpart{Injektivität}
\proofstep{1}{F\colon A\to B}{\rA}
  \proofstep{2}{\PowMap{F}\colon\powerset(A)\bij\powerset(B)}{\rA}
  \proofstep{2}{\PowMap{F}\colon\powerset(A)\inj\powerset(B)}{%
    \FormulaRefAuto{F\colon A\bij B\vdash F\colon A\inj B}{2}}
  \proofstep{2}{\PowMap{F}\colon\powerset(A)\sur\powerset(B)}{%
    \FormulaRefAuto{F\colon A\bij B\vdash F\colon A\sur B}{2}}
  \proofstep{1,2}{F\colon A\inj B}{%
    \FormulaRefAuto{PowerMapInjectiveReflectsInjectivity}[thm]{1,3}}

\closeproofpart

\proofpart{Surjektivität}
  \setcounter{proofstepnr}{5}% Fortlaufende Schritte dieses Beweises.
\proofstep{1,2}{F\colon A\sur B}{%
    \FormulaRefAuto{PowerMapSurjectiveReflectsSurjectivity}[thm]{1,4}}

\closeproofpart

\proofpart{Bijektivität}
  \setcounter{proofstepnr}{6}% Fortlaufende Schritte dieses Beweises.
\proofstep{1,2}{F\colon A\bij B}{%
    \FormulaRefAuto{F\colon A\inj B\dsep F\colon A\sur B
      \vdash F\colon A\bij B}{5,6}}

\closeproofpart
\end{tabproofsplit}

\begin{remark}[Externer Hinweis: keine allgemeine Kürzungsregel]
Der eben bewiesene Satz setzt wesentlich voraus, dass die Bijektion der
Potenzmengen gleich \(\PowMap{F}\) ist.  Aus der bloßen Existenz irgendeiner
Bijektion
\[
  G\colon\powerset(A)\bij\powerset(B)
\]
folgt in der üblichen Mengenlehre ZFC nicht allgemein, dass eine Bijektion von
\(A\) nach \(B\) existiert.  Dies ist hier ein externer mengentheoretischer
Hinweis, kein im vorliegenden Begriffsaufbau verwendeter Satz: Unter der
üblichen relativen Konsistenzvoraussetzung gibt es Modelle von ZFC, in denen
verschieden mächtige unendliche Mengen gleichmächtige Potenzmengen besitzen,
und Modelle von ZFC, in denen dies nicht vorkommt.  Die Unabhängigkeit folgt
aus einem Satz von
\href{https://doi.org/10.1016/0003-4843(70)90012-4}{Easton}.  Das
Diagonalargument zeigt nur, dass keine Abbildung von einer Menge auf ihre
gesamte Potenzmenge surjektiv ist; eine Kürzung zweier Potenzmengen liefert es
nicht.
\end{remark}

\FormulaThmDeltaK[Bijektivität auf den nichtleeren Potenzmengen]{%
  \begin{aligned}[t]
    &f\colon A\bij B\vdash{}\\[-2pt]
    &\PowMapNE{f}\colon\PowNE{A}\bij\PowNE{B}
  \end{aligned}%
}{NonemptyPowerMapBijective}{%
  \DeltaRow{Mengen}{A\dsep B}
  \DeltaRow{Funktionen}{f\dsep\PowMap{f}\dsep\PowMapNE{f}}
}
\begin{tabproofwide}
  \proofstepwidestar[1]{f\colon A\bij B}{\rA}
  \proofstepwidestar[1]{f\colon A\to B}{%
    \FormulaRefAuto{Bijektive Funktion}[def]{1}}
  \proofstepwidestar[1]{\PowMap{f}\colon\powerset(A)\bij\powerset(B)}{%
    \FormulaRefAuto{F\colon A\bij B\vdash
      \PowMap{F}\colon\powerset(A)\bij\powerset(B)}[thm]{1}}
  \proofstepwidestar[]{\PowNE{B}\subseteq\powerset(B)}{%
    \FormulaRefAuto{A\setminus B\subseteq A}[thm]}
  \proofstepwidestar[1]{%
    \PowMap{f}\restriction_{(\PowMap{f})^{-1}[\PowNE{B}]}
    \colon(\PowMap{f})^{-1}[\PowNE{B}]\bij\PowNE{B}}{%
    \FormulaRefAuto{F\colon A\bij B\dsep T\subseteq B
      \vdash F\restriction_{F^{-1}[T]}\colon F^{-1}[T]\bij T}[thm]{3,4}}
  \proofstepwidestar[1]{(\PowMap{f})^{-1}[\PowNE{B}]=\PowNE{A}}{%
    \FormulaRefAuto{NonemptyPowerMapPreimage}[thm]{2}}
  \proofstepwidestar[1]{%
    \PowMap{f}\restriction_{\PowNE{A}}
    \colon\PowNE{A}\bij\PowNE{B}}{\rIE{6,5}}
  \proofstepwidestar[1]{%
    \PowMapNE{f}=\PowMap{f}\restriction_{\PowNE{A}}}{%
    \FormulaRefAuto{NonemptyPowerMapDef}[def]{2}}
  \proofstepwidestar[1]{%
    \PowMapNE{f}\colon\PowNE{A}\bij\PowNE{B}}{\rIE{8,7}}
\end{tabproofwide}

\subsection{Entfernte Teilmengenfamilien}

Das punktweise Kriterium für bijektive Einschränkungen macht hier die
Kernidee unmittelbar sichtbar: Man setzt
\(S=\powerset(A)\setminus C\), \(T=\powerset(B)\setminus C\) und betrachtet
die Bijektion \(\PowMap{F}\). Negiert man die beiden
Differenzmitgliedschaften, so bleibt genau die Erhaltung der Zugehörigkeit
zu \(C\).

\FormulaThmDeltaK[Exakte \(C\)-Verträglichkeit der eingeschränkten
Potenzmengenabbildung]{%
  \begin{aligned}[t]
    &F\colon A\bij B\vdash{}\\[-2pt]
    &\Bigl(
      \PowMap{F}\restriction_{\powerset(A)\setminus C}
        \colon\powerset(A)\setminus C\bij\powerset(B)\setminus C
      \leftrightarrow{}\\[-2pt]
    &\hspace{34mm}
      \forall X\in\powerset(A)\,
        \bigl(X\in C\leftrightarrow F[X]\in C\bigr)
    \Bigr)
  \end{aligned}%
}{RestrictedPowerMapBijectiveIffCCompatible}{%
  \DeltaRow{Mengen}{A\dsep B\dsep C\dsep X\dsep Y}
  \DeltaRow{Funktionen}{F\dsep\PowMap{F}}
}
\begin{tabproofsplitwide}
  \proofpartwideindR[Aus der eingeschränkten Bijektivität]{%
    \begin{aligned}[t]
      &F\colon A\bij B\dsep
      \PowMap{F}\restriction_{\powerset(A)\setminus C}
        \colon\powerset(A)\setminus C\bij\powerset(B)\setminus C\\[-2pt]
      &\vdash
      \forall X\in\powerset(A)\,
        \bigl(X\in C\leftrightarrow F[X]\in C\bigr)
    \end{aligned}%
  }{%
    \begin{aligned}[t]
      &F\colon A\bij B\dsep
      \PowMap{F}\restriction_{\powerset(A)\setminus C}
        \colon\powerset(A)\setminus C\bij\powerset(B)\setminus C\\[-2pt]
      &\vdash
      \forall X\in\powerset(A)\,
        \bigl(X\in C\leftrightarrow F[X]\in C\bigr)
    \end{aligned}%
  }[RestrictedPowerMapBijectiveImpliesCCompatibility]
    \proofstepwidestar[1]{F\colon A\bij B}{\rA}
    \proofstepwidestar[2]{%
      \PowMap{F}\restriction_{\powerset(A)\setminus C}
        \colon\powerset(A)\setminus C\bij\powerset(B)\setminus C}{\rA}
    \proofstepwidestar[1]{F\colon A\to B}{%
      \FormulaRefAuto{Bijektive Funktion}[def]{1}}
    \proofstepwidestar[1]{%
      \PowMap{F}\colon\powerset(A)\bij\powerset(B)}{%
      \FormulaRefAuto{F\colon A\bij B\vdash
        \PowMap{F}\colon\powerset(A)\bij\powerset(B)}[thm]{1}}
    \proofstepwidestar[1]{%
      \PowMap{F}\colon\powerset(A)\to\powerset(B)}{%
      \FormulaRefAuto{Bijektive Funktion}[def]{4}}
    \proofstepwidestar[]{\powerset(A)\setminus C\subseteq\powerset(A)}{%
      \FormulaRefAuto{A\setminus B\subseteq A}[thm]}
    \proofstepwidestar[]{\powerset(B)\setminus C\subseteq\powerset(B)}{%
      \FormulaRefAuto{A\setminus B\subseteq A}[thm]}
    \proofstepwidestar[1,2]{%
      \begin{aligned}[t]
        &\forall X\in\powerset(A)\,\\[-2pt]
        &\quad\Bigl(
          X\in\powerset(A)\setminus C
          \leftrightarrow
          \PowMap{F}(X)\in\powerset(B)\setminus C
        \Bigr)
      \end{aligned}}{%
      \FormulaRefAuto{BijectiveRestrictionImpliesPointwiseCompatibility}[thm]{%
        4,6,7,2}}
    \proofstepwidestar[9]{X\in\powerset(A)}{\rA}
    \proofstepwidestar[1,9]{\PowMap{F}(X)\in\powerset(B)}{%
      \FormulaRefAuto{F\colon A\to B\dsep x\in A\vdash F(x)\in B}[thm]{%
        5,9}}
    \proofstepwidestar[9]{( X \notin \powerset ( A ) \setminus C \leftrightarrow X \in C )}{%
      \FormulaRefAuto{RelCompNotMembership}[thm]{9}}
    \proofstepwide[9]{X\in C}{\leftrightarrow}{%
      X\notin\powerset(A)\setminus C}{%
      \FormulaRefAuto{P\leftrightarrow Q\vdash Q\leftrightarrow P}[thm]{%
        11}}
    \proofstepwide[1,2,9]{}{\leftrightarrow}{%
      \PowMap{F}(X)\notin\powerset(B)\setminus C}{%
      \FormulaRefAuto{P \leftrightarrow Q \eqvdash
        \neg P \leftrightarrow \neg Q}[thm]{%
        \rRE{\rUE{8},9}}}
    \proofstepwide[1,9]{}{\leftrightarrow}{\PowMap{F}(X)\in C}{%
      \FormulaRefAuto{RelCompNotMembership}[thm]{10}}
    \proofstepwide[1,2,9]{X\in C}{\leftrightarrow}{%
      \PowMap{F}(X)\in C}{\rChain{12,14}}
    \proofstepwidestar[1,9]{\PowMap{F}(X)=F[X]}{%
      \FormulaRefAuto{F\colon A\to B\dsep X\in\powerset(A)
        \vdash\PowMap{F}(X)=F[X]}[thm]{3,9}}
    \proofstepwidestar[1,2,9]{X\in C\leftrightarrow F[X]\in C}{%
      \rIE{16,15}}
    \proofstepwidestar[1,2]{%
      \forall X\in\powerset(A)\,
        \bigl(X\in C\leftrightarrow F[X]\in C\bigr)}{%
      \rUI{\rRI{9,17}}}
  \closeproofpartwideindR

  \proofpartwideindR[Aus der \(C\)-Verträglichkeit]{%
    \begin{aligned}[t]
      &F\colon A\bij B\dsep
      \forall X\in\powerset(A)\,
        \bigl(X\in C\leftrightarrow F[X]\in C\bigr)\\[-2pt]
      &\vdash
      \PowMap{F}\restriction_{\powerset(A)\setminus C}
        \colon\powerset(A)\setminus C\bij\powerset(B)\setminus C
    \end{aligned}%
  }{%
    \begin{aligned}[t]
      &F\colon A\bij B\dsep
      \forall X\in\powerset(A)\,
        \bigl(X\in C\leftrightarrow F[X]\in C\bigr)\\[-2pt]
      &\vdash
      \PowMap{F}\restriction_{\powerset(A)\setminus C}
        \colon\powerset(A)\setminus C\bij\powerset(B)\setminus C
    \end{aligned}%
  }[CCompatibilityImpliesRestrictedPowerMapBijective]
    \proofstepwidestar[1]{F\colon A\bij B}{\rA}
    \proofstepwidestar[2]{%
      \forall X\in\powerset(A)\,
        \bigl(X\in C\leftrightarrow F[X]\in C\bigr)}{\rA}
    \proofstepwidestar[1]{F\colon A\to B}{%
      \FormulaRefAuto{Bijektive Funktion}[def]{1}}
    \proofstepwidestar[1]{%
      \PowMap{F}\colon\powerset(A)\bij\powerset(B)}{%
      \FormulaRefAuto{F\colon A\bij B\vdash
        \PowMap{F}\colon\powerset(A)\bij\powerset(B)}[thm]{1}}
    \proofstepwidestar[1]{%
      \PowMap{F}\colon\powerset(A)\to\powerset(B)}{%
      \FormulaRefAuto{Bijektive Funktion}[def]{4}}
    \proofstepwidestar[]{\powerset(A)\setminus C\subseteq\powerset(A)}{%
      \FormulaRefAuto{A\setminus B\subseteq A}[thm]}
    \proofstepwidestar[]{\powerset(B)\setminus C\subseteq\powerset(B)}{%
      \FormulaRefAuto{A\setminus B\subseteq A}[thm]}
    \proofstepwidestar[8]{X\in\powerset(A)}{\rA}
    \proofstepwidestar[1,2,8]{X\in C\leftrightarrow F[X]\in C}{%
      \rRE{\rUE{2},8}}
    \proofstepwidestar[1,8]{\PowMap F ( X ) = F [ X ]}{%
      \FormulaRefAuto{F\colon A\to B\dsep X\in\powerset(A)
          \vdash\PowMap{F}(X)=F[X]}[thm]{3,8}}
    \proofstepwidestar[1,8]{F[X]=\PowMap{F}(X)}{%
      \FormulaRefAuto{a=b\vdash b=a}[thm]{%
        10}}
    \proofstepwidestar[1,8]{\PowMap{F}(X)\in\powerset(B)}{%
      \FormulaRefAuto{F\colon A\to B\dsep x\in A\vdash F(x)\in B}[thm]{%
        5,8}}
    \proofstepwide[8]{X\notin\powerset(A)\setminus C}{\leftrightarrow}{%
      X\in C}{\FormulaRefAuto{RelCompNotMembership}[thm]{8}}
    \proofstepwide[1,2,8]{}{\leftrightarrow}{\PowMap{F}(X)\in C}{%
      \rIE{11,9}}
    \proofstepwidestar[1,8]{( \PowMap F ( X ) \notin \powerset ( B ) \setminus C \leftrightarrow \PowMap F ( X ) \in C )}{%
      \FormulaRefAuto{RelCompNotMembership}[thm]{12}}
    \proofstepwide[1,8]{\PowMap{F}(X)\in C}{\leftrightarrow}{%
      \PowMap{F}(X)\notin\powerset(B)\setminus C}{%
      \FormulaRefAuto{P\leftrightarrow Q\vdash Q\leftrightarrow P}[thm]{%
        15}}
    \proofstepwide[1,2,8]{X\notin\powerset(A)\setminus C}{%
      \leftrightarrow}{%
      \PowMap{F}(X)\notin\powerset(B)\setminus C}{\rChain{13,16}}
    \proofstepwidestar[1,2,8]{%
      X\in\powerset(A)\setminus C
      \leftrightarrow
      \PowMap{F}(X)\in\powerset(B)\setminus C}{%
      \FormulaRefAuto{P \leftrightarrow Q \eqvdash
        \neg P \leftrightarrow \neg Q}[thm]{17}}
    \proofstepwidestar[1,2]{%
      \begin{aligned}[t]
        &\forall X\in\powerset(A)\,\\[-2pt]
        &\quad\Bigl(
          X\in\powerset(A)\setminus C
          \leftrightarrow
          \PowMap{F}(X)\in\powerset(B)\setminus C
        \Bigr)
      \end{aligned}}{%
      \rUI{\rRI{8,18}}}
    \proofstepwidestar[1,2]{%
      \PowMap{F}\restriction_{\powerset(A)\setminus C}
        \colon\powerset(A)\setminus C\bij\powerset(B)\setminus C}{%
      \FormulaRefAuto{PointwiseCompatibilityImpliesBijectiveRestriction}[thm]{%
        4,6,7,19}}
  \closeproofpartwideindR

  \proofpartwide{Schluss}
    \proofstepwidestar[1]{F\colon A\bij B}{\rA}
    \proofstepwidestar[2]{%
      \PowMap{F}\restriction_{\powerset(A)\setminus C}
        \colon\powerset(A)\setminus C\bij\powerset(B)\setminus C}{\rA}
    \proofstepwidestar[1,2]{%
      \forall X\in\powerset(A)\,
        \bigl(X\in C\leftrightarrow F[X]\in C\bigr)}{%
      \FormulaRefAuto{RestrictedPowerMapBijectiveImpliesCCompatibility}[thm]{%
        1,2}}
    \proofstepwidestar[1]{%
      \begin{aligned}[t]
        &\PowMap{F}\restriction_{\powerset(A)\setminus C}
          \colon\powerset(A)\setminus C\bij\powerset(B)\setminus C\\[-2pt]
        &\quad\rightarrow
          \forall X\in\powerset(A)\,
            \bigl(X\in C\leftrightarrow F[X]\in C\bigr)
      \end{aligned}}{\rRI{2,3}}
    \proofstepwidestar[5]{%
      \forall X\in\powerset(A)\,
        \bigl(X\in C\leftrightarrow F[X]\in C\bigr)}{\rA}
    \proofstepwidestar[1,5]{%
      \PowMap{F}\restriction_{\powerset(A)\setminus C}
        \colon\powerset(A)\setminus C\bij\powerset(B)\setminus C}{%
      \FormulaRefAuto{CCompatibilityImpliesRestrictedPowerMapBijective}[thm]{%
        1,5}}
    \proofstepwidestar[1]{%
      \begin{aligned}[t]
        &\forall X\in\powerset(A)\,
          \bigl(X\in C\leftrightarrow F[X]\in C\bigr)\\[-2pt]
        &\quad\rightarrow
          \PowMap{F}\restriction_{\powerset(A)\setminus C}
            \colon\powerset(A)\setminus C\bij\powerset(B)\setminus C
      \end{aligned}}{\rRI{5,6}}
    \proofstepwidestar[1]{%
      \begin{aligned}[t]
        &\PowMap{F}\restriction_{\powerset(A)\setminus C}
          \colon\powerset(A)\setminus C\bij\powerset(B)\setminus C
          \leftrightarrow{}\\[-2pt]
        &\quad
          \forall X\in\powerset(A)\,
            \bigl(X\in C\leftrightarrow F[X]\in C\bigr)
      \end{aligned}}{\rLRI{4,7}}
  \closeproofpartwide
\end{tabproofsplitwide}

Davon zu unterscheiden ist die Rückwirkung von Eigenschaften der
eingeschränkten Potenzmengenabbildung auf die Grundabbildung. Erhalten
gebliebene Einermengen liefern dafür hinreichende Bedingungen. Die folgenden
Beweise machen sichtbar, an welcher Stelle die jeweilige
Einermengenbedingung benötigt wird.

\FormulaThmDeltaK[Eingeschränkte Injektivität wird an Einermengen
reflektiert]{%
  \begin{aligned}[t]
    &F\colon A\to B\dsep
      \forall a\in A\,\bigl(\{a\}\notin C\bigr)\dsep{}\\[-2pt]
    &\PowMap{F}\restriction_{\powerset(A)\setminus C}
      \colon\powerset(A)\setminus C\inj\powerset(B)\setminus C\dsep{}\\[-2pt]
    &x,y\in A\dsep F(x)=F(y)\vdash x=y
  \end{aligned}%
}{RestrictedPowerMapInjectivityPointwiseReflection}{%
  \DeltaRow{Mengen}{A\dsep B\dsep C\dsep x\dsep y}
  \DeltaRow{Funktionen}{F\dsep\PowMap{F}}
}
\begin{tabproofwide}
  \proofstepwidestar[1]{F\colon A\to B}{\rA}
  \proofstepwidestar[2]{%
    \forall a\in A\,\bigl(\{a\}\notin C\bigr)}{\rA}
  \proofstepwidestar[3]{%
    \begin{aligned}[t]
      &\PowMap{F}\restriction_{\powerset(A)\setminus C}\colon{}\\[-2pt]
      &\quad\powerset(A)\setminus C\inj\powerset(B)\setminus C
    \end{aligned}}{\rA}
  \proofstepwidestar[4]{x\in A}{\rA}
  \proofstepwidestar[5]{y\in A}{\rA}
  \proofstepwidestar[6]{F(x)=F(y)}{\rA}
  \proofstepwidestar[4]{\{ x \} \subseteq A}{%
      \FormulaRefAuto{a\in A\vdash\{a\}\subseteq A}{4}}
  \proofstepwidestar[4]{\{x\}\in\powerset(A)}{%
    \FormulaRefAuto{x\in\powerset(A)\eqvdash x\subseteq A}{%
      7}}
  \proofstepwidestar[5]{\{ y \} \subseteq A}{%
      \FormulaRefAuto{a\in A\vdash\{a\}\subseteq A}{5}}
  \proofstepwidestar[5]{\{y\}\in\powerset(A)}{%
    \FormulaRefAuto{x\in\powerset(A)\eqvdash x\subseteq A}{%
      9}}
  \proofstepwidestar[2,4]{\{x\}\notin C}{\rRE{\rUE{2},4}}
  \proofstepwidestar[2,5]{\{y\}\notin C}{\rRE{\rUE{2},5}}
  \proofstepwidestar[2,4]{\{x\}\in\powerset(A)\setminus C}{%
    \FormulaRefAuto{x\in A\setminus B\eqvdash x\in A\land x\notin B}{%
      \rAI{8,11}}}
  \proofstepwidestar[2,5]{\{y\}\in\powerset(A)\setminus C}{%
    \FormulaRefAuto{x\in A\setminus B\eqvdash x\in A\land x\notin B}{%
      \rAI{10,12}}}
  \proofstepwidestar[]{\powerset(A)\setminus C\subseteq\powerset(A)}{%
    \FormulaRefAuto{A\setminus B\subseteq A}}
  \proofstepwidestar[1,2,4]{%
    \begin{aligned}[t]
      &\bigl(\PowMap{F}\restriction_{\powerset(A)\setminus C}\bigr)(\{x\})\\[-2pt]
      &\quad=\PowMap{F}(\{x\})
    \end{aligned}}{%
    \FormulaRefAuto{C\subseteq A\dsep x\in C\vdash
      F\restriction_C(x)=F(x)}{15,13}}
  \proofstepwidestar[1,2,5]{%
    \begin{aligned}[t]
      &\bigl(\PowMap{F}\restriction_{\powerset(A)\setminus C}\bigr)(\{y\})\\[-2pt]
      &\quad=\PowMap{F}(\{y\})
    \end{aligned}}{%
    \FormulaRefAuto{C\subseteq A\dsep x\in C\vdash
      F\restriction_C(x)=F(x)}{15,14}}
  \proofstepwidestar[1,4]{F [ \{ x \} ] = \{ F ( x ) \}}{%
      \FormulaRefAuto{x\in A\dsep F\colon A\to B
        \vdash F[\{x\}]=\{F(x)\}}{4,1}}
  \proofstepwidestar[1,4]{\PowMap F ( \{ x \} ) = F [ \{ x \} ]}{%
      \FormulaRefAuto{F\colon A\to B\dsep X\in\powerset(A)
        \vdash\PowMap{F}(X)=F[X]}{1,8}}
  \proofstepwidestar[1,4]{\PowMap{F}(\{x\})=\{F(x)\}}{%
    \FormulaRefAuto{a=b,\,b=c\vdash a=c}{%
      19,
      18}}
  \proofstepwidestar[1,5]{F [ \{ y \} ] = \{ F ( y ) \}}{%
      \FormulaRefAuto{x\in A\dsep F\colon A\to B
        \vdash F[\{x\}]=\{F(x)\}}{5,1}}
  \proofstepwidestar[1,5]{\PowMap F ( \{ y \} ) = F [ \{ y \} ]}{%
      \FormulaRefAuto{F\colon A\to B\dsep X\in\powerset(A)
        \vdash\PowMap{F}(X)=F[X]}{1,10}}
  \proofstepwidestar[1,5]{\PowMap{F}(\{y\})=\{F(y)\}}{%
    \FormulaRefAuto{a=b,\,b=c\vdash a=c}{%
      22,
      21}}
  \proofstepwidestar[6]{\{F(x)\}=\{F(y)\}}{%
    \FormulaRefAuto{\{a\}=\{b\}\eqvdash a=b}{6}}
  \proofstepwidestar[1,2,4,5,6]{%
    \begin{aligned}[t]
      &\bigl(\PowMap{F}\restriction_{\powerset(A)\setminus C}\bigr)(\{x\})\\[-2pt]
      &\quad=\{F(x)\}
    \end{aligned}}{%
    \FormulaRefAuto{a=b\dsep b=c\vdash a=c}{16,20}}
  \proofstepwidestar[1,2,4,5,6]{%
    \begin{aligned}[t]
      &\bigl(\PowMap{F}\restriction_{\powerset(A)\setminus C}\bigr)(\{x\})\\[-2pt]
      &\quad=\{F(y)\}
    \end{aligned}}{%
    \FormulaRefAuto{a=b\dsep b=c\vdash a=c}{25,24}}
  \proofstepwidestar[1,5]{\{ F ( y ) \} = \PowMap F ( \{ y \} )}{%
      \FormulaRefAuto{a=b\vdash b=a}{23}}
  \proofstepwidestar[1,2,4,5,6]{%
    \begin{aligned}[t]
      &\bigl(\PowMap{F}\restriction_{\powerset(A)\setminus C}\bigr)(\{x\})\\[-2pt]
      &\quad=\PowMap{F}(\{y\})
    \end{aligned}}{%
    \FormulaRefAuto{a=b\dsep b=c\vdash a=c}{%
      26,27}}
  \proofstepwidestar[1,2,5]{\PowMap F ( \{ y \} ) = ( \PowMap F \restriction _ { \powerset ( A ) \setminus C } ) ( \{ y \} )}{%
      \FormulaRefAuto{a=b\vdash b=a}{17}}
  \proofstepwidestar[1,2,4,5,6]{%
    \begin{aligned}[t]
      &\bigl(\PowMap{F}\restriction_{\powerset(A)\setminus C}\bigr)(\{x\})\\[-2pt]
      &\quad=\bigl(\PowMap{F}\restriction_{\powerset(A)\setminus C}\bigr)(\{y\})
    \end{aligned}}{%
    \FormulaRefAuto{a=b\dsep b=c\vdash a=c}{%
      28,29}}
  \proofstepwidestar[1,2,3,4,5,6]{\{x\}=\{y\}}{%
    \FormulaRefAuto{F\colon A\inj B\dsep x\in A\dsep y\in A\dsep
      F(x)=F(y)\vdash x=y}{3,13,14,30}}
  \proofstepwidestar[1,2,3,4,5,6]{x=y}{%
    \FormulaRefAuto{\{a\}=\{b\}\eqvdash a=b}{31}}
\end{tabproofwide}

\FormulaThmDeltaK[Eingeschränkte Injektivität erzwingt
Injektivität der Grundabbildung]{%
  \begin{aligned}[t]
    &F\colon A\to B\dsep
      \forall a\in A\,\bigl(\{a\}\notin C\bigr)\dsep{}\\[-2pt]
    &\PowMap{F}\restriction_{\powerset(A)\setminus C}
      \colon\powerset(A)\setminus C\inj\powerset(B)\setminus C
      \vdash F\colon A\inj B
  \end{aligned}%
}{RestrictedPowerMapInjectiveReflectsInjectivity}{%
  \DeltaRow{Mengen}{A\dsep B\dsep C\dsep x\dsep y}
  \DeltaRow{Funktionen}{F\dsep\PowMap{F}}
}
\begin{tabproofwide}
  \proofstepwidestar[1]{F\colon A\to B}{\rA}
  \proofstepwidestar[2]{\forall a\in A\,(\{a\}\notin C)}{\rA}
  \proofstepwidestar[3]{%
    \PowMap{F}\restriction_{\powerset(A)\setminus C}
      \colon\powerset(A)\setminus C\inj\powerset(B)\setminus C}{\rA}
  \proofstepwidestar[1,2,3]{%
    \forall x,y\in A\,\bigl(F(x)=F(y)\rightarrow x=y\bigr)}{%
    \FormulaRefAuto{RestrictedPowerMapInjectivityPointwiseReflection}[thm]{%
      1,2,3}}
  \proofstepwidestar[1,2,3]{F\colon A\inj B}{%
    \FormulaRefAuto{Injektive Funktion}[def]{1,4}}
\end{tabproofwide}

\FormulaThmDeltaK[Eingeschränkte Surjektivität wird an Einermengen
reflektiert]{%
  \begin{aligned}[t]
    &F\colon A\to B\dsep
      \forall b\in B\,\bigl(\{b\}\notin C\bigr)\dsep{}\\[-2pt]
    &\PowMap{F}\restriction_{\powerset(A)\setminus C}
      \colon\powerset(A)\setminus C\sur\powerset(B)\setminus C\dsep{}\\[-2pt]
    &y\in B\vdash\exists x\in A\,F(x)=y
  \end{aligned}%
}{RestrictedPowerMapSurjectivityPointwiseReflection}{%
  \DeltaRow{Mengen}{A\dsep B\dsep C\dsep X\dsep x\dsep y}
  \DeltaRow{Funktionen}{F\dsep\PowMap{F}}
}
\begin{tabproofwide}
  \proofstepwidestar[1]{F\colon A\to B}{\rA}
  \proofstepwidestar[2]{\forall b\in B\,(\{b\}\notin C)}{\rA}
  \proofstepwidestar[3]{%
    \PowMap{F}\restriction_{\powerset(A)\setminus C}
      \colon\powerset(A)\setminus C\sur\powerset(B)\setminus C}{\rA}
  \proofstepwidestar[4]{y\in B}{\rA}
  \proofstepwidestar[4]{\{ y \} \subseteq B}{%
      \FormulaRefAuto{a\in A\vdash\{a\}\subseteq A}{4}}
  \proofstepwidestar[4]{\{y\}\in\powerset(B)}{%
    \FormulaRefAuto{x\in\powerset(A)\eqvdash x\subseteq A}{%
      5}}
  \proofstepwidestar[2,4]{\{y\}\notin C}{\rRE{\rUE{2},4}}
  \proofstepwidestar[2,4]{\{y\}\in\powerset(B)\setminus C}{%
    \FormulaRefAuto{x\in A\setminus B\eqvdash x\in A\land x\notin B}{%
      \rAI{6,7}}}
  \proofstepwidestar[2,3,4]{%
    \exists X\in\powerset(A)\setminus C\,
      \bigl(\PowMap{F}\restriction_{\powerset(A)\setminus C}\bigr)(X)
        =\{y\}}{%
    \FormulaRefAuto{y\in B\vdash\exists x\in A\;F(x)=y}{3,8}}
  \proofstepwidestar[10]{%
    X\in\powerset(A)\setminus C\land
    \bigl(\PowMap{F}\restriction_{\powerset(A)\setminus C}\bigr)(X)
      =\{y\}}{\rA}
  \proofstepwidestar[10]{X\in\powerset(A)\setminus C}{\rAEa{10}}
  \proofstepwidestar[10]{%
    \bigl(\PowMap{F}\restriction_{\powerset(A)\setminus C}\bigr)(X)
      =\{y\}}{\rAEb{10}}
  \proofstepwidestar[]{\powerset(A)\setminus C\subseteq\powerset(A)}{%
    \FormulaRefAuto{A\setminus B\subseteq A}}
  \proofstepwidestar[10]{X\in\powerset(A)}{%
    \FormulaRefAuto{A\subseteq B\dsep x\in A\vdash x\in B}{13,11}}
  \proofstepwide[10]{\{y\}}{=}{%
    \bigl(\PowMap{F}\restriction_{\powerset(A)\setminus C}\bigr)(X)}{%
    \FormulaRefAuto{a=b\vdash b=a}{12}}
  \proofstepwide[10]{}{=}{\PowMap{F}(X)}{%
    \FormulaRefAuto{C\subseteq A\dsep x\in C\vdash
      F\restriction_C(x)=F(x)}{13,11}}
  \proofstepwide[1,10]{}{=}{F[X]}{%
    \FormulaRefAuto{F\colon A\to B\dsep X\in\powerset(A)
      \vdash\PowMap{F}(X)=F[X]}{1,14}}
  \proofstepwide[1,10]{\{y\}}{=}{F[X]}{\rChain{15,17}}
  \proofstepwidestar[]{y\in\{y\}}{%
    \FormulaRefAuto{x\in\{a\}\eqvdash x=a}{\rII}}
  \proofstepwidestar[1,10]{y\in F[X]}{\rIE{18,19}}
  \proofstepwidestar[10]{X\subseteq A}{%
    \FormulaRefAuto{x\in\powerset(A)\eqvdash x\subseteq A}{14}}
  \proofstepwidestar[1,10]{\exists x\in X\,y=F(x)}{%
    \FormulaRefAuto{C\subseteq A\dsep F\colon A\to B\dsep
      y\in F[C]\vdash\exists x\in C\,y=F(x)}{21,1,20}}
  \proofstepwidestar[23]{x\in X\land y=F(x)}{\rA}
  \proofstepwidestar[10,23]{x\in A}{%
    \FormulaRefAuto{A\subseteq B\dsep x\in A\vdash x\in B}{21,\rAEa{23}}}
  \proofstepwidestar[23]{F(x)=y}{%
    \FormulaRefAuto{a=b\vdash b=a}{\rAEb{23}}}
  \proofstepwidestar[10,23]{\exists x\in A\,F(x)=y}{%
    \rEI{\rAI{24,25}}}
  \proofstepwidestar[1,10]{\exists x\in A\,F(x)=y}{%
    \rEE{22,23,26}}
  \proofstepwidestar[1,2,3,4]{\exists x\in A\,F(x)=y}{%
    \rEE{9,10,27}}
\end{tabproofwide}

\FormulaThmDeltaK[Eingeschränkte Surjektivität erzwingt
Surjektivität der Grundabbildung]{%
  \begin{aligned}[t]
    &F\colon A\to B\dsep
      \forall b\in B\,\bigl(\{b\}\notin C\bigr)\dsep{}\\[-2pt]
    &\PowMap{F}\restriction_{\powerset(A)\setminus C}
      \colon\powerset(A)\setminus C\sur\powerset(B)\setminus C
      \vdash F\colon A\sur B
  \end{aligned}%
}{RestrictedPowerMapSurjectiveReflectsSurjectivity}{%
  \DeltaRow{Mengen}{A\dsep B\dsep C\dsep x\dsep y}
  \DeltaRow{Funktionen}{F\dsep\PowMap{F}}
}
\begin{tabproofwide}
  \proofstepwidestar[1]{F\colon A\to B}{\rA}
  \proofstepwidestar[2]{\forall b\in B\,(\{b\}\notin C)}{\rA}
  \proofstepwidestar[3]{%
    \PowMap{F}\restriction_{\powerset(A)\setminus C}
      \colon\powerset(A)\setminus C\sur\powerset(B)\setminus C}{\rA}
  \proofstepwidestar[1,2,3]{\forall y\in B\,\exists x\in A\,F(x)=y}{%
    \FormulaRefAuto{RestrictedPowerMapSurjectivityPointwiseReflection}[thm]{%
      1,2,3}}
  \proofstepwidestar[1,2,3]{F\colon A\sur B}{%
    \FormulaRefAuto{Surjektive Funktion}[def]{1,4}}
\end{tabproofwide}

\FormulaThmDeltaK[Bijektivität außerhalb einer Teilmengenfamilie
erzwingt Bijektivität der Grundabbildung]{%
  \begin{aligned}[t]
    &F\colon A\to B\dsep
      \forall a\in A\,\bigl(\{a\}\notin C\bigr)\dsep
      \forall b\in B\,\bigl(\{b\}\notin C\bigr)\dsep{}\\[-2pt]
    &\PowMap{F}\restriction_{\powerset(A)\setminus C}
      \colon\powerset(A)\setminus C\bij\powerset(B)\setminus C
      \vdash F\colon A\bij B
  \end{aligned}%
}{RestrictedPowerMapBijectiveReflectsBijectivity}{%
  \DeltaRow{Mengen}{A\dsep B\dsep C}
  \DeltaRow{Funktionen}{F\dsep\PowMap{F}}
}
\begin{tabproofsplitwide}
\proofpartwide{Injektivität}
\proofstepwidestar[1]{F\colon A\to B}{\rA}
  \proofstepwidestar[2]{\forall a\in A\,(\{a\}\notin C)}{\rA}
  \proofstepwidestar[3]{\forall b\in B\,(\{b\}\notin C)}{\rA}
  \proofstepwidestar[4]{%
    \begin{aligned}[t]
      &\PowMap{F}\restriction_{\powerset(A)\setminus C}\colon{}\\[-2pt]
      &\quad\powerset(A)\setminus C\bij\powerset(B)\setminus C
    \end{aligned}}{\rA}
  \proofstepwidestar[4]{%
    \begin{aligned}[t]
      &\PowMap{F}\restriction_{\powerset(A)\setminus C}\colon{}\\[-2pt]
      &\quad\powerset(A)\setminus C\inj\powerset(B)\setminus C
    \end{aligned}}{%
    \FormulaRefAuto{F\colon A\bij B\vdash F\colon A\inj B}{4}}
  \proofstepwidestar[4]{%
    \begin{aligned}[t]
      &\PowMap{F}\restriction_{\powerset(A)\setminus C}\colon{}\\[-2pt]
      &\quad\powerset(A)\setminus C\sur\powerset(B)\setminus C
    \end{aligned}}{%
    \FormulaRefAuto{F\colon A\bij B\vdash F\colon A\sur B}{4}}
  \proofstepwidestar[1,2,4]{F\colon A\inj B}{%
    \FormulaRefAuto{RestrictedPowerMapInjectiveReflectsInjectivity}[thm]{%
      1,2,5}}

\closeproofpartwide

\proofpartwide{Surjektivität}
  \setcounter{proofstepnr}{7}% Fortlaufende Schritte dieses Beweises.
\proofstepwidestar[1,3,4]{F\colon A\sur B}{%
    \FormulaRefAuto{RestrictedPowerMapSurjectiveReflectsSurjectivity}[thm]{%
      1,3,6}}

\closeproofpartwide

\proofpartwide{Bijektivität}
  \setcounter{proofstepnr}{8}% Fortlaufende Schritte dieses Beweises.
\proofstepwidestar[1,2,3,4]{F\colon A\bij B}{%
    \FormulaRefAuto{F\colon A\inj B\dsep F\colon A\sur B
      \vdash F\colon A\bij B}{7,8}}

\closeproofpartwide
\end{tabproofsplitwide}

\begin{remark}[Wozu die beiden Einermengenbedingungen dienen]
Der vorangehende Satz verwendet zwei verschiedene Sichtbarkeitsbedingungen:
\[
  \forall a\in A\,\bigl(\{a\}\notin C\bigr),\qquad
  \forall b\in B\,\bigl(\{b\}\notin C\bigr).
\]
Die erste hält die Einermengen des Quellbereichs sichtbar und ermöglicht so
die Rückwirkung auf die Injektivität von \(F\).  Die zweite hält die
Einermengen des Zielbereichs sichtbar und ermöglicht die Rückwirkung auf die
Surjektivität.  Für eine einzelne Abbildung können diese Voraussetzungen
stärker als nötig sein; aus dem allgemeinen Satz lässt sich jedoch keine von
ihnen ersatzlos streichen.

\emph{Ohne die Bedingung im Zielbereich.}
Setze
\[
  A=\varnothing,\qquad B=\{\varnothing\},\qquad
  C=\{B\}=\bigl\{\{\varnothing\}\bigr\}.
\]
Für das einzige Element \(b=\varnothing\) von \(B\) gilt dann
\(\{b\}=B\in C\): Gerade die Einermenge, mit der dieser Zielpunkt erkannt
würde, ist also entfernt.  Zugleich ist
\[
  \powerset(A)\setminus C
  =\{\varnothing\}
  =\powerset(B)\setminus C.
\]
Die eingeschränkte Potenzmengenabbildung besteht daher nur aus
\(\varnothing\mapsto\varnothing\) und ist bijektiv.  Die zugrunde liegende
Abbildung \(F\colon A\to B\) ist dagegen nicht surjektiv.

\emph{Ohne die Bedingung im Quellbereich.}
Setze
\[
  A=\{0,1\},\qquad B=\{2\},\qquad
  C=\bigl\{\{0\},\{1\}\bigr\}
\]
und sei \(F\colon A\to B\) konstant.  Nun liegen \(\{0\}\) und \(\{1\}\) in
\(C\); die beiden verschiedenen Quellpunkte können deshalb nicht mehr durch
ihre Einermengen unterschieden werden.  Tatsächlich bildet die
eingeschränkte Potenzmengenabbildung
\[
  \{\varnothing,\{0,1\}\}
    \quad\text{bijektiv auf}\quad
  \{\varnothing,\{2\}\}
\]
ab, obwohl \(F\) nicht injektiv ist.  Die beiden Beispiele zeigen jeweils
genau den Informationsverlust, den die betreffende Einermengenbedingung
ausschließt.

\emph{Warum die Bijektion von \(F\) induziert sein muss.}
Selbst wenn beide Einermengenbedingungen erfüllt sind, genügt die bloße
Existenz irgendeiner Bijektion zwischen den Restfamilien nicht.  Seien dazu
\[
  A=\{a_0,a_1\},\qquad B=\{b_0,b_1,b_2\},
\]
wobei \(A\cap B=\varnothing\), und sei
\[
  C=\{\varnothing\}\cup
    \Bigl(\powerset(B)\setminus
      \bigl\{\{b_0\},\{b_1\},\{b_2\}\bigr\}\Bigr).
\]
Dann gehört keine Einermenge eines Elements von \(A\) oder \(B\) zu \(C\),
und es gilt
\[
  \powerset(A)\setminus C
    =\bigl\{\{a_0\},\{a_1\},A\bigr\},\qquad
  \powerset(B)\setminus C
    =\bigl\{\{b_0\},\{b_1\},\{b_2\}\bigr\}.
\]
Zwischen diesen beiden dreielementigen Restfamilien existiert eine Bijektion;
zwischen der zweielementigen Menge \(A\) und der dreielementigen Menge \(B\)
existiert dagegen keine.  Dies widerspricht dem Satz nicht: Die beliebige
Bijektion der Restfamilien muss nicht die Einschränkung von \(\PowMap{F}\)
sein.  Gerade diese Bindung an die Grundabbildung \(F\) ist daher wesentlich.
\end{remark}

\section{Zusammengesetzte Bijektionen auf Mengenfamilien}

\ExplSyntaxOn
\tl_new:N \l_mx_band_eight_families_break_args_tl
\NewDocumentCommand{\BandEightFamiliesBreakArgs}{m}
  {
    \group_begin:
      \tl_set:Nn \l_mx_band_eight_families_break_args_tl {#1}
      \tl_replace_all:Nnn
        \l_mx_band_eight_families_break_args_tl {,} {,\allowbreak}
      \tl_use:N \l_mx_band_eight_families_break_args_tl
    \group_end:
  }
\ExplSyntaxOff
Die folgenden drei Konstruktionen fassen ein wiederkehrendes Verfahren
zusammen.  Zunächst werden Bijektionen auf disjunkten Teilen verklebt.  Dann
wird eine veränderliche Restmenge bei festgehaltenem Kern transportiert.
Schließlich wird eine solche Abbildung auf eine disjunkte, unveränderte
Schicht fortgesetzt.  Die Indexmengen und die beteiligten Elemente sind dabei
vollkommen beliebig.

\subsection{Verkleben disjunkter Zweige}

\FormulaThmDeltaK[Bijektives Verkleben disjunkter Zweige]{%
  \begin{aligned}[t]
    &A_0\cap A_1=\varnothing\dsep B_0\cap B_1=\varnothing\dsep
      f_0\colon A_0\bij B_0\dsep f_1\colon A_1\bij B_1\vdash{}\\[-2pt]
    &\CaseFun{A_0}{A_1}{f_0}{f_1}
      \colon A_0\cup A_1\bij B_0\cup B_1
  \end{aligned}%
}{CaseFunBijectiveDisjointTargets}{%
  \DeltaRow{Mengen}{A_0\dsep A_1\dsep B_0\dsep B_1\dsep x\dsep y}
  \DeltaRow{Hilfsmengen und -variablen}{A\dsep B\dsep C\dsep D\dsep u}
  \DeltaRow{Funktionen}{f_0\dsep f_1\dsep
    \CaseFun{A_0}{A_1}{f_0}{f_1}}
  \DeltaRow{Hilfsfunktionen}{f\dsep G}
}
\begin{tabproofsplitwide}
  \proofpartwideindR[Urbildtransport entlang eines surjektiven Zweiges]{%
    \begin{aligned}[t]
      &f\colon A\sur B\dsep A\subseteq C\dsep G\colon C\to D\dsep
        [u\in A]\vdots G(u)=f(u)\dsep y\in B\vdash{}\\[-2pt]
      &\exists x\in C\,G(x)=y
    \end{aligned}%
  }{%
    f\colon A\sur B\dsep A\subseteq C\dsep G\colon C\to D\dsep
    [u\in A]\vdots G(u)=f(u)\dsep y\in B
    \vdash \exists x\in C\,G(x)=y
  }[SurjectiveBranchWitnessTransfer]
    \proofstepwidestar[1]{f\colon A\sur B}{\rA}
    \proofstepwidestar[2]{A\subseteq C}{\rA}
    \proofstepwidestar[3]{G\colon C\to D}{\rA}
    \proofstepwidestar[4]{y\in B}{\rA}
    \proofstepwidestar[1,4]{\exists x\in A\,(f(x)=y)}{%
      \FormulaRefAuto{Surjektivität}[ax]{1,4}}
    \proofstepwidestar[6]{x\in A\land f(x)=y}{\rA}
    \proofstepwidestar[6]{x\in A}{\rAEa{6}}
    \proofstepwidestar[6]{f(x)=y}{\rAEb{6}}
    \proofstepwidestar[3,6]{G(x)=f(x)}{%
      \ensuremath{[u\in A]\vdots G(u)=f(u)(7)}}
    \proofstepwidestar[3,6]{G(x)=y}{%
      \FormulaRefAuto{a=b \dsep c=d \dsep P(b,d)\vdash P(a,c)}[thm]{%
        9,\rII,8}}
    \proofstepwidestar[2,6]{x\in C}{%
      \FormulaRefAuto{A\subseteq B\dsep x\in A\vdash x\in B}[thm]{2,7}}
    \proofstepwidestar[2,3,6]{\exists x\in C\,(G(x)=y)}{%
      \rEI{\rAI{11,10}}}
    \proofstepwidestar[1,2,3,4]{\exists x\in C\,(G(x)=y)}{%
      \rEE{5,6,12}}
  \closeproofpartwideindR

  \proofpartwideindR[Gemeinsame Zieltypisierung der Zweige]{%
    \begin{aligned}[t]
      &f_0\colon A_0\bij B_0\dsep
        f_1\colon A_1\bij B_1\vdash{}\\[-2pt]
      &\bigl(f_0\colon A_0\to B_0\cup B_1\bigr)\land{}\\[-2pt]
      &\quad\bigl(f_1\colon A_1\to B_0\cup B_1\bigr)
    \end{aligned}%
  }{%
    \begin{aligned}[t]
      &f_0\colon A_0\bij B_0\dsep
        f_1\colon A_1\bij B_1\vdash{}\\[-2pt]
      &\bigl(f_0\colon A_0\to B_0\cup B_1\bigr)\land
        \bigl(f_1\colon A_1\to B_0\cup B_1\bigr)
    \end{aligned}%
  }[CaseFunBijectiveBranchesCommonCodomain]
    \proofstepwidestar[1]{f_0\colon A_0\bij B_0}{\rA}
    \proofstepwidestar[2]{f_1\colon A_1\bij B_1}{\rA}
    \proofstepwidestar[1]{f_0\colon A_0\to B_0}{%
      \FormulaRefAuto{Bijektive Funktion}[def]{1}}
    \proofstepwidestar[2]{f_1\colon A_1\to B_1}{%
      \FormulaRefAuto{Bijektive Funktion}[def]{2}}
    \proofstepwidestar[]{B_0\subseteq B_0\cup B_1}{%
      \FormulaRefAuto{A\subseteq A\cup B}[thm]}
    \proofstepwidestar[]{B_1\subseteq B_0\cup B_1}{%
      \FormulaRefAuto{B\subseteq A\cup B}[thm]}
    \proofstepwidestar[1]{f_0\colon A_0\to B_0\cup B_1}{%
      \FormulaRefAuto{F\colon A\to B\dsep B\subseteq C
        \vdash F\colon A\to C}[thm]{3,5}}
    \proofstepwidestar[2]{f_1\colon A_1\to B_0\cup B_1}{%
      \FormulaRefAuto{F\colon A\to B\dsep B\subseteq C
        \vdash F\colon A\to C}[thm]{4,6}}
    \proofstepwidestar[1,2]{%
      \bigl(f_0\colon A_0\to B_0\cup B_1\bigr)\land
      \bigl(f_1\colon A_1\to B_0\cup B_1\bigr)}{\rAI{7,8}}
  \closeproofpartwideindR

  \proofpartwideindR[Surjektives Verkleben]{%
    \begin{aligned}[t]
      &A_0\cap A_1=\varnothing\dsep
        f_0\colon A_0\bij B_0\dsep f_1\colon A_1\bij B_1\vdash{}\\[-2pt]
      &\CaseFun{A_0}{A_1}{f_0}{f_1}
        \colon A_0\cup A_1\sur B_0\cup B_1
    \end{aligned}%
  }{%
    A_0\cap A_1=\varnothing\dsep
    f_0\colon A_0\bij B_0\dsep f_1\colon A_1\bij B_1
    \vdash \CaseFun{A_0}{A_1}{f_0}{f_1}
    \colon A_0\cup A_1\sur B_0\cup B_1
  }[CaseFunBijectiveDisjointTargetsSurjectivity]
    \proofstepwidestar[1]{A_0\cap A_1=\varnothing}{\rA}
    \proofstepwidestar[2]{f_0\colon A_0\bij B_0}{\rA}
    \proofstepwidestar[3]{f_1\colon A_1\bij B_1}{\rA}
    \proofstepwidestar[2,3]{f_0\colon A_0\to B_0\cup B_1\land f_1\colon A_1\to B_0\cup B_1}{%
      \FormulaRefAuto{CaseFunBijectiveBranchesCommonCodomain}[thm]{%
        2,3}}
    \proofstepwidestar[2,3]{f_0\colon A_0\to B_0\cup B_1}{%
      \rAEa{4}}
    \proofstepwidestar[2,3]{f_0\colon A_0\to B_0\cup B_1\land f_1\colon A_1\to B_0\cup B_1}{%
      \FormulaRefAuto{CaseFunBijectiveBranchesCommonCodomain}[thm]{%
        2,3}}
    \proofstepwidestar[2,3]{f_1\colon A_1\to B_0\cup B_1}{%
      \rAEb{6}}
    \proofstepwidestar[2,3]{%
      \CaseFun{A_0}{A_1}{f_0}{f_1}
        \colon A_0\cup A_1\to B_0\cup B_1}{%
      \FormulaRefAuto{f\colon A\to C\dsep g\colon B\to C\vdash
        \CaseFun{A}{B}{f}{g}\colon A\cup B\to C}[thm]{5,7}}
    \proofstepwidestar[2]{f_0\colon A_0\sur B_0}{%
      \FormulaRefAuto{F\colon A\bij B\vdash F\colon A\sur B}[thm]{2}}
    \proofstepwidestar[3]{f_1\colon A_1\sur B_1}{%
      \FormulaRefAuto{F\colon A\bij B\vdash F\colon A\sur B}[thm]{3}}
    \proofstepwidestar[11]{x\in A_0}{\rA}
    \proofstepwidestarbreakkeep[2,3,11]{%
      \CaseFun{A_0}{A_1}{f_0}{f_1}(x)=f_0(x)}{%
      \FormulaRefAuto{f\colon A\to C\dsep g\colon B\to C\dsep x\in A
        \vdash \CaseFun{A}{B}{f}{g}(x)=f(x)}[thm]{5,7,11}}
    \proofstepwidestar[13]{x\in A_1}{\rA}
    \proofstepwidestar[1,2,3,13]{%
      \CaseFun{A_0}{A_1}{f_0}{f_1}(x)=f_1(x)}{%
      \FormulaRefAuto{A\cap B=\varnothing\dsep f\colon A\to C\dsep
        g\colon B\to C\dsep x\in B\vdash
        \CaseFun{A}{B}{f}{g}(x)=g(x)}[thm]{1,5,7,13}}
    \proofstepwidestar[15]{y\in B_0\cup B_1}{\rA}
    \proofstepwidestar[15]{y\in B_0\lor y\in B_1}{%
      \FormulaRefAuto{z \in A \cup B \eqvdash z \in A \lor z \in B}{15}}
    \proofstepwidestar[17]{y\in B_0}{\rA}
    \proofstepwidestar[]{A_0\subseteq A_0\cup A_1}{%
      \FormulaRefAuto{A\subseteq A\cup B}[thm]}
    \proofstepwidestar[2,3,17]{\exists x\in A_0\cup A_1\,
      \CaseFun{A_0}{A_1}{f_0}{f_1}(x)=y}{%
      \FormulaRefAuto{SurjectiveBranchWitnessTransfer}[thm]{%
        9,18,8,11,12,17}}
    \proofstepwidestar[20]{y\in B_1}{\rA}
    \proofstepwidestar[]{A_1\subseteq A_0\cup A_1}{%
      \FormulaRefAuto{B\subseteq A\cup B}[thm]}
    \proofstepwidestar[1,2,3,20]{\exists x\in A_0\cup A_1\,
      \CaseFun{A_0}{A_1}{f_0}{f_1}(x)=y}{%
      \FormulaRefAuto{SurjectiveBranchWitnessTransfer}[thm]{%
        10,21,8,13,14,20}}
    \proofstepwidestar[1,2,3,15]{\exists x\in A_0\cup A_1\,
      \CaseFun{A_0}{A_1}{f_0}{f_1}(x)=y}{%
      \rOE{16,17,19,20,22}}
    \proofstepwidestar[1,2,3]{%
      \begin{aligned}[t]
        &\forall y\in B_0\cup B_1\,
          \exists x\in A_0\cup A_1\,\\[-2pt]
        &\quad\bigl(\CaseFun{A_0}{A_1}{f_0}{f_1}(x)=y\bigr)
      \end{aligned}}{%
      \rUI{\rRI{15,23}}}
    \proofstepwidestar[1,2,3]{%
      \CaseFun{A_0}{A_1}{f_0}{f_1}
        \colon A_0\cup A_1\sur B_0\cup B_1}{%
      \FormulaRefAuto{Surjektive Funktion}[def]{8,24}}
  \closeproofpartwideindR

  \proofpartwide{Schluss}
    \proofstepwidestar[1]{A_0\cap A_1=\varnothing}{\rA}
    \proofstepwidestar[2]{B_0\cap B_1=\varnothing}{\rA}
    \proofstepwidestar[3]{f_0\colon A_0\bij B_0}{\rA}
    \proofstepwidestar[4]{f_1\colon A_1\bij B_1}{\rA}
    \proofstepwidestar[]{B_0\subseteq B_0\cup B_1}{%
      \FormulaRefAuto{A\subseteq A\cup B}[thm]}
    \proofstepwidestar[3]{f_0\colon A_0\inj B_0\cup B_1}{%
      \FormulaRefAuto{F\colon A\inj B\dsep B\subseteq C\vdash
        F\colon A\inj C}[thm]{3,5}}
    \proofstepwidestar[]{B_1\subseteq B_0\cup B_1}{%
      \FormulaRefAuto{B\subseteq A\cup B}[thm]}
    \proofstepwidestar[4]{f_1\colon A_1\inj B_0\cup B_1}{%
      \FormulaRefAuto{F\colon A\inj B\dsep B\subseteq C\vdash
        F\colon A\inj C}[thm]{4,7}}
    \proofstepwidestar[2,3,4]{f_0[A_0]\cap f_1[A_1]=\varnothing}{%
      \FormulaRefAuto{DisjointTargetImagesDisjoint}[thm]{3,4,2}}
    \proofstepwidestar[1,2,3,4]{%
      \CaseFun{A_0}{A_1}{f_0}{f_1}\colon
      A_0\cup A_1\inj B_0\cup B_1}{%
      \FormulaRefAuto{CaseFunInjectiveDisjointImages}[thm]{1,6,8,9}}
    \proofstepwidestar[1,2,3,4]{%
      \CaseFun{A_0}{A_1}{f_0}{f_1}\colon
      A_0\cup A_1\sur B_0\cup B_1}{%
      \FormulaRefAuto{CaseFunBijectiveDisjointTargetsSurjectivity}[thm]{%
        1,3,4}}
    \proofstepwidestar[1,2,3,4]{%
      \CaseFun{A_0}{A_1}{f_0}{f_1}\colon
      A_0\cup A_1\bij B_0\cup B_1}{%
      \FormulaRefAuto{F\colon A\inj B\dsep F\colon A\sur B
        \vdash F\colon A\bij B}[thm]{10,11}}
  \closeproofpartwide
\end{tabproofsplitwide}

\subsection{Boolesche Intervalle mit festem Kern und ihr Transport}
Die folgende Notation ist die auf den Tr\"ager \(P\cup A\) spezialisierte
Intervallnotation aus der Mengenlehre.  Sie hebt den festen unteren Rand
\(P\) und die frei ver\"anderlichen Koordinaten \(A\setminus P\) hervor.

\FormulaDefDeltaK[Mengenfamilie mit festem Kern]{%
  \CoreFamily{P}{A}\coloneqq
  \bigl\{H\in\powerset(P\cup A)\bigm|P\subseteq H\bigr\}%
}{CoreFamilyDef}{%
  \DeltaRow{Mengen}{P\dsep A\dsep H\dsep\CoreFamily{P}{A}}
}

\FormulaThmDeltaK[Kernfamilien sind boolesche Intervalle]{%
  \CoreFamily{P}{A}=\mathfrak{I}[P,P\cup A]%
}{CoreFamilyBooleanInterval}{%
  \DeltaRow{Mengen}{P\dsep A\dsep H\dsep
    \CoreFamily{P}{A}\dsep\mathfrak{I}[P,P\cup A]}
}
\begin{tabproofwide}
  \proofstepwidestar[]{\CoreFamily{P}{A}=%
    \bigl\{H\in\powerset(P\cup A)\bigm|P\subseteq H\bigr\}}{%
    \FormulaRefAuto{CoreFamilyDef}[def]}
  \proofstepwidestar[]{\mathfrak I[P,P\cup A]=\{H\in\powerset(P\cup A)\mid P\subseteq H\}}{%
      \FormulaRefAuto{BooleanSetIntervalDef}[def]}
  \proofstepwidestar[]{\bigl\{H\in\powerset(P\cup A)\bigm|P\subseteq H\bigr\}=\mathfrak{I}[P,P\cup A]}{%
    \FormulaRefAuto{a=b\vdash b=a}[thm]{%
      2}}
  \proofstepwidestar[]{\CoreFamily{P}{A}=\mathfrak{I}[P,P\cup A]}{%
    \rChain{1,3}}
\end{tabproofwide}

\FormulaThmDeltaK[Elementkriterium einer Kernfamilie]{%
  H\in\CoreFamily{P}{A}\eqvdash P\subseteq H\land H\subseteq P\cup A
}{CoreFamilyMembership}{%
  \DeltaRow{Mengen}{P\dsep A\dsep H\dsep\CoreFamily{P}{A}}
}
\begin{tabproofwide}
  \proofstepwidestar[]{H\in\CoreFamily{P}{A}\leftrightarrow H\in\CoreFamily{P}{A}}{%
      \FormulaRefAuto{P\leftrightarrow P}[thm]}
  \proofstepwidestar[]{\CoreFamily{P}{A}=\mathfrak I[P,P\cup A]}{%
      \FormulaRefAuto{CoreFamilyBooleanInterval}[thm]}
  \proofstepwide[]{H\in\CoreFamily{P}{A}}{\leftrightarrow}{%
    H\in\mathfrak{I}[P,P\cup A]}{%
    \rIE{2,
      1}}
  \proofstepwide[]{}{\leftrightarrow}{%
    P\subseteq H\land H\subseteq P\cup A}{%
    \FormulaRefAuto{BooleanSetIntervalMembership}[thm]}
  \proofstepwide[]{H\in\CoreFamily{P}{A}}{\leftrightarrow}{%
    P\subseteq H\land H\subseteq P\cup A}{\rChain{3,4}}
\end{tabproofwide}

\FormulaThmDeltaK[Jede Kernfamilie liegt in der Potenzmenge ihres Trägers]{%
  \CoreFamily{P}{A}\subseteq\powerset(P\cup A)%
}{CoreFamilySubsetPowerset}{%
  \DeltaRow{Mengen}{P\dsep A\dsep H\dsep\CoreFamily{P}{A}}
}
\begin{tabproofwide}
  \proofstepwidestar[1]{H\in\CoreFamily{P}{A}}{\rA}
  \proofstepwidestar[1]{P\subseteq H\land H\subseteq P\cup A}{%
    \FormulaRefAuto{CoreFamilyMembership}[thm]{1}}
  \proofstepwidestar[1]{H\in\powerset(P\cup A)}{%
    \FormulaRefAuto{x\in\powerset(A)\eqvdash x\subseteq A}[thm]{\rAEb{2}}}
  \proofstepwidestar[]{\CoreFamily{P}{A}\subseteq\powerset(P\cup A)}{%
    \FormulaRefAuto{A\subseteq B\coloneq
      \forall x\,(x\in A\rightarrow x\in B)}[def]{\rUI{\rRI{1,3}}}}
\end{tabproofwide}

\FormulaThmDeltaK[Zerlegung einer Kernfamilie an einem adjungierten Punkt]{%
  c\notin P\cup A\vdash
  \begin{aligned}[t]
    &\text{(i)}\;\CoreFamily{P}{A\cup\{c\}}
      =\CoreFamily{P}{A}\cup\CoreFamily{P\cup\{c\}}{A}\\[-2pt]
    &\text{(ii)}\;\CoreFamily{P}{A}\cap
      \CoreFamily{P\cup\{c\}}{A}=\varnothing
  \end{aligned}%
}{CoreFamilyAdjoinedSplit}{%
  \DeltaRow{Mengen}{P\dsep A\dsep c\dsep H\dsep
    \CoreFamily{P}{A\cup\{c\}}\dsep\CoreFamily{P}{A}\dsep
    \CoreFamily{P\cup\{c\}}{A}}
}
\begin{tabproofsplitwide}
  \proofpartwideindR[Zerlegung]{%
    c\notin P\cup A\vdash
    \CoreFamily{P}{A\cup\{c\}}
      =\CoreFamily{P}{A}\cup\CoreFamily{P\cup\{c\}}{A}%
  }{%
    c\notin P\cup A\vdash
    \CoreFamily{P}{A\cup\{c\}}
      =\CoreFamily{P}{A}\cup\CoreFamily{P\cup\{c\}}{A}%
  }[CoreFamilyAdjoinedSplitCover]
    \multicolumn{6}{@{}l}{\textbf{Hinrichtung: Fallzerlegung nach \(c\in H\)}}\\[.4ex]
    \proofstepwidestar[1]{c\notin P\cup A}{\rA}
    \proofstepwidestar[]{P\cup(A\cup\{c\})=(P\cup A)\cup\{c\}}{%
      \FormulaRefAuto{A\cup(B\cup C)=(A\cup B)\cup C}[thm]}
    \proofstepwidestar[]{A\cup\{c\}=\{c\}\cup A}{%
      \FormulaRefAuto{A\cup B=B\cup A}[thm]}
    \proofstepwidestar[]{P\cup(A\cup\{c\})=P\cup(\{c\}\cup A)}{%
      \rLRS{3}}
    \proofstepwidestar[]{P\cup(\{c\}\cup A)=(P\cup\{c\})\cup A}{%
      \FormulaRefAuto{A\cup(B\cup C)=(A\cup B)\cup C}[thm]}
    \proofstepwidestar[]{P\cup(A\cup\{c\})=(P\cup\{c\})\cup A}{%
      \rChain{4,5}}
    \proofstepwidestar[7]{H\in\CoreFamily{P}{A\cup\{c\}}}{\rA}
    \proofstepwidestar[7]{P\subseteq H\land H\subseteq P\cup(A\cup\{c\})}{%
      \FormulaRefAuto{CoreFamilyMembership}[thm]{7}}
    \proofstepwidestar[7]{c\in H\lor c\notin H}{%
      \FormulaRefAuto{P\lor\neg P}[thm]}
    \proofstepwidestar[10]{c\in H}{\rA}
    \proofstepwidestar[10]{\{c\}\subseteq H}{%
      \FormulaRefAuto{x\in A\vdash\{x\}\subseteq A}[thm]{10}}
    \proofstepwidestar[7,10]{P\cup\{c\}\subseteq H}{%
      \FormulaRefAuto{A\subseteq B\dsep C\subseteq B\vdash
        A\cup C\subseteq B}[thm]{\rAEa{8},11}}
    \proofstepwidestar[7]{H\subseteq(P\cup\{c\})\cup A}{%
      \rIE{6,\rAEb{8}}}
    \proofstepwidestar[7,10]{H\in\CoreFamily{P\cup\{c\}}{A}}{%
      \FormulaRefAuto{CoreFamilyMembership}[thm]{\rAI{12,13}}}
    \proofstepwidestar[7,10]{H\in
      \CoreFamily{P}{A}\cup\CoreFamily{P\cup\{c\}}{A}}{%
      \FormulaRefAuto{x\in A\cup B\eqvdash x\in A\lor x\in B}[thm]{%
        \rOIb{14}}}
    \proofstepwidestar[16]{c\notin H}{\rA}
    \proofstepwidestar[7]{H\subseteq(P\cup A)\cup\{c\}}{%
      \rIE{2,\rAEb{8}}}
    \proofstepwidestar[1,7,16]{H\subseteq P\cup A}{%
      \FormulaRefAuto{SubsetAdjoinWithoutPoint}[thm]{17,16}}
    \proofstepwidestar[1,7,16]{H\in\CoreFamily{P}{A}}{%
      \FormulaRefAuto{CoreFamilyMembership}[thm]{\rAI{\rAEa{8},18}}}
    \proofstepwidestar[1,7,16]{H\in
      \CoreFamily{P}{A}\cup\CoreFamily{P\cup\{c\}}{A}}{%
      \FormulaRefAuto{x\in A\cup B\eqvdash x\in A\lor x\in B}[thm]{%
        \rOIa{19}}}
    \proofstepwidestar[1,7]{H\in
      \CoreFamily{P}{A}\cup\CoreFamily{P\cup\{c\}}{A}}{%
      \rOE{9,10,15,16,20}}
    \proofstepwidestar[1]{\CoreFamily{P}{A\cup\{c\}}\subseteq
      \CoreFamily{P}{A}\cup\CoreFamily{P\cup\{c\}}{A}}{%
      \FormulaRefAuto{A\subseteq B\coloneq
        \forall x\,(x\in A\rightarrow x\in B)}[def]{\rUI{\rRI{7,21}}}}

    \addlinespace[.8ex]
    \multicolumn{6}{@{}l}{\textbf{Rückrichtung: Einbettung beider Teile}}\\[.4ex]
    \proofstepwidestar[23]{H\in
      \CoreFamily{P}{A}\cup\CoreFamily{P\cup\{c\}}{A}}{\rA}
    \proofstepwidestar[23]{H\in\CoreFamily{P}{A}\lor
      H\in\CoreFamily{P\cup\{c\}}{A}}{%
      \FormulaRefAuto{x\in A\cup B\eqvdash x\in A\lor x\in B}[thm]{23}}
    \proofstepwidestar[25]{H\in\CoreFamily{P}{A}}{\rA}
    \proofstepwidestar[25]{P\subseteq H\land H\subseteq P\cup A}{%
      \FormulaRefAuto{CoreFamilyMembership}[thm]{25}}
    \proofstepwidestar[]{P \subseteq P}{%
      \FormulaRefAuto{A\subseteq A}[thm]}
    \proofstepwidestar[]{A\subseteq A\cup\{c\}}{%
      \FormulaRefAuto{A\subseteq A\cup B}[thm]}
    \proofstepwidestar[]{P\cup A\subseteq P\cup(A\cup\{c\})}{%
      \FormulaRefAuto{A\subseteq B\dsep C\subseteq D\vdash
        A\cup C\subseteq B\cup D}[thm]{%
        27,
        28}}
    \proofstepwidestar[25]{H\subseteq P\cup(A\cup\{c\})}{%
      \FormulaRefAuto{A\subseteq B\dsep B\subseteq C\vdash A\subseteq C}[thm]{%
        \rAEb{26},29}}
    \proofstepwidestar[25]{H\in\CoreFamily{P}{A\cup\{c\}}}{%
      \FormulaRefAuto{CoreFamilyMembership}[thm]{\rAI{\rAEa{26},30}}}
    \proofstepwidestar[32]{H\in\CoreFamily{P\cup\{c\}}{A}}{\rA}
    \proofstepwidestar[32]{P\cup\{c\}\subseteq H\land
      H\subseteq(P\cup\{c\})\cup A}{%
      \FormulaRefAuto{CoreFamilyMembership}[thm]{32}}
    \proofstepwidestar[]{P\subseteq P\cup\{c\}}{%
      \FormulaRefAuto{A\subseteq A\cup B}[thm]}
    \proofstepwidestar[32]{P\subseteq H}{%
      \FormulaRefAuto{A\subseteq B\dsep B\subseteq C\vdash A\subseteq C}[thm]{%
        34,\rAEa{33}}}
    \proofstepwidestar[]{(P\cup\{c\})\cup A=P\cup(\{c\}\cup A)}{%
      \FormulaRefAuto{(A\cup B)\cup C=A\cup(B\cup C)}[thm]}
    \proofstepwidestar[]{\{c\}\cup A=A\cup\{c\}}{%
      \FormulaRefAuto{A\cup B=B\cup A}[thm]}
    \proofstepwidestar[]{P\cup(\{c\}\cup A)=P\cup(A\cup\{c\})}{%
      \rLRS{37}}
    \proofstepwidestar[]{(P\cup\{c\})\cup A=P\cup(A\cup\{c\})}{%
      \rChain{36,38}}
    \proofstepwidestar[32]{H\subseteq P\cup(A\cup\{c\})}{%
      \rIE{39,\rAEb{33}}}
    \proofstepwidestar[32]{H\in\CoreFamily{P}{A\cup\{c\}}}{%
      \FormulaRefAuto{CoreFamilyMembership}[thm]{\rAI{35,40}}}
    \proofstepwidestar[23]{H\in\CoreFamily{P}{A\cup\{c\}}}{%
      \rOE{24,25,31,32,41}}
    \proofstepwidestar[]{%
      \CoreFamily{P}{A}\cup\CoreFamily{P\cup\{c\}}{A}
      \subseteq\CoreFamily{P}{A\cup\{c\}}}{%
      \FormulaRefAuto{A\subseteq B\coloneq
        \forall x\,(x\in A\rightarrow x\in B)}[def]{\rUI{\rRI{23,42}}}}
    \proofstepwidestar[1]{\CoreFamily{P}{A\cup\{c\}}
      =\CoreFamily{P}{A}\cup\CoreFamily{P\cup\{c\}}{A}}{%
      \FormulaRefAuto{A\subseteq B\dsep B\subseteq A\vdash A=B}[thm]{22,43}}
  \closeproofpartwideindR

  \proofpartwideindR[Disjunktheit]{%
    c\notin P\cup A\vdash
    \CoreFamily{P}{A}\cap\CoreFamily{P\cup\{c\}}{A}=\varnothing
  }{%
    c\notin P\cup A\vdash
    \CoreFamily{P}{A}\cap\CoreFamily{P\cup\{c\}}{A}=\varnothing
  }[CoreFamilyAdjoinedSplitDisjoint]
    \proofstepwidestar[1]{c\notin P\cup A}{\rA}
    \proofstepwidestar[2]{H\in
      \CoreFamily{P}{A}\cap\CoreFamily{P\cup\{c\}}{A}}{\rA}
    \proofstepwidestar[2]{H\in\CoreFamily{P}{A}\land
      H\in\CoreFamily{P\cup\{c\}}{A}}{%
      \FormulaRefAuto{x\in A\cap B\eqvdash x\in A\land x\in B}[thm]{2}}
    \proofstepwidestar[2]{P\subseteq H\land H\subseteq P\cup A}{%
      \FormulaRefAuto{CoreFamilyMembership}[thm]{\rAEa{3}}}
    \proofstepwidestar[2]{H\subseteq P\cup A}{%
      \rAEb{4}}
    \proofstepwidestar[2]{P\cup\{c\}\subseteq H\land H\subseteq(P\cup\{c\})\cup A}{%
      \FormulaRefAuto{CoreFamilyMembership}[thm]{\rAEb{3}}}
    \proofstepwidestar[2]{P\cup\{c\}\subseteq H}{%
      \rAEa{6}}
    \proofstepwidestar[]{c\in\{c\}}{%
      \FormulaRefAuto{a\in\{a\}}[thm]}
    \proofstepwidestar[]{c\in P\cup\{c\}}{%
      \FormulaRefAuto{x\in A\cup B\eqvdash x\in A\lor x\in B}[thm]{%
        \rOIb{8}}}
    \proofstepwidestar[2]{c\in H}{%
      \FormulaRefAuto{A\subseteq B\dsep x\in A\vdash x\in B}[thm]{7,9}}
    \proofstepwidestar[2]{c\in P\cup A}{%
      \FormulaRefAuto{A\subseteq B\dsep x\in A\vdash x\in B}[thm]{5,10}}
    \proofstepwidestar[1,2]{\bot}{\rBI{1,11}}
    \proofstepwidestar[1]{%
      \CoreFamily{P}{A}\cap\CoreFamily{P\cup\{c\}}{A}=\varnothing}{%
      \FormulaRefAuto{\forall x\,x\notin A\vdash A=\varnothing}[thm]{%
        \rUI{\rCI{2,12}}}}
  \closeproofpartwideindR
\end{tabproofsplitwide}

\FormulaThmDeltaK[Typisierung des Kerntransports]{%
  \begin{aligned}[t]
    &P\cap A=\varnothing\dsep Q\cap B=\varnothing\dsep
      f\colon A\to B\dsep H\in\CoreFamily{P}{A}\vdash{}\\[-2pt]
    &Q\cup f[H\setminus P]\in\CoreFamily{Q}{B}
  \end{aligned}%
}{CoreTransportTarget}{%
  \DeltaRow{Mengen}{P\dsep Q\dsep A\dsep B\dsep H\dsep x}
  \DeltaRow{Funktionen}{f}
}
\begin{tabproofwide}
  \proofstepwidestar[1]{P\cap A=\varnothing}{\rA}
  \proofstepwidestar[2]{Q\cap B=\varnothing}{\rA}
  \proofstepwidestar[3]{f\colon A\to B}{\rA}
  \proofstepwidestar[4]{H\in\CoreFamily{P}{A}}{\rA}
  \proofstepwidestar[4]{P\subseteq H\land H\subseteq P\cup A}{%
    \FormulaRefAuto{CoreFamilyMembership}[thm]{4}}
  \proofstepwidestar[4]{P\subseteq H}{\rAEa{5}}
  \proofstepwidestar[4]{H\subseteq P\cup A}{\rAEb{5}}
  \proofstepwidestar[8]{x\in H\setminus P}{\rA}
  \proofstepwidestar[8]{x\in H\land x\notin P}{%
    \FormulaRefAuto{x\in A\setminus B\eqvdash x\in A\land x\notin B}[thm]{8}}
  \proofstepwidestar[4,8]{x\in P\cup A}{%
    \FormulaRefAuto{A\subseteq B\dsep x\in A\vdash x\in B}[thm]{7,\rAEa{9}}}
  \proofstepwidestar[4,8]{x\in P\lor x\in A}{%
    \FormulaRefAuto{z \in A \cup B \eqvdash z \in A \lor z \in B}[thm]{10}}
  \proofstepwidestar[4,8]{x\in A}{%
    \FormulaRefAuto{P\lor Q\dsep\neg P\vdash Q}[thm]{11,\rAEb{9}}}
  \proofstepwidestar[4]{H\setminus P\subseteq A}{%
    \FormulaRefAuto{A\subseteq B \coloneq
      \forall x\,(x\in A\rightarrow x\in B)}[def]{\rUI{\rRI{8,12}}}}
  \proofstepwidestar[3,4]{f[H\setminus P]\subseteq B}{%
    \FormulaRefAuto{C\subseteq A\dsep F\colon A\to B
      \vdash F[C]\subseteq B}[thm]{13,3}}
  \proofstepwidestar[]{Q \subseteq Q}{%
      \FormulaRefAuto{A\subseteq A}[thm]}
  \proofstepwidestar[3,4]{Q\cup f[H\setminus P]\subseteq Q\cup B}{%
    \FormulaRefAuto{A\subseteq B\dsep C\subseteq D\vdash
      A\cup C\subseteq B\cup D}[thm]{%
      15,14}}
  \proofstepwidestar[]{Q\subseteq Q\cup f[H\setminus P]}{%
    \FormulaRefAuto{A\subseteq A\cup B}[thm]}
  \proofstepwidestarbreakkeep[3,4]{Q\cup f[H\setminus P]\in\CoreFamily{Q}{B}}{%
    \FormulaRefAuto{CoreFamilyMembership}[thm]{\rAI{17,16}}}
\end{tabproofwide}

\newpage
\FormulaDefDeltaK[Kerntransport]{%
  P\cap A=\varnothing\dsep Q\cap B=\varnothing\dsep f\colon A\to B
  \vdash
  \begin{aligned}[t]
    &\CoreTransport{P}{Q}{f}\colon
      \CoreFamily{P}{A}\to\CoreFamily{Q}{B},\\[-2pt]
    &\forall H\in\CoreFamily{P}{A}\,
      \Bigl(\CoreTransport{P}{Q}{f}(H)
        \coloneqq Q\cup f[H\setminus P]\Bigr)
  \end{aligned}%
}{CoreTransportDef}{%
  \DeltaRow{Mengen}{P\dsep Q\dsep A\dsep B\dsep H}
  \DeltaRow{Funktionen}{f\dsep\CoreTransport{P}{Q}{f}}
}
\begin{tabproof}
  \proofstep{1}{P\cap A=\varnothing}{\rA}
  \proofstep{2}{Q\cap B=\varnothing}{\rA}
  \proofstep{3}{f\colon A\to B}{\rA}
  \proofstep{4}{H\in\CoreFamily{P}{A}}{\rA}
  \proofstep{1,2,3,4}{Q\cup f[H\setminus P]\in\CoreFamily{Q}{B}}{%
    \FormulaRefAuto{CoreTransportTarget}[thm]{1,2,3,4}}
\end{tabproof}

\FormulaThmDeltaK[Typ und Grundgleichung des Kerntransports]{%
  \begin{aligned}[t]
    &\text{(i)}\;P\cap A=\varnothing\dsep Q\cap B=\varnothing\dsep
      f\colon A\to B\vdash \CoreTransport{P}{Q}{f}\colon
      \CoreFamily{P}{A}\to\CoreFamily{Q}{B}\\[-2pt]
    &\text{(ii)}\;P\cap A=\varnothing\dsep Q\cap B=\varnothing\dsep
      f\colon A\to B\dsep H\in\CoreFamily{P}{A}\vdash
      \CoreTransport{P}{Q}{f}(H)=Q\cup f[H\setminus P]
  \end{aligned}%
}{CoreTransportValue}{%
  \DeltaRow{Mengen}{P\dsep Q\dsep A\dsep B\dsep H}
  \DeltaRow{Funktionen}{f\dsep\CoreTransport{P}{Q}{f}}
}
\begin{tabproofsplitwide}
  \proofpartwideindR[Typ]{%
    P\cap A=\varnothing\dsep Q\cap B=\varnothing\dsep f\colon A\to B
    \vdash \CoreTransport{P}{Q}{f}\colon
      \CoreFamily{P}{A}\to\CoreFamily{Q}{B}%
  }{%
    P\cap A=\varnothing\dsep Q\cap B=\varnothing\dsep f\colon A\to B
    \vdash \CoreTransport{P}{Q}{f}\colon
      \CoreFamily{P}{A}\to\CoreFamily{Q}{B}%
  }[CoreTransportMapping]
    \proofstepwidestar[1]{P\cap A=\varnothing}{\rA}
    \proofstepwidestar[2]{Q\cap B=\varnothing}{\rA}
    \proofstepwidestar[3]{f\colon A\to B}{\rA}
    \proofstepwidestar[1,2,3]{\CoreTransport{P}{Q}{f}\colon
      \CoreFamily{P}{A}\to\CoreFamily{Q}{B}}{%
      \FormulaRefAuto{CoreTransportDef}[def]{1,2,3}}
  \closeproofpartwideindR

  \proofpartwideindR[Grundgleichung]{%
    P\cap A=\varnothing\dsep Q\cap B=\varnothing\dsep f\colon A\to B
    \dsep H\in\CoreFamily{P}{A}\vdash
    \CoreTransport{P}{Q}{f}(H)=Q\cup f[H\setminus P]%
  }{%
    P\cap A=\varnothing\dsep Q\cap B=\varnothing\dsep f\colon A\to B
    \dsep H\in\CoreFamily{P}{A}\vdash
    \CoreTransport{P}{Q}{f}(H)=Q\cup f[H\setminus P]%
  }[CoreTransportEvaluation]
    \proofstepwidestar[1]{P\cap A=\varnothing}{\rA}
    \proofstepwidestar[2]{Q\cap B=\varnothing}{\rA}
    \proofstepwidestar[3]{f\colon A\to B}{\rA}
    \proofstepwidestar[4]{H\in\CoreFamily{P}{A}}{\rA}
    \proofstepwidestar[1,2,3]{%
      \forall K\in\CoreFamily{P}{A}\,
        \bigl(\CoreTransport{P}{Q}{f}(K)=Q\cup f[K\setminus P]\bigr)}{%
      \FormulaRefAuto{CoreTransportDef}[def]{1,2,3}}
    \proofstepwidestar[1,2,3,4]{\CoreTransport{P}{Q}{f}(H)=
      Q\cup f[H\setminus P]}{\rRE{\rUE{5},4}}
  \closeproofpartwideindR
\end{tabproofsplitwide}

\FormulaThmDeltaK[Rest einer Kernfamilie liegt in der variablen Schicht]{%
  P\cap A=\varnothing\dsep H\in\CoreFamily{P}{A}
  \vdash H\setminus P\subseteq A
}{CoreFamilyRemainder}{%
  \DeltaRow{Mengen}{P\dsep A\dsep H\dsep x}
}
\begin{tabproof}
  \proofstep{1}{P\cap A=\varnothing}{\rA}
  \proofstep{2}{H\in\CoreFamily{P}{A}}{\rA}
  \proofstep{2}{P\subseteq H\land H\subseteq P\cup A}{%
      \FormulaRefAuto{CoreFamilyMembership}[thm]{2}}
  \proofstep{2}{H\subseteq P\cup A}{%
    \rAEb{3}}
  \proofstep{5}{x\in H\setminus P}{\rA}
  \proofstep{5}{x\in H\land x\notin P}{%
    \FormulaRefAuto{x\in A\setminus B\eqvdash x\in A\land x\notin B}[thm]{5}}
  \proofstep{2,5}{x\in P\cup A}{%
    \FormulaRefAuto{A\subseteq B\dsep x\in A\vdash x\in B}[thm]{4,\rAEa{6}}}
  \proofstep{2,5}{x\in P\lor x\in A}{%
    \FormulaRefAuto{z \in A \cup B \eqvdash z \in A \lor z \in B}[thm]{7}}
  \proofstep{2,5}{x\in A}{%
    \FormulaRefAuto{P\lor Q\dsep\neg P\vdash Q}[thm]{8,\rAEb{6}}}
  \proofstep{1,2}{H\setminus P\subseteq A}{%
    \FormulaRefAuto{A\subseteq B \coloneq
      \forall x\,(x\in A\rightarrow x\in B)}[def]{\rUI{\rRI{5,9}}}}
\end{tabproof}

\FormulaThmDeltaK[Hin- und Rücktransport einer Kernfamilie]{%
  \begin{aligned}[t]
    &P\cap A=\varnothing\dsep Q\cap B=\varnothing\dsep
      f\colon A\to B\dsep g\colon B\to A\dsep
      \forall x\in A\,g(f(x))=x\dsep{}\\[-2pt]
    &H\in\CoreFamily{P}{A}\vdash
      \CoreTransport{Q}{P}{g}
      \bigl(\CoreTransport{P}{Q}{f}(H)\bigr)=H
  \end{aligned}%
}{CoreTransportRoundTrip}{%
  \DeltaRow{Mengen}{P\dsep Q\dsep A\dsep B\dsep H\dsep K\dsep x}
  \DeltaRow{Funktionen}{f\dsep g\dsep
    \CoreTransport{P}{Q}{f}\dsep\CoreTransport{Q}{P}{g}}
}
\begin{tabproofwide}
  \proofstepwidestar[1]{P\cap A=\varnothing}{\rA}
  \proofstepwidestar[2]{Q\cap B=\varnothing}{\rA}
  \proofstepwidestar[3]{f\colon A\to B}{\rA}
  \proofstepwidestar[4]{g\colon B\to A}{\rA}
  \proofstepwidestar[5]{\forall x\in A\,g(f(x))=x}{\rA}
  \proofstepwidestar[6]{H\in\CoreFamily{P}{A}}{\rA}
  \proofstepwidestar[1,6]{H\setminus P\subseteq A}{%
    \FormulaRefAuto{CoreFamilyRemainder}[thm]{1,6}}
  \proofstepwidestar[3,6]{f[H\setminus P]\subseteq B}{%
    \FormulaRefAuto{C\subseteq A\dsep F\colon A\to B
      \vdash F[C]\subseteq B}[thm]{7,3}}
  \proofstepwidestar[]{Q \subseteq Q}{%
      \FormulaRefAuto{A\subseteq A}[thm]}
  \proofstepwidestar[2,3,6]{Q\cap f[H\setminus P]=\varnothing}{%
    \FormulaRefAuto{SubsetsOfDisjointSets}[thm]{%
      9,8,2}}
  \proofstepwidestar[1,2,3,6]{%
    \CoreTransport{P}{Q}{f}(H)=Q\cup f[H\setminus P]}{%
    \FormulaRefAuto{CoreTransportEvaluation}[thm]{1,2,3,6}}
  \proofstepwidestar[1,2,3]{\CoreTransport{P}{Q}{f}\colon\CoreFamily{P}{A}\to\CoreFamily{Q}{B}}{%
      \FormulaRefAuto{CoreTransportMapping}[thm]{1,2,3}}
  \proofstepwidestar[1,2,3,6]{%
    \CoreTransport{P}{Q}{f}(H)\in\CoreFamily{Q}{B}}{%
    \FormulaRefAuto{F\colon A\to B\dsep x\in A\vdash F(x)\in B}[thm]{%
      12,6}}
  \proofstepwidestarbreakkeep[1,2,3,4,6]{%
    \CoreTransport{Q}{P}{g}(\CoreTransport{P}{Q}{f}(H))
    =P\cup g[\CoreTransport{P}{Q}{f}(H)\setminus Q]}{%
    \FormulaRefAuto{CoreTransportEvaluation}[thm]{2,1,4,13}}
  \proofstepwidestarbreakkeep[2,3,6]{%
    (Q\cup f[H\setminus P])\setminus Q=f[H\setminus P]}{%
    \FormulaRefAuto{RemoveDisjointUnionSummand}[thm]{10}}
  \proofstepwidestarbreakkeep[1,2,3,6]{%
    \CoreTransport{P}{Q}{f}(H)\setminus Q=f[H\setminus P]}{%
    \rIE{11,15}}
  \proofstepwidestarbreakkeep[3,4,5,6]{g[f[H\setminus P]]=H\setminus P}{%
    \FormulaRefAuto{LeftInverseImages}[thm]{3,4,5,7}}
  \proofstepwidestar[1,2,3,4,5,6]{%
    \CoreTransport{Q}{P}{g}(\CoreTransport{P}{Q}{f}(H))
    =P\cup(H\setminus P)}{\rChain{14,\rLRS{16},\rLRS{17}}}
  \proofstepwidestar[6]{P\subseteq H\land H\subseteq P\cup A}{%
      \FormulaRefAuto{CoreFamilyMembership}[thm]{6}}
  \proofstepwidestar[6]{H=P\cup(H\setminus P)}{%
    \FormulaRefAuto{UnionWithRelativeRemainder}[thm]{%
      \rAEa{19}}}
  \proofstepwidestar[1,2,3,4,5,6]{%
    \CoreTransport{Q}{P}{g}(\CoreTransport{P}{Q}{f}(H))=H}{%
    \rChain{18,\rLRS{20}}}
\end{tabproofwide}

\newpage
\FormulaThmDeltaK[Bijektiver Transport von Kernfamilien]{%
  \begin{aligned}[t]
    &P\cap A=\varnothing\dsep Q\cap B=\varnothing\dsep
      f\colon A\bij B\vdash{}\\[-2pt]
    &\CoreTransport{P}{Q}{f}\colon
      \CoreFamily{P}{A}\bij\CoreFamily{Q}{B}
  \end{aligned}%
}{CoreFamilyTransportBijective}{%
  \DeltaRow{Mengen}{P\dsep Q\dsep A\dsep B\dsep H\dsep K\dsep J}
  \DeltaRow{Funktionen}{f\dsep\CoreTransport{P}{Q}{f}\dsep
    \CoreTransport{Q}{P}{f^{-1}}}
}
\begin{tabproofsplitwide}
  \proofpartwideindR[Injektivität]{%
    \begin{aligned}[t]
      &P\cap A=\varnothing\dsep Q\cap B=\varnothing\dsep
        f\colon A\bij B\dsep H,K\in\CoreFamily{P}{A}\dsep{}\\[-2pt]
      &\CoreTransport{P}{Q}{f}(H)=\CoreTransport{P}{Q}{f}(K)
        \vdash H=K
    \end{aligned}%
  }{%
    P\cap A=\varnothing\dsep Q\cap B=\varnothing\dsep
    f\colon A\bij B\dsep H,K\in\CoreFamily{P}{A}\dsep
    \CoreTransport{P}{Q}{f}(H)=\CoreTransport{P}{Q}{f}(K)
    \vdash H=K
  }[CoreFamilyTransportInjective]
    \proofstepwidestar[1]{P\cap A=\varnothing}{\rA}
    \proofstepwidestar[2]{Q\cap B=\varnothing}{\rA}
    \proofstepwidestar[3]{f\colon A\bij B}{\rA}
    \proofstepwidestar[4]{H,K\in\CoreFamily{P}{A}}{\rA}
    \proofstepwidestar[5]{\CoreTransport{P}{Q}{f}(H)=
      \CoreTransport{P}{Q}{f}(K)}{\rA}
    \proofstepwidestar[3]{f^{-1}\colon B\to A}{%
      \FormulaRefAuto{F\colon A\bij B\vdash
        F^{-1}\colon B\to A}[thm]{3}}
    \proofstepwidestar[3]{f\colon A\to B}{%
      \FormulaRefAuto{Bijektive Funktion}[def]{3}}
    \proofstepwidestar[3]{\forall x\in A\,f^{-1}(f(x))=x}{%
      \FormulaRefAuto{x\in A \vdash F^{-1}(F(x))=x}[thm]{3}}
    \proofstepwidestar[1,2,3]{\CoreTransport{P}{Q}{f}\colon
      \CoreFamily{P}{A}\to\CoreFamily{Q}{B}}{%
      \FormulaRefAuto{CoreTransportMapping}[thm]{1,2,7}}
    \proofstepwidestar[1,2,3]{\CoreTransport{Q}{P}{f^{-1}}\colon
      \CoreFamily{Q}{B}\to\CoreFamily{P}{A}}{%
      \FormulaRefAuto{CoreTransportMapping}[thm]{2,1,6}}
    \proofstepwidestar[1,2,3,4]{\CoreTransport{P}{Q}{f}(H)
      \in\CoreFamily{Q}{B}}{%
      \FormulaRefAuto{F\colon A\to B\dsep x\in A\vdash F(x)\in B}[thm]{%
        9,\rAEa{4}}}
    \proofstepwidestar[1,2,3,4]{\CoreTransport{P}{Q}{f}(K)
      \in\CoreFamily{Q}{B}}{%
      \FormulaRefAuto{F\colon A\to B\dsep x\in A\vdash F(x)\in B}[thm]{%
        9,\rAEb{4}}}
    \proofstepwidestarbreakkeep[1,2,3,4,5]{%
      \CoreTransport{Q}{P}{f^{-1}}(\CoreTransport{P}{Q}{f}(H))
      =\CoreTransport{Q}{P}{f^{-1}}(\CoreTransport{P}{Q}{f}(K))}{%
      \FormulaRefAuto{F\colon A\to B\dsep x\in A\dsep y\in A\dsep x=y\vdash F(x)=F(y)}[thm]{%
        10,11,12,5}}
    \proofstepwidestar[1,2,3,4]{%
      \CoreTransport{Q}{P}{f^{-1}}(\CoreTransport{P}{Q}{f}(H))=H}{%
      \FormulaRefAuto{CoreTransportRoundTrip}[thm]{1,2,7,6,8,\rAEa{4}}}
    \proofstepwidestar[1,2,3,4]{%
      \CoreTransport{Q}{P}{f^{-1}}(\CoreTransport{P}{Q}{f}(K))=K}{%
      \FormulaRefAuto{CoreTransportRoundTrip}[thm]{1,2,7,6,8,\rAEb{4}}}
    \proofstepwidestar[1,2,3,4,5]{H=K}{%
      \FormulaRefAuto{a=b\dsep a=c\dsep b=d\vdash c=d}[thm]{%
        13,14,15}}
  \closeproofpartwideindR

  \proofpartwideindR[Surjektivität]{%
    \begin{aligned}[t]
      &P\cap A=\varnothing\dsep Q\cap B=\varnothing\dsep
        f\colon A\bij B\dsep J\in\CoreFamily{Q}{B}\vdash{}\\[-2pt]
      &\exists H\in\CoreFamily{P}{A}\,
        \CoreTransport{P}{Q}{f}(H)=J
    \end{aligned}%
  }{%
    P\cap A=\varnothing\dsep Q\cap B=\varnothing\dsep
    f\colon A\bij B\dsep J\in\CoreFamily{Q}{B}\vdash
    \exists H\in\CoreFamily{P}{A}\,
    \CoreTransport{P}{Q}{f}(H)=J
  }[CoreFamilyTransportSurjective]
    \proofstepwidestar[1]{P\cap A=\varnothing}{\rA}
    \proofstepwidestar[2]{Q\cap B=\varnothing}{\rA}
    \proofstepwidestar[3]{f\colon A\bij B}{\rA}
    \proofstepwidestar[4]{J\in\CoreFamily{Q}{B}}{\rA}
    \proofstepwidestar[3]{f^{-1}\colon B\bij A}{%
      \FormulaRefAuto{F\colon A\bij B\vdash F^{-1}\colon B\bij A}[thm]{3}}
    \proofstepwidestar[3]{f^{-1}\colon B\to A}{%
      \FormulaRefAuto{Bijektive Funktion}[def]{5}}
    \proofstepwidestar[3]{f\colon A\to B}{%
      \FormulaRefAuto{Bijektive Funktion}[def]{3}}
    \proofstepwidestar[2,1,3]{\CoreTransport{Q}{P}{f^{-1}}\colon
      \CoreFamily{Q}{B}\to\CoreFamily{P}{A}}{%
      \FormulaRefAuto{CoreTransportMapping}[thm]{2,1,6}}
    \proofstepwidestar[1,2,3,4]{\CoreTransport{Q}{P}{f^{-1}}(J)
      \in\CoreFamily{P}{A}}{%
      \FormulaRefAuto{F\colon A\to B\dsep x\in A\vdash F(x)\in B}[thm]{8,4}}
    \proofstepwidestar[3]{\forall y\in B\,f(f^{-1}(y))=y}{%
      \FormulaRefAuto{F\colon A\bij B\dsep y\in B\vdash
        F(F^{-1}(y))=y}[thm]{3}}
    \proofstepwidestar[1,2,3,4]{%
      \CoreTransport{P}{Q}{f}(\CoreTransport{Q}{P}{f^{-1}}(J))=J}{%
      \FormulaRefAuto{CoreTransportRoundTrip}[thm]{2,1,6,7,10,4}}
    \proofstepwidestar[1,2,3,4]{\exists H\in\CoreFamily{P}{A}\,
      \CoreTransport{P}{Q}{f}(H)=J}{\rEI{\rAI{9,11}}}
  \closeproofpartwideindR

  \proofpartwide{Schluss}
    \proofstepwidestar[1]{P\cap A=\varnothing}{\rA}
    \proofstepwidestar[2]{Q\cap B=\varnothing}{\rA}
    \proofstepwidestar[3]{f\colon A\bij B}{\rA}
    \proofstepwidestar[3]{f \colon A \to B}{%
      \FormulaRefAuto{Bijektive Funktion}[def]{3}}
    \proofstepwidestar[1,2,3]{\CoreTransport{P}{Q}{f}\colon
      \CoreFamily{P}{A}\to\CoreFamily{Q}{B}}{%
      \FormulaRefAuto{CoreTransportMapping}[thm]{1,2,
        4}}
    \proofstepwidestar[6]{H,K\in\CoreFamily{P}{A}}{\rA}
    \proofstepwidestar[7]{\CoreTransport{P}{Q}{f}(H)=
      \CoreTransport{P}{Q}{f}(K)}{\rA}
    \proofstepwidestar[1,2,3,6,7]{H=K}{%
      \FormulaRefAuto{CoreFamilyTransportInjective}[thm]{1,2,3,6,7}}
    \proofstepwidestar[1,2,3]{\CoreTransport{P}{Q}{f}\colon
      \CoreFamily{P}{A}\inj\CoreFamily{Q}{B}}{%
      \FormulaRefAuto{Injektive Funktion}[def]{5,
        \rUI{\rRI{6,\rRI{7,8}}}}}
    \proofstepwidestar[10]{J\in\CoreFamily{Q}{B}}{\rA}
    \proofstepwidestarbreakkeep[1,2,3,10]{\exists H\in\CoreFamily{P}{A}\,
      \CoreTransport{P}{Q}{f}(H)=J}{%
      \FormulaRefAuto{CoreFamilyTransportSurjective}[thm]{1,2,3,10}}
    \proofstepwidestar[1,2,3]{\CoreTransport{P}{Q}{f}\colon
      \CoreFamily{P}{A}\sur\CoreFamily{Q}{B}}{%
      \FormulaRefAuto{Surjektive Funktion}[def]{5,\rUI{\rRI{10,11}}}}
    \proofstepwidestar[1,2,3]{\CoreTransport{P}{Q}{f}\colon
      \CoreFamily{P}{A}\bij\CoreFamily{Q}{B}}{%
      \FormulaRefAuto{F\colon A\inj B\dsep F\colon A\sur B
        \vdash F\colon A\bij B}[thm]{9,12}}
  \closeproofpartwide
\end{tabproofsplitwide}

\FormulaThmDeltaK[Bijektive Absorption einer fremden Kernfamilie]{%
  \begin{aligned}[t]
    &P\cap A=\varnothing\dsep c\notin P\cup A\dsep
      Q\cap B=\varnothing\dsep
      \sigma\colon A\bij A\cup\{c\}\dsep
      \theta\colon A\bij B\dsep{}\\[-2pt]
    &\CoreFamily{P}{A\cup\{c\}}\cap\CoreFamily{Q}{B}=\varnothing
      \vdash{}\\[-2pt]
    &\CaseFun
      {\CoreFamily{P}{A}}
      {\CoreFamily{P\cup\{c\}}{A}}
      {\CoreTransport{P}{P}{\sigma}}
      {\CoreTransport{P\cup\{c\}}{Q}{\theta}}
      \colon\CoreFamily{P}{A\cup\{c\}}
      \bij\CoreFamily{P}{A\cup\{c\}}\cup\CoreFamily{Q}{B}
  \end{aligned}%
}{CoreFamilyAbsorptionBijective}{%
  \DeltaRow{Mengen}{P\dsep Q\dsep A\dsep B\dsep c}
  \DeltaRow{Kernfamilien}{\CoreFamily{P}{A}\dsep
    \CoreFamily{P\cup\{c\}}{A}\dsep
    \CoreFamily{P}{A\cup\{c\}}\dsep\CoreFamily{Q}{B}}
  \DeltaRow{Funktionen}{\sigma\dsep\theta}
  \DeltaRow{Kerntransporte}{\CoreTransport{P}{P}{\sigma}\dsep
    \CoreTransport{P\cup\{c\}}{Q}{\theta}}
  \DeltaRow{Fallfunktion}{\CaseFun
    {\CoreFamily{P}{A}}
    {\CoreFamily{P\cup\{c\}}{A}}
    {\CoreTransport{P}{P}{\sigma}}
    {\CoreTransport{P\cup\{c\}}{Q}{\theta}}}
}
\begin{tabproofwide}
  \proofstepwidestar[1]{P\cap A=\varnothing}{\rA}
  \proofstepwidestar[2]{c\notin P\cup A}{\rA}
  \proofstepwidestar[3]{Q\cap B=\varnothing}{\rA}
  \proofstepwidestar[4]{\sigma\colon A\bij A\cup\{c\}}{\rA}
  \proofstepwidestar[5]{\theta\colon A\bij B}{\rA}
  \proofstepwidestar[6]{%
    \CoreFamily{P}{A\cup\{c\}}\cap\CoreFamily{Q}{B}=\varnothing}{\rA}
  \proofstepwidestar[2]{c\notin P\land c\notin A}{%
    \FormulaRefAuto{OutsideUnionComponents}[thm]{2}}
  \proofstepwidestar[2]{c\notin P}{\rAEa{7}}
  \proofstepwidestar[2]{c\notin A}{\rAEb{7}}
  \proofstepwidestar[2]{P\cap\{c\}=\varnothing}{%
    \FormulaRefAuto{AvoidedSingletonDisjoint}[thm]{8}}
  \proofstepwidestar[]{P\cap(A\cup\{c\})=(P\cap A)\cup(P\cap\{c\})}{%
      \FormulaRefAuto{A\cap(B\cup C)=(A\cap B)\cup(A\cap C)}[thm]}
  \proofstepwidestar[]{\varnothing \cup \varnothing = \varnothing}{%
      \FormulaRefAuto{\varnothing\cup\varnothing=\varnothing}[thm]}
  \proofstepwidestar[1,2]{P\cap(A\cup\{c\})=\varnothing}{%
    \rChain{%
      11,
      \rLRS{1},
      \rLRS{10},
      12}}
  \proofstepwidestar[2]{(P\cup\{c\})\cap A=P\cap A}{%
    \FormulaRefAuto{AdjoinOutsideIntersection}[thm]{9}}
  \proofstepwidestar[1,2]{(P\cup\{c\})\cap A=\varnothing}{%
    \rChain{14,1}}
  \proofstepwidestar[2]{%
    \CoreFamily{P}{A\cup\{c\}}
      =\CoreFamily{P}{A}\cup\CoreFamily{P\cup\{c\}}{A}}{%
    \FormulaRefAuto{CoreFamilyAdjoinedSplitCover}[thm]{2}}
  \proofstepwidestar[2]{%
    \CoreFamily{P}{A}\cap\CoreFamily{P\cup\{c\}}{A}=\varnothing}{%
    \FormulaRefAuto{CoreFamilyAdjoinedSplitDisjoint}[thm]{2}}
  \proofstepwidestar[1,2,4]{\CoreTransport{P}{P}{\sigma}\colon
    \CoreFamily{P}{A}\bij\CoreFamily{P}{A\cup\{c\}}}{%
    \FormulaRefAuto{CoreFamilyTransportBijective}[thm]{1,13,4}}
  \proofstepwidestar[1,2,3,5]{\CoreTransport{P\cup\{c\}}{Q}{\theta}\colon
    \CoreFamily{P\cup\{c\}}{A}\bij\CoreFamily{Q}{B}}{%
    \FormulaRefAuto{CoreFamilyTransportBijective}[thm]{15,3,5}}
  \proofstepwidestar[1,2,3,4,5,6]{%
    \begin{aligned}[t]
      &\CaseFun
        {\CoreFamily{P}{A}}
        {\CoreFamily{P\cup\{c\}}{A}}
        {\CoreTransport{P}{P}{\sigma}}
        {\CoreTransport{P\cup\{c\}}{Q}{\theta}}\\[-2pt]
      &\quad\colon\CoreFamily{P}{A}\cup\CoreFamily{P\cup\{c\}}{A}\\[-2pt]
      &\qquad\bij
        \CoreFamily{P}{A\cup\{c\}}\cup\CoreFamily{Q}{B}
    \end{aligned}}{%
    \FormulaRefAuto{CaseFunBijectiveDisjointTargets}[thm]{17,6,18,19}}
  \proofstepwidestar[1,2,3,4,5,6]{%
    \begin{aligned}[t]
      &\CaseFun
        {\CoreFamily{P}{A}}
        {\CoreFamily{P\cup\{c\}}{A}}
        {\CoreTransport{P}{P}{\sigma}}
        {\CoreTransport{P\cup\{c\}}{Q}{\theta}}\\[-2pt]
      &\quad\colon\CoreFamily{P}{A\cup\{c\}}\\[-2pt]
      &\qquad\bij
        \CoreFamily{P}{A\cup\{c\}}\cup\CoreFamily{Q}{B}
    \end{aligned}}{%
    \rIE{16,20}}
\end{tabproofwide}

\newpage
\FormulaThmDeltaK[Gesch\"utzte Fortsetzung einer lokalen Potenzmengenabsorption]{%
  \begin{aligned}[t]
    &k\notin D\dsep
      \mathcal C\subseteq\powerset(D)\dsep
      \varnothing\notin\mathcal C\dsep{}\\[-2pt]
    &\psi\colon\mathcal C\bij
      \mathcal C\cup\mathfrak{I}[\{k\},D\cup\{k\}]\dsep
      \mathcal O=\powerset(D)\setminus\mathcal C\dsep{}\\[-2pt]
    &\eta=\CaseFun{\mathcal O}{\mathcal C}{\Id_{\mathcal O}}{\psi}
      \vdash{}\\[-2pt]
    &\begin{aligned}[t]
      &\text{(i)}\;\eta\colon\powerset(D)\bij\powerset(D\cup\{k\})\\[-2pt]
      &\text{(ii)}\;\forall H\in\powerset(D)\,
        \bigl(H\notin\mathcal C\rightarrow\eta(H)=H\bigr)\\[-2pt]
      &\text{(iii)}\;\eta(\varnothing)=\varnothing
    \end{aligned}
  \end{aligned}%
}{ProtectedPowerSetExtension}{%
  \DeltaRow{Mengen}{D\dsep k\dsep H\dsep
    \mathcal C\dsep\mathcal O\dsep
    \mathfrak{I}[\{k\},D\cup\{k\}]}
  \DeltaRow{Funktionen}{\psi\dsep\eta\dsep\Id_{\mathcal O}}
  \DeltaRow{Fallfunktion}{%
    \CaseFun{\mathcal O}{\mathcal C}{\Id_{\mathcal O}}{\psi}}
}
\begin{tabproofsplitwide}
  \proofpartwideindR[Bijektivität]{%
    \begin{aligned}[t]
      &k\notin D\dsep\mathcal C\subseteq\powerset(D)\dsep
        \varnothing\notin\mathcal C\dsep{}\\[-2pt]
      &\psi\colon\mathcal C\bij
        \mathcal C\cup\mathfrak{I}[\{k\},D\cup\{k\}]\dsep
        \mathcal O=\powerset(D)\setminus\mathcal C\dsep{}\\[-2pt]
      &\eta=\CaseFun{\mathcal O}{\mathcal C}{\Id_{\mathcal O}}{\psi}
        \vdash\eta\colon\powerset(D)\bij\powerset(D\cup\{k\})
    \end{aligned}%
  }{%
    k\notin D\dsep\mathcal C\subseteq\powerset(D)\dsep
    \varnothing\notin\mathcal C\dsep
    \psi\colon\mathcal C\bij
      \mathcal C\cup\mathfrak{I}[\{k\},D\cup\{k\}]\dsep
    \mathcal O=\powerset(D)\setminus\mathcal C\dsep
    \eta=\CaseFun{\mathcal O}{\mathcal C}{\Id_{\mathcal O}}{\psi}
    \vdash\eta\colon\powerset(D)\bij\powerset(D\cup\{k\})%
  }[ProtectedPowerSetExtensionBijective]
    \proofstepwidestar[1]{k\notin D}{\rA}
    \proofstepwidestar[2]{\mathcal C\subseteq\powerset(D)}{\rA}
    \proofstepwidestar[3]{\varnothing\notin\mathcal C}{\rA}
    \proofstepwidestar[4]{\psi\colon\mathcal C\bij
      \mathcal C\cup\mathfrak{I}[\{k\},D\cup\{k\}]}{\rA}
    \proofstepwidestar[5]{\mathcal O=\powerset(D)\setminus\mathcal C}{\rA}
    \proofstepwidestar[6]{%
      \eta=\CaseFun{\mathcal O}{\mathcal C}{\Id_{\mathcal O}}{\psi}}{\rA}
    \proofstepwidestar[]{\powerset(D\cup\{k\})
      =\powerset(D)\cup\mathfrak{I}[\{k\},D\cup\{k\}]}{%
      \FormulaRefAuto{PowerSetAdjoinedPointPartitionCover}[thm]}
    \proofstepwidestar[1]{%
      \powerset(D)\cap\mathfrak{I}[\{k\},D\cup\{k\}]=\varnothing}{%
      \FormulaRefAuto{PowerSetAdjoinedPointPartitionDisjoint}[thm]{1}}
    \proofstepwidestar[5]{\mathcal O\cap\mathcal C=\varnothing}{%
      \FormulaRefAuto{A=B\setminus C\vdash A\cap C=\varnothing}[thm]{5}}
    \proofstepwidestar[5]{\mathcal O\subseteq\powerset(D)}{%
      \FormulaRefAuto{A\setminus B\subseteq A}[thm]{5}}
    \proofstepwidestar[]{%
      \mathfrak{I}[\{k\},D\cup\{k\}]
        \subseteq\mathfrak{I}[\{k\},D\cup\{k\}]}{%
      \FormulaRefAuto{A\subseteq A}[thm]}
    \proofstepwidestar[1,5]{%
      \mathcal O\cap\mathfrak{I}[\{k\},D\cup\{k\}]=\varnothing}{%
      \FormulaRefAuto{SubsetsOfDisjointSets}[thm]{10,11,8}}
    \proofstepwidestar[]{\mathcal O\cap(\mathcal C\cup\mathfrak I[\{k\},D\cup\{k\}])=(\mathcal O\cap\mathcal C)\cup(\mathcal O\cap\mathfrak I[\{k\},D\cup\{k\}])}{%
      \FormulaRefAuto{A\cap(B\cup C)=(A\cap B)\cup(A\cap C)}[thm]}
    \proofstepwidestar[]{\varnothing \cup \varnothing = \varnothing}{%
      \FormulaRefAuto{\varnothing\cup\varnothing=\varnothing}[thm]}
    \proofstepwidestar[1,5]{%
      \mathcal O\cap
        \bigl(\mathcal C\cup\mathfrak{I}[\{k\},D\cup\{k\}]\bigr)
      =\varnothing}{%
      \rChain{%
        13,
        \rLRS{9},\rLRS{12},
        14}}
    \proofstepwidestar[]{\Id_{\mathcal O}\colon\mathcal O\bij\mathcal O}{%
      \FormulaRefAuto{\Id_A\colon A\bij A}[thm]}
    \proofstepwidestar[1,4,5]{%
      \CaseFun{\mathcal O}{\mathcal C}{\Id_{\mathcal O}}{\psi}
        \colon\mathcal O\cup\mathcal C
        \bij\mathcal O\cup
          \bigl(\mathcal C\cup\mathfrak{I}[\{k\},D\cup\{k\}]\bigr)}{%
      \FormulaRefAuto{CaseFunBijectiveDisjointTargets}[thm]{9,15,16,4}}
    \proofstepwidestar[]{(\powerset(D)\setminus\mathcal C)\cup\mathcal C=\mathcal C\cup(\powerset(D)\setminus\mathcal C)}{%
      \FormulaRefAuto{A\cup B=B\cup A}[thm]}
    \proofstepwidestar[2]{\powerset ( D ) = \mathcal C \cup ( \powerset ( D ) \setminus \mathcal C )}{%
      \FormulaRefAuto{UnionWithRelativeRemainder}[thm]{2}}
    \proofstepwidestar[2]{\mathcal C \cup ( \powerset ( D ) \setminus \mathcal C ) = \powerset ( D )}{%
      \FormulaRefAuto{a=b\vdash b=a}[thm]{%
          19}}
    \proofstepwidestar[2,5]{\mathcal O\cup\mathcal C=\powerset(D)}{%
      \rChain{%
        \rLRS{5},
        18,
        20}}
    \proofstepwidestar[]{\powerset ( D ) \cup \mathfrak I [ \{ k \} , D \cup \{ k \} ] = \powerset ( D \cup \{ k \} )}{%
      \FormulaRefAuto{a=b\vdash b=a}[thm]{7}}
    \proofstepwidestar[]{(\mathcal O\cup\mathcal C)\cup\mathfrak I[\{k\},D\cup\{k\}]=\mathcal O\cup(\mathcal C\cup\mathfrak I[\{k\},D\cup\{k\}])}{%
      \FormulaRefAuto{(A\cup B)\cup C=A\cup(B\cup C)}[thm]}
    \proofstepwidestar[]{\mathcal O\cup(\mathcal C\cup\mathfrak I[\{k\},D\cup\{k\}])=(\mathcal O\cup\mathcal C)\cup\mathfrak I[\{k\},D\cup\{k\}]}{%
      \FormulaRefAuto{a=b\vdash b=a}[thm]{%
          23}}
    \proofstepwidestar[1,2,5]{%
      \mathcal O\cup
        \bigl(\mathcal C\cup\mathfrak{I}[\{k\},D\cup\{k\}]\bigr)
      =\powerset(D\cup\{k\})}{%
      \rChain{%
        24,
        \rLRS{21},
        22}}
    \proofstepwidestar[1,2,4,5]{%
      \CaseFun{\mathcal O}{\mathcal C}{\Id_{\mathcal O}}{\psi}
        \colon\powerset(D)\bij\powerset(D\cup\{k\})}{%
      \rIE{21,25,17}}
    \proofstepwidestar[1,2,4,5,6]{%
      \eta\colon\powerset(D)\bij\powerset(D\cup\{k\})}{\rIE{6,26}}
  \closeproofpartwideindR

  \proofpartwideindR[Fixierung außerhalb der geschützten Familie]{%
    \begin{aligned}[t]
      &k\notin D\dsep\mathcal C\subseteq\powerset(D)\dsep
        \varnothing\notin\mathcal C\dsep{}\\[-2pt]
      &\psi\colon\mathcal C\bij
        \mathcal C\cup\mathfrak{I}[\{k\},D\cup\{k\}]\dsep
        \mathcal O=\powerset(D)\setminus\mathcal C\dsep{}\\[-2pt]
      &\eta=\CaseFun{\mathcal O}{\mathcal C}{\Id_{\mathcal O}}{\psi}
        \vdash
        \forall H\in\powerset(D)\,
          \bigl(H\notin\mathcal C\rightarrow\eta(H)=H\bigr)
    \end{aligned}%
  }{%
    k\notin D\dsep\mathcal C\subseteq\powerset(D)\dsep
    \varnothing\notin\mathcal C\dsep
    \psi\colon\mathcal C\bij
      \mathcal C\cup\mathfrak{I}[\{k\},D\cup\{k\}]\dsep
    \mathcal O=\powerset(D)\setminus\mathcal C\dsep
    \eta=\CaseFun{\mathcal O}{\mathcal C}{\Id_{\mathcal O}}{\psi}
    \vdash\forall H\in\powerset(D)\,
      \bigl(H\notin\mathcal C\rightarrow\eta(H)=H\bigr)%
  }[ProtectedPowerSetExtensionFixesOutside]
    \proofstepwidestar[1]{k\notin D}{\rA}
    \proofstepwidestar[2]{\mathcal C\subseteq\powerset(D)}{\rA}
    \proofstepwidestar[3]{\varnothing\notin\mathcal C}{\rA}
    \proofstepwidestar[4]{\psi\colon\mathcal C\bij
      \mathcal C\cup\mathfrak{I}[\{k\},D\cup\{k\}]}{\rA}
    \proofstepwidestar[5]{\mathcal O=\powerset(D)\setminus\mathcal C}{\rA}
    \proofstepwidestar[6]{%
      \eta=\CaseFun{\mathcal O}{\mathcal C}{\Id_{\mathcal O}}{\psi}}{\rA}
    \proofstepwidestar[]{\Id_{\mathcal O}\colon\mathcal O\bij\mathcal O}{%
      \FormulaRefAuto{\Id_A\colon A\bij A}[thm]}
    \proofstepwidestar[]{\mathcal O\subseteq
      \mathcal O\cup
        \bigl(\mathcal C\cup\mathfrak{I}[\{k\},D\cup\{k\}]\bigr)}{%
      \FormulaRefAuto{A\subseteq A\cup B}[thm]}
    \proofstepwidestar[]{\Id_{\mathcal O}\colon\mathcal O\to\mathcal O}{%
      \FormulaRefAuto{Bijektive Funktion}[def]{7}}
    \proofstepwidestar[]{\Id_{\mathcal O}\colon\mathcal O\to
      \mathcal O\cup
        \bigl(\mathcal C\cup\mathfrak{I}[\{k\},D\cup\{k\}]\bigr)}{%
      \FormulaRefAuto{F\colon A\to B\dsep B\subseteq C
        \vdash F\colon A\to C}[thm]{9,8}}
    \proofstepwidestar[4]{\psi\colon\mathcal C\to
      \mathcal C\cup\mathfrak{I}[\{k\},D\cup\{k\}]}{%
      \FormulaRefAuto{Bijektive Funktion}[def]{4}}
    \proofstepwidestar[]{%
      \mathcal C\cup\mathfrak{I}[\{k\},D\cup\{k\}]
      \subseteq\mathcal O\cup
        \bigl(\mathcal C\cup\mathfrak{I}[\{k\},D\cup\{k\}]\bigr)}{%
      \FormulaRefAuto{B\subseteq A\cup B}[thm]}
    \proofstepwidestar[4]{\psi\colon\mathcal C\to
      \mathcal O\cup
        \bigl(\mathcal C\cup\mathfrak{I}[\{k\},D\cup\{k\}]\bigr)}{%
      \FormulaRefAuto{F\colon A\to B\dsep B\subseteq C
        \vdash F\colon A\to C}[thm]{11,12}}
    \proofstepwidestar[14]{H\in\powerset(D)}{\rA}
    \proofstepwidestar[15]{H\notin\mathcal C}{\rA}
    \proofstepwidestar[5,14,15]{H\in\mathcal O}{%
      \FormulaRefAuto{x\in A\setminus B\eqvdash
        x\in A\land x\notin B}[thm]{\rAI{14,15},5}}
    \proofstepwidestar[4,5,14,15]{%
      \CaseFun{\mathcal O}{\mathcal C}{\Id_{\mathcal O}}{\psi}(H)
        =\Id_{\mathcal O}(H)}{%
      \FormulaRefAuto{f\colon A\to C\dsep g\colon B\to C\dsep x\in A
        \vdash\CaseFun{A}{B}{f}{g}(x)=f(x)}[thm]{10,13,16}}
    \proofstepwidestar[4,5,6,14,15]{\eta(H)=\Id_{\mathcal O}(H)}{%
      \rIE{6,17}}
    \proofstepwidestar[5,14,15]{\Id_{\mathcal O}(H)=H}{%
      \FormulaRefAuto{x\in A\vdash\Id_A(x)=x}[thm]{16}}
    \proofstepwidestar[4,5,6,14,15]{\eta(H)=H}{\rChain{18,19}}
    \proofstepwidestar[4,5,6,14]{%
      H\notin\mathcal C\rightarrow\eta(H)=H}{\rRI{15,20}}
    \proofstepwidestar[4,5,6]{%
      \forall H\in\powerset(D)\,
        \bigl(H\notin\mathcal C\rightarrow\eta(H)=H\bigr)}{%
      \rUI{\rRI{14,21}}}
  \closeproofpartwideindR

  \proofpartwideindR[Fixierung der leeren Menge]{%
    \begin{aligned}[t]
      &k\notin D\dsep\mathcal C\subseteq\powerset(D)\dsep
        \varnothing\notin\mathcal C\dsep{}\\[-2pt]
      &\psi\colon\mathcal C\bij
        \mathcal C\cup\mathfrak{I}[\{k\},D\cup\{k\}]\dsep
        \mathcal O=\powerset(D)\setminus\mathcal C\dsep{}\\[-2pt]
      &\eta=\CaseFun{\mathcal O}{\mathcal C}{\Id_{\mathcal O}}{\psi}
        \vdash\eta(\varnothing)=\varnothing
    \end{aligned}%
  }{%
    k\notin D\dsep\mathcal C\subseteq\powerset(D)\dsep
    \varnothing\notin\mathcal C\dsep
    \psi\colon\mathcal C\bij
      \mathcal C\cup\mathfrak{I}[\{k\},D\cup\{k\}]\dsep
    \mathcal O=\powerset(D)\setminus\mathcal C\dsep
    \eta=\CaseFun{\mathcal O}{\mathcal C}{\Id_{\mathcal O}}{\psi}
    \vdash\eta(\varnothing)=\varnothing%
  }[ProtectedPowerSetExtensionFixesEmpty]
    \proofstepwidestar[1]{k\notin D}{\rA}
    \proofstepwidestar[2]{\mathcal C\subseteq\powerset(D)}{\rA}
    \proofstepwidestar[3]{\varnothing\notin\mathcal C}{\rA}
    \proofstepwidestar[4]{\psi\colon\mathcal C\bij
      \mathcal C\cup\mathfrak{I}[\{k\},D\cup\{k\}]}{\rA}
    \proofstepwidestar[5]{\mathcal O=\powerset(D)\setminus\mathcal C}{\rA}
    \proofstepwidestar[6]{%
      \eta=\CaseFun{\mathcal O}{\mathcal C}{\Id_{\mathcal O}}{\psi}}{\rA}
    \proofstepwidestar[1,2,3,4,5,6]{%
      \forall H\in\powerset(D)\,
        \bigl(H\notin\mathcal C\rightarrow\eta(H)=H\bigr)}{%
      \FormulaRefAuto{ProtectedPowerSetExtensionFixesOutside}[thm]{1,2,3,4,5,6}}
    \proofstepwidestar[]{\varnothing\subseteq D}{%
      \FormulaRefAuto{\varnothing\subseteq A}[thm]}
    \proofstepwidestar[]{\varnothing\in\powerset(D)}{%
      \FormulaRefAuto{x\in\powerset(A)\eqvdash x\subseteq A}[thm]{%
        8}}
    \proofstepwidestar[1,2,3,4,5,6]{%
      \varnothing\notin\mathcal C\rightarrow
        \eta(\varnothing)=\varnothing}{\rRE{\rUE{7},9}}
    \proofstepwidestar[1,2,3,4,5,6]{%
      \eta(\varnothing)=\varnothing}{\rRE{10,3}}
  \closeproofpartwideindR
\end{tabproofsplitwide}

Die lokale Fortsetzung ver\"andert damit nur die gesch\"utzte Familie
\(\mathcal C\).  Die folgende Schichtfortsetzung h\"alt zus\"atzlich eine
disjunkte Tr\"agerschicht punktweise fest.

\subsection{Fortsetzung über eine disjunkte unveränderte Schicht}

\FormulaThmDeltaK[Typisierung der Schichtfortsetzung]{%
  \begin{aligned}[t]
    &g\colon\powerset(D)\to\powerset(E)\dsep
      X\in\powerset(A\cup D)\vdash{}\\[-2pt]
    &(X\cap A)\cup g(X\cap D)\in\powerset(A\cup E)
  \end{aligned}%
}{LayerPowerMapTarget}{%
  \DeltaRow{Mengen}{A\dsep D\dsep E\dsep X}
  \DeltaRow{Funktionen}{g}
}
\begin{tabproofwide}
  \proofstepwidestar[1]{g\colon\powerset(D)\to\powerset(E)}{\rA}
  \proofstepwidestar[2]{X\in\powerset(A\cup D)}{\rA}
  \proofstepwidestar[]{X\cap D\subseteq D}{%
    \FormulaRefAuto{A\cap B\subseteq B}[thm]}
  \proofstepwidestar[]{X\cap D\in\powerset(D)}{%
    \FormulaRefAuto{x \in \powerset(A)\eqvdash x \subseteq A}[thm]{3}}
  \proofstepwidestar[1]{g(X\cap D)\in\powerset(E)}{%
    \FormulaRefAuto{F\colon A\to B\dsep x\in A\vdash F(x)\in B}[thm]{1,4}}
  \proofstepwidestar[1]{g(X\cap D)\subseteq E}{%
    \FormulaRefAuto{x \in \powerset(A)\eqvdash x \subseteq A}[thm]{5}}
  \proofstepwidestar[]{X\cap A\subseteq A}{%
    \FormulaRefAuto{A\cap B\subseteq B}[thm]}
  \proofstepwidestar[1]{(X\cap A)\cup g(X\cap D)\subseteq A\cup E}{%
    \FormulaRefAuto{A\subseteq B\dsep C\subseteq D\vdash
      A\cup C\subseteq B\cup D}[thm]{7,6}}
  \proofstepwidestar[1]{%
    (X\cap A)\cup g(X\cap D)\in\powerset(A\cup E)}{%
    \FormulaRefAuto{x \in \powerset(A)\eqvdash x \subseteq A}[thm]{8}}
\end{tabproofwide}

\FormulaDefDeltaK[Schichtfortsetzung einer Potenzmengenabbildung]{%
  g\colon\powerset(D)\to\powerset(E)\vdash
  \begin{aligned}[t]
    &\LayerPowerMap{A}{D}{E}{g}\colon
      \powerset(A\cup D)\to\powerset(A\cup E),\\[-2pt]
    &\forall X\in\powerset(A\cup D)\,
      \Bigl(\LayerPowerMap{A}{D}{E}{g}(X)\coloneqq
        (X\cap A)\cup g(X\cap D)\Bigr)
  \end{aligned}%
}{LayerPowerMapDef}{%
  \DeltaRow{Mengen}{A\dsep D\dsep E\dsep X}
  \DeltaRow{Funktionen}{g\dsep\LayerPowerMap{A}{D}{E}{g}}
}
\begin{tabproof}
  \proofstep{1}{g\colon\powerset(D)\to\powerset(E)}{\rA}
  \proofstep{2}{X\in\powerset(A\cup D)}{\rA}
  \proofstep{1,2}{(X\cap A)\cup g(X\cap D)\in\powerset(A\cup E)}{%
    \FormulaRefAuto{LayerPowerMapTarget}[thm]{1,2}}
\end{tabproof}

\FormulaThmDeltaK[Typ und Grundgleichung der Schichtfortsetzung]{%
  \begin{aligned}[t]
    &\text{(i)}\;g\colon\powerset(D)\to\powerset(E)\vdash
      \LayerPowerMap{A}{D}{E}{g}\colon
      \powerset(A\cup D)\to\powerset(A\cup E)\\[-2pt]
    &\text{(ii)}\;g\colon\powerset(D)\to\powerset(E)\dsep
      X\in\powerset(A\cup D)\vdash
      \LayerPowerMap{A}{D}{E}{g}(X)=(X\cap A)\cup g(X\cap D)
  \end{aligned}%
}{LayerPowerMapValue}{%
  \DeltaRow{Mengen}{A\dsep D\dsep E\dsep X}
  \DeltaRow{Funktionen}{g\dsep\LayerPowerMap{A}{D}{E}{g}}
}
\begin{tabproofsplitwide}
  \proofpartwideindR[Typ]{%
    g\colon\powerset(D)\to\powerset(E)\vdash
    \LayerPowerMap{A}{D}{E}{g}\colon
      \powerset(A\cup D)\to\powerset(A\cup E)%
  }{%
    g\colon\powerset(D)\to\powerset(E)\vdash
    \LayerPowerMap{A}{D}{E}{g}\colon
      \powerset(A\cup D)\to\powerset(A\cup E)%
  }[LayerPowerMapMapping]
    \proofstepwidestar[1]{g\colon\powerset(D)\to\powerset(E)}{\rA}
    \proofstepwidestar[1]{\LayerPowerMap{A}{D}{E}{g}\colon
      \powerset(A\cup D)\to\powerset(A\cup E)}{%
      \FormulaRefAuto{LayerPowerMapDef}[def]{1}}
  \closeproofpartwideindR

  \proofpartwideindR[Grundgleichung]{%
    g\colon\powerset(D)\to\powerset(E)\dsep
    X\in\powerset(A\cup D)\vdash
    \LayerPowerMap{A}{D}{E}{g}(X)=(X\cap A)\cup g(X\cap D)%
  }{%
    g\colon\powerset(D)\to\powerset(E)\dsep
    X\in\powerset(A\cup D)\vdash
    \LayerPowerMap{A}{D}{E}{g}(X)=(X\cap A)\cup g(X\cap D)%
  }[LayerPowerMapEvaluation]
    \proofstepwidestar[1]{g\colon\powerset(D)\to\powerset(E)}{\rA}
    \proofstepwidestar[2]{X\in\powerset(A\cup D)}{\rA}
    \proofstepwidestar[1]{%
      \forall Y\in\powerset(A\cup D)\,
        \bigl(\LayerPowerMap{A}{D}{E}{g}(Y)
          =(Y\cap A)\cup g(Y\cap D)\bigr)}{%
      \FormulaRefAuto{LayerPowerMapDef}[def]{1}}
    \proofstepwidestar[1,2]{\LayerPowerMap{A}{D}{E}{g}(X)=
      (X\cap A)\cup g(X\cap D)}{\rRE{\rUE{3},2}}
  \closeproofpartwideindR
\end{tabproofsplitwide}

\FormulaThmDeltaK[Zerlegung einer Menge in zwei Schichten]{%
  X\in\powerset(A\cup D)\vdash
  X=(X\cap A)\cup(X\cap D)
}{LayerSetDecomposition}{%
  \DeltaRow{Mengen}{A\dsep D\dsep X}
}
\begin{tabproofwide}
  \proofstepwidestar[1]{X\in\powerset(A\cup D)}{\rA}
  \proofstepwidestar[1]{X\subseteq A\cup D}{%
    \FormulaRefAuto{x \in \powerset(A)\eqvdash x \subseteq A}[thm]{1}}
  \proofstepwidestar[1]{X\cap(A\cup D)=X}{%
    \FormulaRefAuto{A \subseteq B \eqvdash A \cap B = A}[thm]{2}}
  \proofstepwidestarbreakkeep[]{X\cap(A\cup D)=(X\cap A)\cup(X\cap D)}{%
    \FormulaRefAuto{A \cap (B \cup C) = (A \cap B) \cup (A \cap C)}[thm]}
  \proofstepwidestar[1]{X=(X\cap A)\cup(X\cap D)}{%
    \rChain{\rLRS{3},4}}
\end{tabproofwide}

\FormulaThmDeltaK[Die feste Schicht bleibt unverändert]{%
  \begin{aligned}[t]
    &A\cap E=\varnothing\dsep g\colon\powerset(D)\to\powerset(E)
      \dsep X\in\powerset(A\cup D)\vdash{}\\[-2pt]
    &\LayerPowerMap{A}{D}{E}{g}(X)\cap A=X\cap A
  \end{aligned}%
}{LayerPowerMapStablePart}{%
  \DeltaRow{Mengen}{A\dsep D\dsep E\dsep X\dsep x}
  \DeltaRow{Funktionen}{g\dsep\LayerPowerMap{A}{D}{E}{g}}
}
\begin{tabproofwide}
  \proofstepwidestar[1]{A\cap E=\varnothing}{\rA}
  \proofstepwidestar[2]{g\colon\powerset(D)\to\powerset(E)}{\rA}
  \proofstepwidestar[3]{X\in\powerset(A\cup D)}{\rA}
  \proofstepwidestar[2,3]{\LayerPowerMap{A}{D}{E}{g}(X)=
    (X\cap A)\cup g(X\cap D)}{\FormulaRefAuto{LayerPowerMapEvaluation}[thm]{2,3}}
  \proofstepwidestar[]{X\cap D\subseteq D}{%
      \FormulaRefAuto{A\cap B\subseteq B}[thm]}
  \proofstepwidestar[]{X\cap D\in\powerset(D)}{%
      \FormulaRefAuto{x \in \powerset(A)\eqvdash x \subseteq A}[thm]{%
          5}}
  \proofstepwidestar[2]{g(X\cap D)\in\powerset(E)}{%
      \FormulaRefAuto{F\colon A\to B\dsep x\in A\vdash F(x)\in B}[thm]{%
        2,6}}
  \proofstepwidestarbreakkeep[2]{g(X\cap D)\subseteq E}{%
    \FormulaRefAuto{x \in \powerset(A)\eqvdash x \subseteq A}[thm]{%
      7}}
  \proofstepwidestar[]{A \subseteq A}{%
      \FormulaRefAuto{A\subseteq A}[thm]}
  \proofstepwidestar[1,2]{A\cap g(X\cap D)=\varnothing}{%
    \FormulaRefAuto{SubsetsOfDisjointSets}[thm]{%
      9,8,1}}
  \proofstepwidestar[1,2]{g(X\cap D)\cap A=\varnothing}{%
    \FormulaRefAuto{A\cap B=B\cap A}[thm]{10}}
  \proofstepwidestar[]{X\cap A\subseteq A}{%
      \FormulaRefAuto{A\cap B\subseteq B}[thm]}
  \proofstepwidestar[]{(X\cap A)\cap A=X\cap A}{%
    \FormulaRefAuto{A \subseteq B \eqvdash A \cap B = A}[thm]{%
      12}}
  \proofstepwidestarbreakkeep[]{((X\cap A)\cup g(X\cap D))\cap A
    =((X\cap A)\cap A)\cup(g(X\cap D)\cap A)}{%
    \FormulaRefAuto{(A \cup B) \cap C = (A \cap C) \cup (B \cap C)}[thm]}
  \proofstepwidestar[]{(X\cap A)\cup\varnothing=X\cap A}{%
      \FormulaRefAuto{A\cup\varnothing=A}[thm]}
  \proofstepwidestarbreakkeep[1,2]{((X\cap A)\cup g(X\cap D))\cap A=X\cap A}{%
    \rChain{14,\rLRS{13},\rLRS{11},
      15}}
  \proofstepwidestar[1,2,3]{\LayerPowerMap{A}{D}{E}{g}(X)\cap A=X\cap A}{%
    \rIE{4,16}}
\end{tabproofwide}

\FormulaThmDeltaK[Die veränderliche Zielschicht]{%
  \begin{aligned}[t]
    &A\cap E=\varnothing\dsep g\colon\powerset(D)\to\powerset(E)
      \dsep X\in\powerset(A\cup D)\vdash{}\\[-2pt]
    &\LayerPowerMap{A}{D}{E}{g}(X)\cap E=g(X\cap D)
  \end{aligned}%
}{LayerPowerMapComponent}{%
  \DeltaRow{Mengen}{A\dsep D\dsep E\dsep X}
  \DeltaRow{Funktionen}{g\dsep\LayerPowerMap{A}{D}{E}{g}}
}
\begin{tabproofwide}
  \proofstepwidestar[1]{A\cap E=\varnothing}{\rA}
  \proofstepwidestar[2]{g\colon\powerset(D)\to\powerset(E)}{\rA}
  \proofstepwidestar[3]{X\in\powerset(A\cup D)}{\rA}
  \proofstepwidestar[2,3]{\LayerPowerMap{A}{D}{E}{g}(X)=
    (X\cap A)\cup g(X\cap D)}{\FormulaRefAuto{LayerPowerMapEvaluation}[thm]{2,3}}
  \proofstepwidestar[]{E \subseteq E}{%
      \FormulaRefAuto{A\subseteq A}[thm]}
  \proofstepwidestar[]{X\cap A\subseteq A}{%
      \FormulaRefAuto{A\cap B\subseteq B}[thm]}
  \proofstepwidestarbreakkeep[1]{(X\cap A)\cap E=\varnothing}{%
    \FormulaRefAuto{SubsetsOfDisjointSets}[thm]{%
      6,
      5,1}}
  \proofstepwidestar[]{X\cap D\subseteq D}{%
      \FormulaRefAuto{A\cap B\subseteq B}[thm]}
  \proofstepwidestar[]{X\cap D\in\powerset(D)}{%
      \FormulaRefAuto{x \in \powerset(A)\eqvdash x \subseteq A}[thm]{%
          8}}
  \proofstepwidestar[2]{g(X\cap D)\in\powerset(E)}{%
      \FormulaRefAuto{F\colon A\to B\dsep x\in A\vdash F(x)\in B}[thm]{%
        2,9}}
  \proofstepwidestarbreakkeep[2]{g(X\cap D)\subseteq E}{%
    \FormulaRefAuto{x \in \powerset(A)\eqvdash x \subseteq A}[thm]{%
      10}}
  \proofstepwidestar[2]{g(X\cap D)\cap E=g(X\cap D)}{%
    \FormulaRefAuto{A \subseteq B \eqvdash A \cap B = A}[thm]{11}}
  \proofstepwidestarbreakkeep[]{((X\cap A)\cup g(X\cap D))\cap E
    =((X\cap A)\cap E)\cup(g(X\cap D)\cap E)}{%
    \FormulaRefAuto{(A \cup B) \cap C = (A \cap C) \cup (B \cap C)}[thm]}
  \proofstepwidestar[]{\varnothing\cup g(X\cap D)=g(X\cap D)}{%
      \FormulaRefAuto{\varnothing\cup A=A}[thm]}
  \proofstepwidestar[1,2]{((X\cap A)\cup g(X\cap D))\cap E=g(X\cap D)}{%
    \rChain{13,\rLRS{7},\rLRS{12},
      14}}
  \proofstepwidestar[1,2,3]{\LayerPowerMap{A}{D}{E}{g}(X)\cap E
    =g(X\cap D)}{\rIE{4,15}}
\end{tabproofwide}

\FormulaThmDeltaK[Hin- und Rückweg der Schichtfortsetzung]{%
  \begin{aligned}[t]
    &A\cap D=\varnothing\dsep A\cap E=\varnothing\dsep
      g\colon\powerset(D)\to\powerset(E)\dsep
      h\colon\powerset(E)\to\powerset(D)\dsep{}\\[-2pt]
    &\forall K\in\powerset(D)\,h(g(K))=K\dsep
      X\in\powerset(A\cup D)\vdash{}\\[-2pt]
    &\LayerPowerMap{A}{E}{D}{h}
      (\LayerPowerMap{A}{D}{E}{g}(X))=X
  \end{aligned}%
}{LayerPowerMapRoundTrip}{%
  \DeltaRow{Mengen}{A\dsep D\dsep E\dsep X\dsep K}
  \DeltaRow{Funktionen}{g\dsep h\dsep
    \LayerPowerMap{A}{D}{E}{g}\dsep\LayerPowerMap{A}{E}{D}{h}}
}
\begin{tabproofwide}
  \proofstepwidestar[1]{A\cap D=\varnothing}{\rA}
  \proofstepwidestar[2]{A\cap E=\varnothing}{\rA}
  \proofstepwidestar[3]{g\colon\powerset(D)\to\powerset(E)}{\rA}
  \proofstepwidestar[4]{h\colon\powerset(E)\to\powerset(D)}{\rA}
  \proofstepwidestar[5]{\forall K\in\powerset(D)\,h(g(K))=K}{\rA}
  \proofstepwidestar[6]{X\in\powerset(A\cup D)}{\rA}
  \proofstepwidestar[3]{\LayerPowerMap{A}{D}{E}{g}\colon\powerset(A\cup D)\to\powerset(A\cup E)}{%
      \FormulaRefAuto{LayerPowerMapMapping}[thm]{3}}
  \proofstepwidestar[3,6]{\LayerPowerMap{A}{D}{E}{g}(X)
    \in\powerset(A\cup E)}{%
    \FormulaRefAuto{F\colon A\to B\dsep x\in A\vdash F(x)\in B}[thm]{%
      7,6}}
  \proofstepwidestar[2,3,6]{\LayerPowerMap{A}{D}{E}{g}(X)\cap A=X\cap A}{%
    \FormulaRefAuto{LayerPowerMapStablePart}[thm]{2,3,6}}
  \proofstepwidestarbreakkeep[2,3,6]{\LayerPowerMap{A}{D}{E}{g}(X)\cap E=g(X\cap D)}{%
    \FormulaRefAuto{LayerPowerMapComponent}[thm]{2,3,6}}
  \proofstepwidestarbreakkeep[4,8]{%
    \begin{aligned}[t]
      &\LayerPowerMap{A}{E}{D}{h}(\LayerPowerMap{A}{D}{E}{g}(X))\\[-2pt]
      &\quad=(\LayerPowerMap{A}{D}{E}{g}(X)\cap A)\cup\\[-2pt]
      &\qquad h(\LayerPowerMap{A}{D}{E}{g}(X)\cap E)
    \end{aligned}}{%
    \FormulaRefAuto{LayerPowerMapEvaluation}[thm]{4,8}}
  \proofstepwidestar[]{X\cap D\subseteq D}{%
      \FormulaRefAuto{A\cap B\subseteq B}[thm]}
  \proofstepwidestar[]{X\cap D\in\powerset(D)}{%
    \FormulaRefAuto{x \in \powerset(A)\eqvdash x \subseteq A}[thm]{%
      12}}
  \proofstepwidestar[5]{h(g(X\cap D))=X\cap D}{\rRE{\rUE{5},13}}
  \proofstepwidestarbreakkeep[1,2,3,4,5,6]{%
    \begin{aligned}[t]
      &\LayerPowerMap{A}{E}{D}{h}(\LayerPowerMap{A}{D}{E}{g}(X))\\[-2pt]
      &\quad=(X\cap A)\cup(X\cap D)
    \end{aligned}}{\rChain{11,\rLRS{9},\rLRS{10},\rLRS{14}}}
  \proofstepwidestar[6]{X=(X\cap A)\cup(X\cap D)}{%
    \FormulaRefAuto{LayerSetDecomposition}[thm]{6}}
  \proofstepwidestar[1,2,3,4,5,6]{%
    \LayerPowerMap{A}{E}{D}{h}(\LayerPowerMap{A}{D}{E}{g}(X))=X}{%
    \rChain{15,\rLRS{16}}}
\end{tabproofwide}

\FormulaThmDeltaK[Bijektive Schichtfortsetzung]{%
  \begin{aligned}[t]
    &A\cap D=\varnothing\dsep A\cap E=\varnothing\dsep
      g\colon\powerset(D)\bij\powerset(E)\vdash{}\\[-2pt]
    &\LayerPowerMap{A}{D}{E}{g}\colon
      \powerset(A\cup D)\bij\powerset(A\cup E)
  \end{aligned}%
}{LayerPowerMapBijective}{%
  \DeltaRow{Mengen}{A\dsep D\dsep E\dsep X\dsep Y\dsep Z}
  \DeltaRow{Funktionen}{g\dsep\LayerPowerMap{A}{D}{E}{g}\dsep
    \LayerPowerMap{A}{E}{D}{g^{-1}}}
}
\begin{tabproofsplitwide}
  \proofpartwideindR[Injektivität]{%
    \begin{aligned}[t]
      &A\cap D=\varnothing\dsep A\cap E=\varnothing\dsep
        g\colon\powerset(D)\bij\powerset(E)\dsep X,Y\in\powerset(A\cup D)
        \dsep{}\\[-2pt]
      &\LayerPowerMap{A}{D}{E}{g}(X)=
        \LayerPowerMap{A}{D}{E}{g}(Y)\vdash X=Y
    \end{aligned}%
  }{%
    A\cap D=\varnothing\dsep A\cap E=\varnothing\dsep
    g\colon\powerset(D)\bij\powerset(E)\dsep X,Y\in\powerset(A\cup D)
    \dsep\LayerPowerMap{A}{D}{E}{g}(X)=
    \LayerPowerMap{A}{D}{E}{g}(Y)\vdash X=Y
  }[LayerPowerMapInjective]
    \proofstepwidestar[1]{A\cap D=\varnothing}{\rA}
    \proofstepwidestar[2]{A\cap E=\varnothing}{\rA}
    \proofstepwidestar[3]{g\colon\powerset(D)\bij\powerset(E)}{\rA}
    \proofstepwidestar[4]{X,Y\in\powerset(A\cup D)}{\rA}
    \proofstepwidestar[5]{\LayerPowerMap{A}{D}{E}{g}(X)=
      \LayerPowerMap{A}{D}{E}{g}(Y)}{\rA}
    \proofstepwidestar[3]{g^{-1}\colon\powerset(E)\to\powerset(D)}{%
      \FormulaRefAuto{F\colon A\bij B\vdash
        F^{-1}\colon B\to A}[thm]{3}}
    \proofstepwidestar[3]{g\colon\powerset(D)\to\powerset(E)}{%
      \FormulaRefAuto{Bijektive Funktion}[def]{3}}
    \proofstepwidestar[3]{\forall K\in\powerset(D)\,g^{-1}(g(K))=K}{%
      \FormulaRefAuto{x\in A \vdash F^{-1}(F(x))=x}[thm]{3}}
    \proofstepwidestar[3]{\LayerPowerMap{A}{D}{E}{g}\colon
      \powerset(A\cup D)\to\powerset(A\cup E)}{%
      \FormulaRefAuto{LayerPowerMapMapping}[thm]{7}}
    \proofstepwidestar[3]{\LayerPowerMap{A}{E}{D}{g^{-1}}\colon
      \powerset(A\cup E)\to\powerset(A\cup D)}{%
      \FormulaRefAuto{LayerPowerMapMapping}[thm]{6}}
    \proofstepwidestarbreakkeep[3,4]{\LayerPowerMap{A}{D}{E}{g}(X)
      \in\powerset(A\cup E)}{%
      \FormulaRefAuto{F\colon A\to B\dsep x\in A\vdash F(x)\in B}[thm]{%
        9,\rAEa{4}}}
    \proofstepwidestar[3,4]{\LayerPowerMap{A}{D}{E}{g}(Y)
      \in\powerset(A\cup E)}{%
      \FormulaRefAuto{F\colon A\to B\dsep x\in A\vdash F(x)\in B}[thm]{%
        9,\rAEb{4}}}
    \proofstepwidestarbreakkeep[1,2,3,4,5]{%
      \LayerPowerMap{A}{E}{D}{g^{-1}}(\LayerPowerMap{A}{D}{E}{g}(X))
      =\LayerPowerMap{A}{E}{D}{g^{-1}}(\LayerPowerMap{A}{D}{E}{g}(Y))}{%
      \FormulaRefAuto{F\colon A\to B\dsep x\in A\dsep y\in A\dsep x=y\vdash F(x)=F(y)}[thm]{%
        10,11,12,5}}
    \proofstepwidestarbreakkeep[1,2,3,4]{%
      \LayerPowerMap{A}{E}{D}{g^{-1}}(\LayerPowerMap{A}{D}{E}{g}(X))=X}{%
      \FormulaRefAuto{LayerPowerMapRoundTrip}[thm]{%
        1,2,7,6,8,\rAEa{4}}}
    \proofstepwidestarbreakkeep[1,2,3,4]{%
      \LayerPowerMap{A}{E}{D}{g^{-1}}(\LayerPowerMap{A}{D}{E}{g}(Y))=Y}{%
      \FormulaRefAuto{LayerPowerMapRoundTrip}[thm]{%
        1,2,7,6,8,\rAEb{4}}}
    \proofstepwidestar[1,2,3,4,5]{X=Y}{%
      \FormulaRefAuto{a=b\dsep a=c\dsep b=d\vdash c=d}[thm]{13,14,15}}
  \closeproofpartwideindR

  \proofpartwideindR[Surjektivität]{%
    \begin{aligned}[t]
      &A\cap D=\varnothing\dsep A\cap E=\varnothing\dsep
        g\colon\powerset(D)\bij\powerset(E)\dsep
        Z\in\powerset(A\cup E)\vdash{}\\[-2pt]
      &\exists X\in\powerset(A\cup D)\,
        \LayerPowerMap{A}{D}{E}{g}(X)=Z
    \end{aligned}%
  }{%
    A\cap D=\varnothing\dsep A\cap E=\varnothing\dsep
    g\colon\powerset(D)\bij\powerset(E)\dsep
    Z\in\powerset(A\cup E)\vdash
    \exists X\in\powerset(A\cup D)\,
    \LayerPowerMap{A}{D}{E}{g}(X)=Z
  }[LayerPowerMapSurjective]
    \proofstepwidestar[1]{A\cap D=\varnothing}{\rA}
    \proofstepwidestar[2]{A\cap E=\varnothing}{\rA}
    \proofstepwidestar[3]{g\colon\powerset(D)\bij\powerset(E)}{\rA}
    \proofstepwidestar[4]{Z\in\powerset(A\cup E)}{\rA}
    \proofstepwidestar[3]{g^{-1}\colon\powerset(E)\bij\powerset(D)}{%
      \FormulaRefAuto{F\colon A\bij B\vdash F^{-1}\colon B\bij A}[thm]{3}}
    \proofstepwidestar[3]{g ^ { - 1 } \colon \powerset ( E ) \to \powerset ( D )}{%
      \FormulaRefAuto{Bijektive Funktion}[def]{5}}
    \proofstepwidestarbreakkeep[3]{\LayerPowerMap{A}{E}{D}{g^{-1}}\colon
      \powerset(A\cup E)\to\powerset(A\cup D)}{%
      \FormulaRefAuto{LayerPowerMapMapping}[thm]{%
        6}}
    \proofstepwidestar[3,4]{\LayerPowerMap{A}{E}{D}{g^{-1}}(Z)
      \in\powerset(A\cup D)}{%
      \FormulaRefAuto{F\colon A\to B\dsep x\in A\vdash F(x)\in B}[thm]{7,4}}
    \proofstepwidestar[3]{\forall L\in\powerset(E)\,g(g^{-1}(L))=L}{%
      \FormulaRefAuto{F\colon A\bij B\dsep y\in B\vdash F(F^{-1}(y))=y}[thm]{3}}
    \proofstepwidestar[3]{{ g ^ { - 1 } } \colon \powerset ( E ) \to \powerset ( D )}{%
      \FormulaRefAuto{Bijektive Funktion}[def]{5}}
    \proofstepwidestar[3]{g \colon \powerset ( D ) \to \powerset ( E )}{%
      \FormulaRefAuto{Bijektive Funktion}[def]{3}}
    \proofstepwidestar[1,2,3,4]{%
      \LayerPowerMap{A}{D}{E}{g}(\LayerPowerMap{A}{E}{D}{g^{-1}}(Z))=Z}{%
      \FormulaRefAuto{LayerPowerMapRoundTrip}[thm]{%
        2,1,10,
        11,9,4}}
    \proofstepwidestar[1,2,3,4]{\exists X\in\powerset(A\cup D)\,
      \LayerPowerMap{A}{D}{E}{g}(X)=Z}{\rEI{\rAI{8,12}}}
  \closeproofpartwideindR

  \newpage
  \proofpartwide{Schluss}
    \proofstepwidestar[1]{A\cap D=\varnothing}{\rA}
    \proofstepwidestar[2]{A\cap E=\varnothing}{\rA}
    \proofstepwidestar[3]{g\colon\powerset(D)\bij\powerset(E)}{\rA}
    \proofstepwidestar[3]{g \colon \powerset ( D ) \to \powerset ( E )}{%
      \FormulaRefAuto{Bijektive Funktion}[def]{3}}
    \proofstepwidestarbreakkeep[3]{\LayerPowerMap{A}{D}{E}{g}\colon
      \powerset(A\cup D)\to\powerset(A\cup E)}{%
      \FormulaRefAuto{LayerPowerMapMapping}[thm]{4}}
    \proofstepwidestar[6]{X,Y\in\powerset(A\cup D)}{\rA}
    \proofstepwidestar[7]{\LayerPowerMap{A}{D}{E}{g}(X)=
      \LayerPowerMap{A}{D}{E}{g}(Y)}{\rA}
    \proofstepwidestar[1,2,3,6,7]{X=Y}{%
      \FormulaRefAuto{LayerPowerMapInjective}[thm]{1,2,3,6,7}}
    \proofstepwidestarbreakkeep[1,2,3]{\LayerPowerMap{A}{D}{E}{g}\colon
      \powerset(A\cup D)\inj\powerset(A\cup E)}{%
      \FormulaRefAuto{Injektive Funktion}[def]{5,\rUI{\rRI{6,\rRI{7,8}}}}}
    \proofstepwidestar[10]{Z\in\powerset(A\cup E)}{\rA}
    \proofstepwidestarbreakkeep[1,2,3,10]{\exists X\in\powerset(A\cup D)\,
      \LayerPowerMap{A}{D}{E}{g}(X)=Z}{%
      \FormulaRefAuto{LayerPowerMapSurjective}[thm]{1,2,3,10}}
    \proofstepwidestarbreakkeep[1,2,3]{\LayerPowerMap{A}{D}{E}{g}\colon
      \powerset(A\cup D)\sur\powerset(A\cup E)}{%
      \FormulaRefAuto{Surjektive Funktion}[def]{5,\rUI{\rRI{10,11}}}}
    \proofstepwidestarbreakkeep[1,2,3]{\LayerPowerMap{A}{D}{E}{g}\colon
      \powerset(A\cup D)\bij\powerset(A\cup E)}{%
      \FormulaRefAuto{F\colon A\inj B\dsep F\colon A\sur B
        \vdash F\colon A\bij B}[thm]{9,12}}
  \closeproofpartwide
\end{tabproofsplitwide}

\FormulaThmDeltaK[Fixierung bei unveränderter variabler Schicht]{%
  \begin{aligned}[t]
    &g\colon\powerset(D)\to\powerset(E)\dsep
      X\in\powerset(A\cup D)\dsep g(X\cap D)=X\cap D\vdash{}\\[-2pt]
    &\LayerPowerMap{A}{D}{E}{g}(X)=X
  \end{aligned}%
}{LayerPowerMapFixesStableUnion}{%
  \DeltaRow{Mengen}{A\dsep D\dsep E\dsep X}
  \DeltaRow{Funktionen}{g\dsep\LayerPowerMap{A}{D}{E}{g}}
}
\begin{tabproofwide}
  \proofstepwidestar[1]{g\colon\powerset(D)\to\powerset(E)}{\rA}
  \proofstepwidestar[2]{X\in\powerset(A\cup D)}{\rA}
  \proofstepwidestar[3]{g(X\cap D)=X\cap D}{\rA}
  \proofstepwidestarbreakkeep[1,2]{\LayerPowerMap{A}{D}{E}{g}(X)=
    (X\cap A)\cup g(X\cap D)}{\FormulaRefAuto{LayerPowerMapEvaluation}[thm]{1,2}}
  \proofstepwidestar[1,2,3]{\LayerPowerMap{A}{D}{E}{g}(X)=
    (X\cap A)\cup(X\cap D)}{\rLRS{3,4}}
  \proofstepwidestar[2]{X=(X\cap A)\cup(X\cap D)}{%
    \FormulaRefAuto{LayerSetDecomposition}[thm]{2}}
  \proofstepwidestar[1,2,3]{\LayerPowerMap{A}{D}{E}{g}(X)=X}{%
    \rChain{5,\rLRS{6}}}
\end{tabproofwide}

\newpage
\FormulaThmDeltaK[Teilmengen der festen Schicht werden fixiert]{%
  \begin{aligned}[t]
    &A\cap D=\varnothing\dsep g\colon\powerset(D)\to\powerset(E)
      \dsep g(\varnothing)=\varnothing\dsep L\subseteq A\dsep{}\\[-2pt]
    &X\in\powerset(A\cup D)\dsep X\subseteq L\vdash
      \LayerPowerMap{A}{D}{E}{g}(X)=X
  \end{aligned}%
}{LayerPowerMapFixesStableSubsets}{%
  \DeltaRow{Mengen}{A\dsep D\dsep E\dsep L\dsep X}
  \DeltaRow{Funktionen}{g\dsep\LayerPowerMap{A}{D}{E}{g}}
}
\begin{tabproofwide}
  \proofstepwidestar[1]{A\cap D=\varnothing}{\rA}
  \proofstepwidestar[2]{g\colon\powerset(D)\to\powerset(E)}{\rA}
  \proofstepwidestar[3]{g(\varnothing)=\varnothing}{\rA}
  \proofstepwidestar[4]{L\subseteq A}{\rA}
  \proofstepwidestar[5]{X\in\powerset(A\cup D)}{\rA}
  \proofstepwidestar[6]{X\subseteq L}{\rA}
  \proofstepwidestar[4,6]{X\subseteq A}{%
    \FormulaRefAuto{A\subseteq B\dsep B\subseteq C\vdash A\subseteq C}[thm]{6,4}}
  \proofstepwidestar[]{D \subseteq D}{%
      \FormulaRefAuto{A\subseteq A}[thm]}
  \proofstepwidestar[1,4,6]{X\cap D=\varnothing}{%
    \FormulaRefAuto{SubsetsOfDisjointSets}[thm]{7,
      8,1}}
  \proofstepwidestar[2,3,4,6]{g(X\cap D)=X\cap D}{%
    \rIE{9,\rIE{3,\rII}}}
  \proofstepwidestarbreakkeep[1,2,3,4,5,6]{\LayerPowerMap{A}{D}{E}{g}(X)=X}{%
    \FormulaRefAuto{LayerPowerMapFixesStableUnion}[thm]{2,5,10}}
\end{tabproofwide}

\FormulaThmDeltaK[Erhaltung von Teilmengen der festen Schicht]{%
  \begin{aligned}[t]
    &A\cap D=\varnothing\dsep A\cap E=\varnothing\dsep
      g\colon\powerset(D)\bij\powerset(E)\dsep
      g(\varnothing)=\varnothing\dsep{}\\[-2pt]
    &L\subseteq A\dsep X\in\powerset(A\cup D)\vdash
      \bigl(X\subseteq L\leftrightarrow
      \LayerPowerMap{A}{D}{E}{g}(X)\subseteq L\bigr)
  \end{aligned}%
}{LayerPowerMapPreservesSubsets}{%
  \DeltaRow{Mengen}{A\dsep D\dsep E\dsep L\dsep X}
  \DeltaRow{Funktionen}{g\dsep\LayerPowerMap{A}{D}{E}{g}\dsep
    \LayerPowerMap{A}{E}{D}{g^{-1}}}
}
\begin{tabproofsplitwide}
  \proofpartwideindR[Hinrichtung]{%
    \begin{aligned}[t]
      &A\cap D=\varnothing\dsep g\colon\powerset(D)\to\powerset(E)
        \dsep g(\varnothing)=\varnothing\dsep L\subseteq A\dsep{}\\[-2pt]
      &X\in\powerset(A\cup D)\dsep X\subseteq L\vdash
        \LayerPowerMap{A}{D}{E}{g}(X)\subseteq L
    \end{aligned}%
  }{%
    A\cap D=\varnothing\dsep g\colon\powerset(D)\to\powerset(E)
    \dsep g(\varnothing)=\varnothing\dsep L\subseteq A\dsep
    X\in\powerset(A\cup D)\dsep X\subseteq L\vdash
    \LayerPowerMap{A}{D}{E}{g}(X)\subseteq L
  }[LayerPowerMapPreservesSubsetsForward]
    \proofstepwidestar[1]{A\cap D=\varnothing}{\rA}
    \proofstepwidestar[2]{g\colon\powerset(D)\to\powerset(E)}{\rA}
    \proofstepwidestar[3]{g(\varnothing)=\varnothing}{\rA}
    \proofstepwidestar[4]{L\subseteq A}{\rA}
    \proofstepwidestar[5]{X\in\powerset(A\cup D)}{\rA}
    \proofstepwidestar[6]{X\subseteq L}{\rA}
    \proofstepwidestarbreakkeep[1,2,3,4,5,6]{\LayerPowerMap{A}{D}{E}{g}(X)=X}{%
      \FormulaRefAuto{LayerPowerMapFixesStableSubsets}[thm]{1,2,3,4,5,6}}
    \proofstepwidestar[1,2,3,4,5,6]{\LayerPowerMap{A}{D}{E}{g}(X)\subseteq L}{%
      \rIE{7,6}}
  \closeproofpartwideindR

  \proofpartwideindR[Rückrichtung]{%
    \begin{aligned}[t]
      &A\cap D=\varnothing\dsep A\cap E=\varnothing\dsep
        g\colon\powerset(D)\bij\powerset(E)\dsep
        g(\varnothing)=\varnothing\dsep L\subseteq A\dsep{}\\[-2pt]
      &X\in\powerset(A\cup D)\dsep
        \LayerPowerMap{A}{D}{E}{g}(X)\subseteq L\vdash X\subseteq L
    \end{aligned}%
  }{%
    A\cap D=\varnothing\dsep A\cap E=\varnothing\dsep
    g\colon\powerset(D)\bij\powerset(E)\dsep
    g(\varnothing)=\varnothing\dsep L\subseteq A\dsep
    X\in\powerset(A\cup D)\dsep
    \LayerPowerMap{A}{D}{E}{g}(X)\subseteq L\vdash X\subseteq L
  }[LayerPowerMapPreservesSubsetsBackward]
    \proofstepwidestar[1]{A\cap D=\varnothing}{\rA}
    \proofstepwidestar[2]{A\cap E=\varnothing}{\rA}
    \proofstepwidestar[3]{g\colon\powerset(D)\bij\powerset(E)}{\rA}
    \proofstepwidestar[4]{g(\varnothing)=\varnothing}{\rA}
    \proofstepwidestar[5]{L\subseteq A}{\rA}
    \proofstepwidestar[6]{X\in\powerset(A\cup D)}{\rA}
    \proofstepwidestar[7]{\LayerPowerMap{A}{D}{E}{g}(X)\subseteq L}{\rA}
    \proofstepwidestar[3]{g^{-1}\colon\powerset(E)\to\powerset(D)}{%
      \FormulaRefAuto{F\colon A\bij B\vdash
        F^{-1}\colon B\to A}[thm]{3}}
    \proofstepwidestar[]{\varnothing\subseteq D}{%
      \FormulaRefAuto{\varnothing\subseteq A}[thm]}
    \proofstepwidestar[]{\varnothing\in\powerset(D)}{%
      \FormulaRefAuto{x\in\powerset(A)\eqvdash x\subseteq A}[thm]{9}}
    \proofstepwidestar[3]{g^{-1}(g(\varnothing))=\varnothing}{%
      \FormulaRefAuto{x\in A \vdash F^{-1}(F(x))=x}[thm]{3,
        10}}
    \proofstepwidestar[3,4]{g^{-1}(\varnothing)=\varnothing}{\rIE{4,11}}
    \proofstepwidestar[3]{g \colon \powerset ( D ) \to \powerset ( E )}{%
      \FormulaRefAuto{Bijektive Funktion}[def]{3}}
    \proofstepwidestar[3]{\LayerPowerMap{A}{D}{E}{g}\colon\powerset(A\cup D)\to\powerset(A\cup E)}{%
      \FormulaRefAuto{LayerPowerMapMapping}[thm]{%
          13}}
    \proofstepwidestar[3,6]{\LayerPowerMap{A}{D}{E}{g}(X)
      \in\powerset(A\cup E)}{%
      \FormulaRefAuto{F\colon A\to B\dsep x\in A\vdash F(x)\in B}[thm]{%
        14,6}}
    \proofstepwidestarbreakkeep[2,3,4,5,6,7]{%
      \LayerPowerMap{A}{E}{D}{g^{-1}}
      (\LayerPowerMap{A}{D}{E}{g}(X))=
      \LayerPowerMap{A}{D}{E}{g}(X)}{%
      \FormulaRefAuto{LayerPowerMapFixesStableSubsets}[thm]{%
        2,8,12,5,15,7}}
    \proofstepwidestar[3]{\forall K\in\powerset(D)\,g^{-1}(g(K))=K}{%
      \FormulaRefAuto{x\in A \vdash F^{-1}(F(x))=x}[thm]{3}}
    \proofstepwidestar[3]{g \colon \powerset ( D ) \to \powerset ( E )}{%
      \FormulaRefAuto{Bijektive Funktion}[def]{3}}
    \proofstepwidestar[1,2,3,6]{%
      \LayerPowerMap{A}{E}{D}{g^{-1}}
      (\LayerPowerMap{A}{D}{E}{g}(X))=X}{%
      \FormulaRefAuto{LayerPowerMapRoundTrip}[thm]{1,2,
        18,8,17,6}}
    \proofstepwidestar[1,2,3,4,5,6,7]{X=
      \LayerPowerMap{A}{D}{E}{g}(X)}{%
      \FormulaRefAuto{a=b\dsep a=c\vdash b=c}[thm]{19,16}}
    \proofstepwidestar[1,2,3,4,5,6,7]{X\subseteq L}{\rIE{20,7}}
  \closeproofpartwideindR

  \proofpartwide{Schluss}
    \proofstepwidestar[1]{A\cap D=\varnothing}{\rA}
    \proofstepwidestar[2]{A\cap E=\varnothing}{\rA}
    \proofstepwidestar[3]{g\colon\powerset(D)\bij\powerset(E)}{\rA}
    \proofstepwidestar[4]{g(\varnothing)=\varnothing}{\rA}
    \proofstepwidestar[5]{L\subseteq A}{\rA}
    \proofstepwidestar[6]{X\in\powerset(A\cup D)}{\rA}
    \proofstepwidestar[7]{X\subseteq L}{\rA}
    \proofstepwidestar[3]{g \colon \powerset ( D ) \to \powerset ( E )}{%
      \FormulaRefAuto{Bijektive Funktion}[def]{3}}
    \proofstepwidestar[1,3,4,5,6,7]{%
      \LayerPowerMap{A}{D}{E}{g}(X)\subseteq L}{%
      \FormulaRefAuto{LayerPowerMapPreservesSubsetsForward}[thm]{%
        1,8,4,5,6,7}}
    \proofstepwidestar[10]{\LayerPowerMap{A}{D}{E}{g}(X)\subseteq L}{\rA}
    \proofstepwidestar[1,2,3,4,5,6,10]{X\subseteq L}{%
      \FormulaRefAuto{LayerPowerMapPreservesSubsetsBackward}[thm]{%
        1,2,3,4,5,6,10}}
    \proofstepwidestarbreakkeep[1,2,3,4,5,6]{X\subseteq L\leftrightarrow
      \LayerPowerMap{A}{D}{E}{g}(X)\subseteq L}{%
      \rBI{\rRI{7,9},\rRI{10,11}}}
  \closeproofpartwide
\end{tabproofsplitwide}

\newpage
\FormulaThmDeltaK[Urbild nichtleerer Teilmengen unter einer nulltreuen Bijektion]{%
  F\colon\powerset(A)\bij\powerset(B)\dsep F(\varnothing)=\varnothing
  \vdash F^{-1}[\PowNE{B}]=\PowNE{A}
}{FixedEmptyBijectionNonemptyPreimage}{%
  \DeltaRow{Mengen}{A\dsep B\dsep X}
  \DeltaRow{Funktionen}{F}
}
\begin{tabproofsplitwide}
  \proofpartwideindR[Erste Elementrichtung]{%
    \begin{aligned}[t]
      &F\colon\powerset(A)\bij\powerset(B)\dsep F(\varnothing)=\varnothing
        \dsep X\in F^{-1}[\PowNE{B}]\vdash{}\\[-2pt]
      &X\in\PowNE{A}
    \end{aligned}%
  }{%
    F\colon\powerset(A)\bij\powerset(B)\dsep F(\varnothing)=\varnothing
    \dsep X\in F^{-1}[\PowNE{B}]\vdash X\in\PowNE{A}
  }[FixedEmptyBijectionNonemptyPreimageForward]
    \proofstepwidestar[1]{F\colon\powerset(A)\bij\powerset(B)}{\rA}
    \proofstepwidestar[2]{F(\varnothing)=\varnothing}{\rA}
    \proofstepwidestar[3]{X\in F^{-1}[\PowNE{B}]}{\rA}
    \proofstepwidestar[1]{F\colon\powerset(A)\to\powerset(B)}{%
      \FormulaRefAuto{Bijektive Funktion}[def]{1}}
    \proofstepwidestar[]{\PowNE{B}\subseteq\powerset(B)}{%
      \FormulaRefAuto{A\setminus B\subseteq A}[thm]}
    \proofstepwidestarbreakkeep[1]{X\in F^{-1}[\PowNE{B}]\leftrightarrow
      (X\in\powerset(A)\land F(X)\in\PowNE{B})}{%
      \FormulaRefAuto{F\colon A\to B\dsep C\subseteq B
        \vdash\bigl(x\in F^{-1}[C]\leftrightarrow
        (x\in A\land F(x)\in C)\bigr)}[thm]{4,5}}
    \proofstepwidestar[1,3]{X\in\powerset(A)\land F(X)\in\PowNE{B}}{%
      \FormulaRefAuto{P\leftrightarrow Q, P\vdash Q}[thm]{6,3}}
    \proofstepwidestar[1,3]{X\in\powerset(A)}{\rAEa{7}}
    \proofstepwidestar[1,3]{F(X)\subseteq B\land F(X)\neq\varnothing}{%
      \FormulaRefAuto{NonemptyPowersetMembership}[thm]{\rAEb{7}}}
    \proofstepwidestar[1,3]{F(X)\neq\varnothing}{%
      \rAEb{9}}
    \proofstepwidestar[11]{X=\varnothing}{\rA}
    \proofstepwidestar[2,11]{F(X)=\varnothing}{\rIE{11,2}}
    \proofstepwidestar[1,2,3,11]{\bot}{\rBI{10,12}}
    \proofstepwidestar[1,2,3]{X\neq\varnothing}{\rCI{11,13}}
    \proofstepwidestar[1,3]{X\subseteq A}{%
      \FormulaRefAuto{x\in\powerset(A)\eqvdash x\subseteq A}[thm]{8}}
    \proofstepwidestar[1,2,3]{X\in\PowNE{A}}{%
      \FormulaRefAuto{NonemptyPowersetMembership}[thm]{%
        \rAI{15,14}}}
  \closeproofpartwideindR

  \proofpartwideindR[Zweite Elementrichtung]{%
    \begin{aligned}[t]
      &F\colon\powerset(A)\bij\powerset(B)\dsep F(\varnothing)=\varnothing
        \dsep X\in\PowNE{A}\vdash{}\\[-2pt]
      &X\in F^{-1}[\PowNE{B}]
    \end{aligned}%
  }{%
    F\colon\powerset(A)\bij\powerset(B)\dsep F(\varnothing)=\varnothing
    \dsep X\in\PowNE{A}\vdash X\in F^{-1}[\PowNE{B}]
  }[FixedEmptyBijectionNonemptyPreimageBackward]
    \proofstepwidestar[1]{F\colon\powerset(A)\bij\powerset(B)}{\rA}
    \proofstepwidestar[2]{F(\varnothing)=\varnothing}{\rA}
    \proofstepwidestar[3]{X\in\PowNE{A}}{\rA}
    \proofstepwidestar[1]{F\colon\powerset(A)\to\powerset(B)}{%
      \FormulaRefAuto{Bijektive Funktion}[def]{1}}
    \proofstepwidestar[1]{F\colon\powerset(A)\inj\powerset(B)}{%
      \FormulaRefAuto{F\colon A\bij B\vdash F\colon A\inj B}[thm]{1}}
    \proofstepwidestar[3]{X\subseteq A\land X\neq\varnothing}{%
      \FormulaRefAuto{NonemptyPowersetMembership}[thm]{3}}
    \proofstepwidestar[3]{X\in\powerset(A)}{%
      \FormulaRefAuto{x\in\powerset(A)\eqvdash x\subseteq A}[thm]{\rAEa{6}}}
    \proofstepwidestar[]{\varnothing\subseteq A}{%
      \FormulaRefAuto{\varnothing\subseteq A}[thm]}
    \proofstepwidestar[]{\varnothing\in\powerset(A)}{%
      \FormulaRefAuto{x\in\powerset(A)\eqvdash x\subseteq A}[thm]{%
        8}}
    \proofstepwidestar[10]{F(X)=\varnothing}{\rA}
    \proofstepwidestar[2,10]{F(X)=F(\varnothing)}{\rIE{2,10}}
    \proofstepwidestar[1,2,3,10]{X=\varnothing}{%
      \FormulaRefAuto{Injektivität}[ax]{5,7,9,11}}
    \proofstepwidestar[1,2,3,10]{\bot}{\rBI{\rAEb{6},12}}
    \proofstepwidestar[1,2,3]{F(X)\neq\varnothing}{\rCI{10,13}}
    \proofstepwidestar[1,3]{F(X)\in\powerset(B)}{%
      \FormulaRefAuto{F\colon A\to B\dsep x\in A\vdash F(x)\in B}[thm]{4,7}}
    \proofstepwidestar[1,3]{F(X)\subseteq B}{%
      \FormulaRefAuto{x\in\powerset(A)\eqvdash x\subseteq A}[thm]{15}}
    \proofstepwidestar[1,2,3]{F(X)\in\PowNE{B}}{%
      \FormulaRefAuto{NonemptyPowersetMembership}[thm]{%
        \rAI{16,14}}}
    \proofstepwidestar[]{\PowNE{B}\subseteq\powerset(B)}{%
      \FormulaRefAuto{A\setminus B\subseteq A}[thm]}
    \proofstepwidestarbreakkeep[1]{X\in F^{-1}[\PowNE{B}]\leftrightarrow
      (X\in\powerset(A)\land F(X)\in\PowNE{B})}{%
      \FormulaRefAuto{F\colon A\to B\dsep C\subseteq B
        \vdash\bigl(x\in F^{-1}[C]\leftrightarrow
        (x\in A\land F(x)\in C)\bigr)}[thm]{4,18}}
    \proofstepwidestar[1,2,3]{X\in F^{-1}[\PowNE{B}]}{%
      \FormulaRefAuto{P\leftrightarrow Q, Q\vdash P}[thm]{19,\rAI{7,17}}}
  \closeproofpartwideindR

  \proofpartwide{Schluss}
    \proofstepwidestar[1]{F\colon\powerset(A)\bij\powerset(B)}{\rA}
    \proofstepwidestar[2]{F(\varnothing)=\varnothing}{\rA}
    \proofstepwidestar[3]{X\in F^{-1}[\PowNE{B}]}{\rA}
    \proofstepwidestar[1,2,3]{X\in\PowNE{A}}{%
      \FormulaRefAuto{FixedEmptyBijectionNonemptyPreimageForward}[thm]{1,2,3}}
    \proofstepwidestarbreakkeep[1,2]{F^{-1}[\PowNE{B}]\subseteq\PowNE{A}}{%
      \FormulaRefAuto{A\subseteq B\coloneq
        \forall x\,(x\in A\rightarrow x\in B)}[def]{\rUI{\rRI{3,4}}}}
    \proofstepwidestar[6]{X\in\PowNE{A}}{\rA}
    \proofstepwidestar[1,2,6]{X\in F^{-1}[\PowNE{B}]}{%
      \FormulaRefAuto{FixedEmptyBijectionNonemptyPreimageBackward}[thm]{1,2,6}}
    \proofstepwidestarbreakkeep[1,2]{\PowNE{A}\subseteq F^{-1}[\PowNE{B}]}{%
      \FormulaRefAuto{A\subseteq B\coloneq
        \forall x\,(x\in A\rightarrow x\in B)}[def]{\rUI{\rRI{6,7}}}}
    \proofstepwidestarbreakkeep[1,2]{F^{-1}[\PowNE{B}]=\PowNE{A}}{%
      \FormulaRefAuto{A\subseteq B\dsep B\subseteq A\vdash A=B}[thm]{5,8}}
  \closeproofpartwide
\end{tabproofsplitwide}

\FormulaThmDeltaK[Bijektive Fortsetzung auf nichtleere Teilmengen]{%
  \begin{aligned}[t]
    &A\cap D=\varnothing\dsep A\cap E=\varnothing\dsep
      g\colon\powerset(D)\bij\powerset(E)\dsep
      g(\varnothing)=\varnothing\vdash{}\\[-2pt]
    &\LayerPowerMap{A}{D}{E}{g}\restriction_{\PowNE{A\cup D}}
      \colon\PowNE{A\cup D}\bij\PowNE{A\cup E}
  \end{aligned}%
}{LayerPowerMapNonemptyBijective}{%
  \DeltaRow{Mengen}{A\dsep D\dsep E}
  \DeltaRow{Funktionen}{g\dsep\LayerPowerMap{A}{D}{E}{g}}
}
\begin{tabproofwide}
  \proofstepwidestar[1]{A\cap D=\varnothing}{\rA}
  \proofstepwidestar[2]{A\cap E=\varnothing}{\rA}
  \proofstepwidestar[3]{g\colon\powerset(D)\bij\powerset(E)}{\rA}
  \proofstepwidestar[4]{g(\varnothing)=\varnothing}{\rA}
  \proofstepwidestar[1,2,3]{\LayerPowerMap{A}{D}{E}{g}\colon
    \powerset(A\cup D)\bij\powerset(A\cup E)}{%
    \FormulaRefAuto{LayerPowerMapBijective}[thm]{1,2,3}}
  \proofstepwidestar[]{\varnothing\subseteq A\cup D}{%
      \FormulaRefAuto{\varnothing\subseteq A}[thm]}
  \proofstepwidestar[]{\varnothing\in\powerset(A\cup D)}{%
    \FormulaRefAuto{x\in\powerset(A)\eqvdash x\subseteq A}[thm]{%
      6}}
  \proofstepwidestar[3]{g \colon \powerset ( D ) \to \powerset ( E )}{%
      \FormulaRefAuto{Bijektive Funktion}[def]{3}}
  \proofstepwidestar[3,7]{\LayerPowerMap{A}{D}{E}{g}(\varnothing)=
    (\varnothing\cap A)\cup g(\varnothing\cap D)}{%
    \FormulaRefAuto{LayerPowerMapEvaluation}[thm]{%
      8,7}}
  \proofstepwidestar[]{\varnothing\cup\varnothing=\varnothing}{%
      \FormulaRefAuto{\varnothing\cup A=A}[thm]}
  \proofstepwidestar[]{\varnothing\cap A=\varnothing}{%
      \FormulaRefAuto{\varnothing\cap A=\varnothing}[thm]}
  \proofstepwidestar[]{\varnothing\cap D=\varnothing}{%
      \FormulaRefAuto{\varnothing\cap A=\varnothing}[thm]}
  \proofstepwidestarbreakkeep[3,4]{\LayerPowerMap{A}{D}{E}{g}(\varnothing)=\varnothing}{%
    \rChain{9,\rLRS{11},
      \rLRS{12},
      \rLRS{4},10}}
  \proofstepwidestar[1,2,3,4]{%
    (\LayerPowerMap{A}{D}{E}{g})^{-1}[\PowNE{A\cup E}]
    =\PowNE{A\cup D}}{%
    \FormulaRefAuto{FixedEmptyBijectionNonemptyPreimage}[thm]{5,13}}
  \proofstepwidestar[]{\PowNE{A\cup E}\subseteq\powerset(A\cup E)}{%
    \FormulaRefAuto{A\setminus B\subseteq A}[thm]}
  \proofstepwidestarbreakkeep[1,2,3]{%
    \begin{aligned}[t]
      &\LayerPowerMap{A}{D}{E}{g}\restriction_{
        (\LayerPowerMap{A}{D}{E}{g})^{-1}[\PowNE{A\cup E}]}\\[-2pt]
      &\quad\colon(\LayerPowerMap{A}{D}{E}{g})^{-1}[\PowNE{A\cup E}]\\[-2pt]
      &\quad\bij\PowNE{A\cup E}
    \end{aligned}}{%
    \FormulaRefAuto{F\colon A\bij B \dsep T\subseteq B \vdash
      F\restriction_{F^{-1}[T]}\colon F^{-1}[T]\bij T}[thm]{5,15}}
  \proofstepwidestar[1,2,3,4]{%
    \LayerPowerMap{A}{D}{E}{g}\restriction_{\PowNE{A\cup D}}
    \colon\PowNE{A\cup D}\bij\PowNE{A\cup E}}{\rIE{14,16}}
\end{tabproofwide}


\section{Erweiterung einer Bijektion}

\FormulaThmDeltaK[Frische Erweiterung stimmt auf dem alten Bereich überein]{%
  \begin{aligned}[t]
    &B\subseteq C\dsep F\colon A\to B\dsep a\notin A\dsep c\in C\dsep x\in A\vdash{}\\
    &\ExtMap{\iota_{B,C}\circ F}{a}{c}(x)=F(x)
  \end{aligned}%
}{FreshExtMapOnBase}{%
  \DeltaRow{Mengen}{A\dsep B\dsep C\dsep a\dsep c\dsep x}%
  \DeltaRow{Funktionen}{F}%
}
\begin{tabproof}
  \proofstep{1}{B\subseteq C}{\rA}
  \proofstep{2}{F\colon A\to B}{\rA}
  \proofstep{3}{a\notin A}{\rA}
  \proofstep{4}{c\in C}{\rA}
  \proofstep{5}{x\in A}{\rA}
  \proofstep{1,2}{\iota_{B,C}\circ F\colon A\to C}{%
    \FormulaRefAuto{B\subseteq C \dsep F\colon A\to B \vdash \iota_{B,C}\circ F\colon A\to C}{1,2}}
  \proofstep{1,2,3,4}{%
    \forall x\in A\,\ExtMap{\iota_{B,C}\circ F}{a}{c}(x)=(\iota_{B,C}\circ F)(x)}{%
    \FormulaRefAuto{a\notin A \dsep f\colon A\to M \dsep b\in M \vdash
      \forall x\in A\,\ExtMap{f}{a}{b}(x)=f(x)}{3,6,4}}
  \proofstep{1,2,3,4,5}{%
    \ExtMap{\iota_{B,C}\circ F}{a}{c}(x)=(\iota_{B,C}\circ F)(x)}{%
    \rRE{\rUE{7},5}}
  \proofstep{1,2,5}{(\iota_{B,C}\circ F)(x)=F(x)}{%
    \FormulaRefAuto{InclusionCompositionValue}[thm]{1,2,5}}
  \proofstep{1,2,3,4,5}{\ExtMap{\iota_{B,C}\circ F}{a}{c}(x)=F(x)}{%
    \FormulaRefAuto{a=b,\,b=c\vdash a=c}{8,9}}
\end{tabproof}

\FormulaThmDeltaK[Frische Erweiterung trifft den neuen Wert]{%
  B\subseteq C\dsep F\colon A\to B\dsep a\notin A\dsep c\in C
  \vdash \ExtMap{\iota_{B,C}\circ F}{a}{c}(a)=c%
}{FreshExtMapAtNew}{%
  \DeltaRow{Mengen}{A\dsep B\dsep C\dsep a\dsep c}%
  \DeltaRow{Funktionen}{F}%
}
\begin{tabproof}
  \proofstep{1}{B\subseteq C}{\rA}
  \proofstep{2}{F\colon A\to B}{\rA}
  \proofstep{3}{a\notin A}{\rA}
  \proofstep{4}{c\in C}{\rA}
  \proofstep{1,2}{\iota_{B,C}\circ F\colon A\to C}{%
    \FormulaRefAuto{B\subseteq C \dsep F\colon A\to B \vdash \iota_{B,C}\circ F\colon A\to C}{1,2}}
  \proofstep{1,2,3,4}{\ExtMap{\iota_{B,C}\circ F}{a}{c}(a)=c}{%
    \FormulaRefAuto{a\notin A \dsep f\colon A\to M \dsep b\in M \vdash
      \ExtMap{f}{a}{b}(a)=b}{3,5,4}}
\end{tabproof}



\FormulaThmDeltaK[Frische Erweiterung einer Bijektion]{%
  \begin{aligned}[t]
    &F\colon A\bij B\dsep a\notin A\dsep b\notin B\vdash{}\\
    &\ExtMap{\iota_{B,B\cup\{b\}}\circ F}{a}{b}
      \colon A\cup\{a\}\bij B\cup\{b\}
  \end{aligned}%
}{FreshExtMapBijective}{%
  \DeltaRow{Mengen}{A\dsep B\dsep a\dsep b\dsep x\dsep y\dsep z}%
  \DeltaRow{Funktionen}{F}%
}
\begin{tabproofsplitwide}
  \proofpartwideindR[Alt--Alt-Fall der Injektivität]{%
    \begin{aligned}[t]
      &F\colon A\bij B\dsep a\notin A\dsep b\notin B\dsep x\in A\dsep y\in A\dsep{}\\
      &\ExtMap{\iota_{B,B\cup\{b\}}\circ F}{a}{b}(x)
       =\ExtMap{\iota_{B,B\cup\{b\}}\circ F}{a}{b}(y)\vdash x=y
    \end{aligned}%
  }{%
    F\colon A\bij B\dsep a\notin A\dsep b\notin B\dsep x\in A\dsep y\in A\dsep
    \ExtMap{\iota_{B,B\cup\{b\}}\circ F}{a}{b}(x)
    =\ExtMap{\iota_{B,B\cup\{b\}}\circ F}{a}{b}(y)\vdash x=y
  }[FreshExtMapInjectiveCaseAA]
    \proofstepwidestar[1]{F\colon A\bij B}{\rA}
    \proofstepwidestar[2]{a\notin A}{\rA}
    \proofstepwidestar[3]{b\notin B}{\rA}
    \proofstepwidestar[4]{x\in A}{\rA}
    \proofstepwidestar[5]{y\in A}{\rA}
    \proofstepwidestar[6]{%
      \begin{aligned}[t]
        &\ExtMap{\iota_{B,B\cup\{b\}}\circ F}{a}{b}(x)\\[-2pt]
        &\qquad=\ExtMap{\iota_{B,B\cup\{b\}}\circ F}{a}{b}(y)
      \end{aligned}}{\rA}
    \proofstepwidestar[]{B\subseteq B\cup\{b\}}{%
      \FormulaRefAuto{A\subseteq A\cup B}}
    \proofstepwidestar[]{b\in B\cup\{b\}}{\FormulaRefAuto{a\in A\cup\{a\}}}
    \proofstepwidestar[1]{F\colon A\to B}{\FormulaRefAuto{Bijektive Funktion}{1}}
    \proofstepwidestar[1,2,4]{%
      \ExtMap{\iota_{B,B\cup\{b\}}\circ F}{a}{b}(x)=F(x)}{%
      \FormulaRefAuto{FreshExtMapOnBase}[thm]{7,9,2,8,4}}
    \proofstepwidestar[1,2,5]{%
      \ExtMap{\iota_{B,B\cup\{b\}}\circ F}{a}{b}(y)=F(y)}{%
      \FormulaRefAuto{FreshExtMapOnBase}[thm]{7,9,2,8,5}}
    \proofstepwidestar[1,2,4,5,6]{F(x)=F(y)}{%
      \FormulaRefAuto{a = b\dsep a = c\dsep b=d\vdash c = d}{6,10,11}}
    \proofstepwidestar[1]{F\colon A\inj B}{%
      \FormulaRefAuto{F\colon A\bij B \vdash F\colon A\inj B}{1}}
    \proofstepwidestar[1,2,4,5,6]{x=y}{%
      \FormulaRefAuto{F\colon A\inj B\dsep x\in A\dsep y\in A\dsep
        F(x)=F(y)\vdash x=y}{13,4,5,12}}
  \closeproofpartwideindR

  \proofpartwideindR[Alt--Neu-Fall der Injektivität]{%
    \begin{aligned}[t]
      &F\colon A\bij B\dsep a\notin A\dsep b\notin B\dsep x\in A\dsep y=a\dsep{}\\
      &\ExtMap{\iota_{B,B\cup\{b\}}\circ F}{a}{b}(x)
       =\ExtMap{\iota_{B,B\cup\{b\}}\circ F}{a}{b}(y)\vdash x=y
    \end{aligned}%
  }{%
    F\colon A\bij B\dsep a\notin A\dsep b\notin B\dsep x\in A\dsep y=a\dsep
    \ExtMap{\iota_{B,B\cup\{b\}}\circ F}{a}{b}(x)
    =\ExtMap{\iota_{B,B\cup\{b\}}\circ F}{a}{b}(y)\vdash x=y
  }[FreshExtMapInjectiveCaseAN]
    \proofstepwidestar[1]{F\colon A\bij B}{\rA}
    \proofstepwidestar[2]{a\notin A}{\rA}
    \proofstepwidestar[3]{b\notin B}{\rA}
    \proofstepwidestar[4]{x\in A}{\rA}
    \proofstepwidestar[5]{y=a}{\rA}
    \proofstepwidestar[6]{%
      \begin{aligned}[t]
        &\ExtMap{\iota_{B,B\cup\{b\}}\circ F}{a}{b}(x)\\[-2pt]
        &\qquad=\ExtMap{\iota_{B,B\cup\{b\}}\circ F}{a}{b}(y)
      \end{aligned}}{\rA}
    \proofstepwidestar[]{B\subseteq B\cup\{b\}}{\FormulaRefAuto{A\subseteq A\cup B}}
    \proofstepwidestar[]{b\in B\cup\{b\}}{\FormulaRefAuto{a\in A\cup\{a\}}}
    \proofstepwidestar[1]{F\colon A\to B}{\FormulaRefAuto{Bijektive Funktion}{1}}
    \proofstepwidestar[1,2,4]{%
      \ExtMap{\iota_{B,B\cup\{b\}}\circ F}{a}{b}(x)=F(x)}{%
      \FormulaRefAuto{FreshExtMapOnBase}[thm]{7,9,2,8,4}}
    \proofstepwidestar[1,2]{%
      \ExtMap{\iota_{B,B\cup\{b\}}\circ F}{a}{b}(a)=b}{%
      \FormulaRefAuto{FreshExtMapAtNew}[thm]{7,9,2,8}}
    \proofstepwidestar[1,2,5]{%
      \ExtMap{\iota_{B,B\cup\{b\}}\circ F}{a}{b}(y)=b}{\rIE{5,11}}
    \proofstepwidestar[1,2,4,5,6]{F(x)=b}{%
      \FormulaRefAuto{a = b\dsep a = c\dsep b=d\vdash c = d}{6,10,12}}
    \proofstepwidestar[1,4]{F(x)\in B}{%
      \FormulaRefAuto{F\colon A\bij B\dsep x\in A\vdash F(x)\in B}{1,4}}
    \proofstepwidestar[1,2,4,5,6]{b = F ( x )}{%
      \FormulaRefAuto{a=b\vdash b=a}{13}}
    \proofstepwidestar[1,2,3,4,5,6]{b\in B}{%
      \rIE{15,14}}
    \proofstepwidestar[1,2,3,4,5,6]{\bot}{\rBI{16,3}}
    \proofstepwidestar[1,2,3,4,5,6]{x=y}{\rCE{17}}
  \closeproofpartwideindR

  \proofpartwideindR[Alt-Fall der Injektivität]{%
    \begin{aligned}[t]
      &F\colon A\bij B\dsep a\notin A\dsep b\notin B\dsep
      x\in A\dsep y\in A\cup\{a\}\dsep{}\\
      &\ExtMap{\iota_{B,B\cup\{b\}}\circ F}{a}{b}(x)
       =\ExtMap{\iota_{B,B\cup\{b\}}\circ F}{a}{b}(y)\vdash x=y
    \end{aligned}%
  }{%
    F\colon A\bij B\dsep a\notin A\dsep b\notin B\dsep
    x\in A\dsep y\in A\cup\{a\}\dsep
    \ExtMap{\iota_{B,B\cup\{b\}}\circ F}{a}{b}(x)
    =\ExtMap{\iota_{B,B\cup\{b\}}\circ F}{a}{b}(y)\vdash x=y
  }[FreshExtMapInjectiveCaseA]
    \proofstepwidestar[1]{F\colon A\bij B}{\rA}
    \proofstepwidestar[2]{a\notin A}{\rA}
    \proofstepwidestar[3]{b\notin B}{\rA}
    \proofstepwidestar[4]{x\in A}{\rA}
    \proofstepwidestar[5]{y\in A\cup\{a\}}{\rA}
    \proofstepwidestar[6]{%
      \begin{aligned}[t]
        &\ExtMap{\iota_{B,B\cup\{b\}}\circ F}{a}{b}(x)\\[-2pt]
        &\qquad=\ExtMap{\iota_{B,B\cup\{b\}}\circ F}{a}{b}(y)
      \end{aligned}}{\rA}
    \proofstepwidestar[5]{y\in A\lor y=a}{%
      \FormulaRefAuto{x\in A\cup\{a\} \vdash x\in A\lor x=a}{5}}
    \proofstepwidestar[8]{y\in A}{\rA}
    \proofstepwidestar[1,2,3,4,6,8]{x=y}{%
      \FormulaRefAuto{FreshExtMapInjectiveCaseAA}{1,2,3,4,8,6}}
    \proofstepwidestar[10]{y=a}{\rA}
    \proofstepwidestar[1,2,3,4,6,10]{x=y}{%
      \FormulaRefAuto{FreshExtMapInjectiveCaseAN}{1,2,3,4,10,6}}
    \proofstepwidestar[1,2,3,4,5,6]{x=y}{\rOE{7,8,9,10,11}}
  \closeproofpartwideindR

  \proofpartwideindR[Neu-Fall der Injektivität]{%
    \begin{aligned}[t]
      &F\colon A\bij B\dsep a\notin A\dsep b\notin B\dsep
      x=a\dsep y\in A\cup\{a\}\dsep{}\\
      &\ExtMap{\iota_{B,B\cup\{b\}}\circ F}{a}{b}(x)
       =\ExtMap{\iota_{B,B\cup\{b\}}\circ F}{a}{b}(y)\vdash x=y
    \end{aligned}%
  }{%
    F\colon A\bij B\dsep a\notin A\dsep b\notin B\dsep
    x=a\dsep y\in A\cup\{a\}\dsep
    \ExtMap{\iota_{B,B\cup\{b\}}\circ F}{a}{b}(x)
    =\ExtMap{\iota_{B,B\cup\{b\}}\circ F}{a}{b}(y)\vdash x=y
  }[FreshExtMapInjectiveCaseN]
    \proofstepwidestar[1]{F\colon A\bij B}{\rA}
    \proofstepwidestar[2]{a\notin A}{\rA}
    \proofstepwidestar[3]{b\notin B}{\rA}
    \proofstepwidestar[4]{x=a}{\rA}
    \proofstepwidestar[5]{y\in A\cup\{a\}}{\rA}
    \proofstepwidestar[6]{%
      \ExtMap{\iota_{B,B\cup\{b\}}\circ F}{a}{b}(x)
      =\ExtMap{\iota_{B,B\cup\{b\}}\circ F}{a}{b}(y)}{\rA}
    \proofstepwidestar[5]{y\in A\lor y=a}{%
      \FormulaRefAuto{x\in A\cup\{a\} \vdash x\in A\lor x=a}{5}}
    \proofstepwidestar[8]{y\in A}{\rA}
    \proofstepwidestar[6]{%
      \begin{aligned}[t]
        &\ExtMap{\iota_{B,B\cup\{b\}}\circ F}{a}{b}(y)\\[-2pt]
        &\qquad=\ExtMap{\iota_{B,B\cup\{b\}}\circ F}{a}{b}(x)
      \end{aligned}}{\FormulaRefAuto{a=b\vdash b=a}{6}}
    \proofstepwidestar[1,2,3,4,6,8]{y=x}{%
      \FormulaRefAuto{FreshExtMapInjectiveCaseAN}{1,2,3,8,4,9}}
    \proofstepwidestar[1,2,3,4,6,8]{x=y}{%
      \FormulaRefAuto{a=b\vdash b=a}{10}}
    \proofstepwidestar[12]{y=a}{\rA}
    \proofstepwidestar[4,12]{x=y}{%
      \FormulaRefAuto{a = b, c = b\vdash a = c}{4,12}}
    \proofstepwidestar[1,2,3,4,5,6]{x=y}{\rOE{7,8,11,12,13}}
  \closeproofpartwideindR

  \proofpartwideindR[Punktweise Injektivität]{%
    \begin{aligned}[t]
      &F\colon A\bij B\dsep a\notin A\dsep b\notin B\dsep
      x\in A\cup\{a\}\dsep y\in A\cup\{a\}\dsep{}\\
      &\ExtMap{\iota_{B,B\cup\{b\}}\circ F}{a}{b}(x)
       =\ExtMap{\iota_{B,B\cup\{b\}}\circ F}{a}{b}(y)\vdash x=y
    \end{aligned}%
  }{%
    F\colon A\bij B\dsep a\notin A\dsep b\notin B\dsep
    x\in A\cup\{a\}\dsep y\in A\cup\{a\}\dsep
    \ExtMap{\iota_{B,B\cup\{b\}}\circ F}{a}{b}(x)
    =\ExtMap{\iota_{B,B\cup\{b\}}\circ F}{a}{b}(y)\vdash x=y
  }[FreshExtMapInjectivePointwise]
    \proofstepwidestar[1]{F\colon A\bij B}{\rA}
    \proofstepwidestar[2]{a\notin A}{\rA}
    \proofstepwidestar[3]{b\notin B}{\rA}
    \proofstepwidestar[4]{x\in A\cup\{a\}}{\rA}
    \proofstepwidestar[5]{y\in A\cup\{a\}}{\rA}
    \proofstepwidestar[6]{%
      \ExtMap{\iota_{B,B\cup\{b\}}\circ F}{a}{b}(x)
      =\ExtMap{\iota_{B,B\cup\{b\}}\circ F}{a}{b}(y)}{\rA}
    \proofstepwidestar[4]{x\in A\lor x=a}{%
      \FormulaRefAuto{x\in A\cup\{a\} \vdash x\in A\lor x=a}{4}}
    \proofstepwidestar[8]{x\in A}{\rA}
    \proofstepwidestar[1,2,3,5,6,8]{x=y}{%
      \FormulaRefAuto{FreshExtMapInjectiveCaseA}{1,2,3,8,5,6}}
    \proofstepwidestar[10]{x=a}{\rA}
    \proofstepwidestar[1,2,3,5,6,10]{x=y}{%
      \FormulaRefAuto{FreshExtMapInjectiveCaseN}{1,2,3,10,5,6}}
    \proofstepwidestar[1,2,3,4,5,6]{x=y}{\rOE{7,8,9,10,11}}
  \closeproofpartwideindR

  \proofpartwideindR[Injektivität der frischen Erweiterung]{%
    F\colon A\bij B\dsep a\notin A\dsep b\notin B\vdash
    \ExtMap{\iota_{B,B\cup\{b\}}\circ F}{a}{b}
    \colon A\cup\{a\}\inj B\cup\{b\}%
  }{%
    F\colon A\bij B\dsep a\notin A\dsep b\notin B\vdash
    \ExtMap{\iota_{B,B\cup\{b\}}\circ F}{a}{b}
    \colon A\cup\{a\}\inj B\cup\{b\}%
  }[FreshExtMapInjective]
    \proofstepwidestar[1]{F\colon A\bij B}{\rA}
    \proofstepwidestar[2]{a\notin A}{\rA}
    \proofstepwidestar[3]{b\notin B}{\rA}
    \proofstepwidestar[]{B\subseteq B\cup\{b\}}{\FormulaRefAuto{A\subseteq A\cup B}}
    \proofstepwidestar[]{b\in B\cup\{b\}}{\FormulaRefAuto{a\in A\cup\{a\}}}
    \proofstepwidestar[1]{F\colon A\to B}{\FormulaRefAuto{Bijektive Funktion}{1}}
    \proofstepwidestar[1]{\iota_{B,B\cup\{b\}}\circ F\colon A\to B\cup\{b\}}{%
      \FormulaRefAuto{B\subseteq C \dsep F\colon A\to B \vdash \iota_{B,C}\circ F\colon A\to C}{4,6}}
    \proofstepwidestar[1,2]{%
      \begin{aligned}[t]
        &\ExtMap{\iota_{B,B\cup\{b\}}\circ F}{a}{b}\\[-2pt]
        &\qquad\colon A\cup\{a\}\to B\cup\{b\}
      \end{aligned}}{%
      \FormulaRefAuto{a\notin A \dsep f\colon A\to M \dsep b\in M \vdash
        \ExtMap{f}{a}{b}\colon (A\cup\{a\})\to M}{2,7,5}}
    \proofstepwidestar[9]{x\in A\cup\{a\}}{\rA}
    \proofstepwidestar[10]{y\in A\cup\{a\}}{\rA}
    \proofstepwidestar[11]{%
      \begin{aligned}[t]
        &\ExtMap{\iota_{B,B\cup\{b\}}\circ F}{a}{b}(x)\\[-2pt]
        &\qquad=\ExtMap{\iota_{B,B\cup\{b\}}\circ F}{a}{b}(y)
      \end{aligned}}{\rA}
    \proofstepwidestar[1,2,3,9,10,11]{x=y}{%
      \FormulaRefAuto{FreshExtMapInjectivePointwise}{1,2,3,9,10,11}}
    \proofstepwidestar[1,2,3,9,10]{%
      \begin{aligned}[t]
        &\ExtMap{\iota_{B,B\cup\{b\}}\circ F}{a}{b}(x)
        =\ExtMap{\iota_{B,B\cup\{b\}}\circ F}{a}{b}(y)\\[-2pt]
        &\qquad\rightarrow x=y
      \end{aligned}}{\rRI{11,12}}
    \proofstepwidestar[1,2,3,9]{%
      \begin{aligned}[t]
        &\forall y\in A\cup\{a\}\,\Bigl(\\[-2pt]
        &\quad\ExtMap{\iota_{B,B\cup\{b\}}\circ F}{a}{b}(x)\\[-2pt]
        &\qquad=\ExtMap{\iota_{B,B\cup\{b\}}\circ F}{a}{b}(y)
          \rightarrow x=y\Bigr)
      \end{aligned}}{\rUI{\rRI{10,13}}}
    \proofstepwidestar[1,2,3]{%
      \begin{aligned}[t]
        &\forall x,y\in A\cup\{a\}\,\Bigl(\\[-2pt]
        &\quad\ExtMap{\iota_{B,B\cup\{b\}}\circ F}{a}{b}(x)\\[-2pt]
        &\qquad=\ExtMap{\iota_{B,B\cup\{b\}}\circ F}{a}{b}(y)\\[-2pt]
        &\hspace{8em}\rightarrow x=y\Bigr)
      \end{aligned}}{\rUI{\rRI{9,14}}}
    \proofstepwidestar[1,2,3]{%
      \begin{aligned}[t]
        &\ExtMap{\iota_{B,B\cup\{b\}}\circ F}{a}{b}\\[-2pt]
        &\qquad\colon A\cup\{a\}\inj B\cup\{b\}
      \end{aligned}}{%
      \FormulaRefAuto{Injektive Funktion}{8,15}}
  \closeproofpartwideindR

  \proofpartwideindR[Alter Zielwert der Surjektivität]{%
    \begin{aligned}[t]
      &F\colon A\bij B\dsep a\notin A\dsep b\notin B\dsep z\in B\vdash{}\\
      &\exists x\in A\cup\{a\}\,
        \ExtMap{\iota_{B,B\cup\{b\}}\circ F}{a}{b}(x)=z
    \end{aligned}%
  }{%
    F\colon A\bij B\dsep a\notin A\dsep b\notin B\dsep z\in B\vdash
    \exists x\in A\cup\{a\}\,
    \ExtMap{\iota_{B,B\cup\{b\}}\circ F}{a}{b}(x)=z
  }[FreshExtMapSurjectiveOldValue]
    \proofstepwidestar[1]{F\colon A\bij B}{\rA}
    \proofstepwidestar[2]{a\notin A}{\rA}
    \proofstepwidestar[3]{b\notin B}{\rA}
    \proofstepwidestar[4]{z\in B}{\rA}
    \proofstepwidestar[]{B\subseteq B\cup\{b\}}{\FormulaRefAuto{A\subseteq A\cup B}}
    \proofstepwidestar[]{b\in B\cup\{b\}}{\FormulaRefAuto{a\in A\cup\{a\}}}
    \proofstepwidestar[1]{F\colon A\to B}{\FormulaRefAuto{Bijektive Funktion}{1}}
    \proofstepwidestar[1]{F\colon A\sur B}{\FormulaRefAuto{Bijektive Funktion}{1}}
    \proofstepwidestar[1,4]{\exists x\in A\,F(x)=z}{%
      \FormulaRefAuto{y\in B\vdash \exists x\in A\; F(x)=y}{4}}
    \proofstepwidestar[10]{x\in A\land F(x)=z}{\rA}
    \proofstepwidestar[10]{x\in A}{\rAEa{10}}
    \proofstepwidestar[10]{F(x)=z}{\rAEb{10}}
    \proofstepwidestar[10]{x\in A\cup\{a\}}{%
      \FormulaRefAuto{z\in A\vdash z\in A\cup B}{11}}
    \proofstepwidestar[1,2,10]{%
      \ExtMap{\iota_{B,B\cup\{b\}}\circ F}{a}{b}(x)=F(x)}{%
      \FormulaRefAuto{FreshExtMapOnBase}[thm]{5,7,2,6,11}}
    \proofstepwidestar[1,2,10]{%
      \ExtMap{\iota_{B,B\cup\{b\}}\circ F}{a}{b}(x)=z}{%
      \FormulaRefAuto{a=b,\,b=c\vdash a=c}{14,12}}
    \proofstepwidestar[1,2,10]{%
      \exists x\in A\cup\{a\}\,
      \ExtMap{\iota_{B,B\cup\{b\}}\circ F}{a}{b}(x)=z}{%
      \rEI{\rAI{13,15}}}
    \proofstepwidestar[1,2,4]{%
      \exists x\in A\cup\{a\}\,
      \ExtMap{\iota_{B,B\cup\{b\}}\circ F}{a}{b}(x)=z}{\rEE{9,10,16}}
  \closeproofpartwideindR

  \proofpartwideindR[Neuer Zielwert der Surjektivität]{%
    \begin{aligned}[t]
      &F\colon A\bij B\dsep a\notin A\dsep b\notin B\dsep z=b\vdash{}\\
      &\exists x\in A\cup\{a\}\,
        \ExtMap{\iota_{B,B\cup\{b\}}\circ F}{a}{b}(x)=z
    \end{aligned}%
  }{%
    F\colon A\bij B\dsep a\notin A\dsep b\notin B\dsep z=b\vdash
    \exists x\in A\cup\{a\}\,
    \ExtMap{\iota_{B,B\cup\{b\}}\circ F}{a}{b}(x)=z
  }[FreshExtMapSurjectiveNewValue]
    \proofstepwidestar[1]{F\colon A\bij B}{\rA}
    \proofstepwidestar[2]{a\notin A}{\rA}
    \proofstepwidestar[3]{b\notin B}{\rA}
    \proofstepwidestar[4]{z=b}{\rA}
    \proofstepwidestar[]{B\subseteq B\cup\{b\}}{\FormulaRefAuto{A\subseteq A\cup B}}
    \proofstepwidestar[]{b\in B\cup\{b\}}{\FormulaRefAuto{a\in A\cup\{a\}}}
    \proofstepwidestar[1]{F\colon A\to B}{\FormulaRefAuto{Bijektive Funktion}{1}}
    \proofstepwidestar[]{a\in A\cup\{a\}}{\FormulaRefAuto{a\in A\cup\{a\}}}
    \proofstepwidestar[1,2]{%
      \ExtMap{\iota_{B,B\cup\{b\}}\circ F}{a}{b}(a)=b}{%
      \FormulaRefAuto{FreshExtMapAtNew}[thm]{5,7,2,6}}
    \proofstepwidestar[4]{b = z}{%
      \FormulaRefAuto{a=b\vdash b=a}{4}}
    \proofstepwidestar[1,2,4]{%
      \ExtMap{\iota_{B,B\cup\{b\}}\circ F}{a}{b}(a)=z}{%
      \rIE{10,9}}
    \proofstepwidestar[1,2,4]{%
      \exists x\in A\cup\{a\}\,
      \ExtMap{\iota_{B,B\cup\{b\}}\circ F}{a}{b}(x)=z}{%
      \rEI{\rAI{8,11}}}
  \closeproofpartwideindR

  \proofpartwideindR[Surjektivität der frischen Erweiterung]{%
    F\colon A\bij B\dsep a\notin A\dsep b\notin B\vdash
    \ExtMap{\iota_{B,B\cup\{b\}}\circ F}{a}{b}
    \colon A\cup\{a\}\sur B\cup\{b\}%
  }{%
    F\colon A\bij B\dsep a\notin A\dsep b\notin B\vdash
    \ExtMap{\iota_{B,B\cup\{b\}}\circ F}{a}{b}
    \colon A\cup\{a\}\sur B\cup\{b\}%
  }[FreshExtMapSurjective]
    \proofstepwidestar[1]{F\colon A\bij B}{\rA}
    \proofstepwidestar[2]{a\notin A}{\rA}
    \proofstepwidestar[3]{b\notin B}{\rA}
    \proofstepwidestar[]{B\subseteq B\cup\{b\}}{\FormulaRefAuto{A\subseteq A\cup B}}
    \proofstepwidestar[]{b\in B\cup\{b\}}{\FormulaRefAuto{a\in A\cup\{a\}}}
    \proofstepwidestar[1]{F\colon A\to B}{\FormulaRefAuto{Bijektive Funktion}{1}}
    \proofstepwidestar[1]{\iota_{B,B\cup\{b\}}\circ F\colon A\to B\cup\{b\}}{%
      \FormulaRefAuto{B\subseteq C \dsep F\colon A\to B \vdash \iota_{B,C}\circ F\colon A\to C}{4,6}}
    \proofstepwidestar[1,2]{%
      \ExtMap{\iota_{B,B\cup\{b\}}\circ F}{a}{b}
      \colon A\cup\{a\}\to B\cup\{b\}}{%
      \FormulaRefAuto{a\notin A \dsep f\colon A\to M \dsep b\in M \vdash
        \ExtMap{f}{a}{b}\colon (A\cup\{a\})\to M}{2,7,5}}
    \proofstepwidestar[9]{z\in B\cup\{b\}}{\rA}
    \proofstepwidestar[9]{z\in B\lor z=b}{%
      \FormulaRefAuto{x\in A\cup\{a\} \vdash x\in A\lor x=a}{9}}
    \proofstepwidestar[11]{z\in B}{\rA}
    \proofstepwidestar[1,2,3,11]{%
      \exists x\in A\cup\{a\}\,
      \ExtMap{\iota_{B,B\cup\{b\}}\circ F}{a}{b}(x)=z}{%
      \FormulaRefAuto{FreshExtMapSurjectiveOldValue}{1,2,3,11}}
    \proofstepwidestar[13]{z=b}{\rA}
    \proofstepwidestar[1,2,3,13]{%
      \exists x\in A\cup\{a\}\,
      \ExtMap{\iota_{B,B\cup\{b\}}\circ F}{a}{b}(x)=z}{%
      \FormulaRefAuto{FreshExtMapSurjectiveNewValue}{1,2,3,13}}
    \proofstepwidestar[1,2,3,9]{%
      \exists x\in A\cup\{a\}\,
      \ExtMap{\iota_{B,B\cup\{b\}}\circ F}{a}{b}(x)=z}{\rOE{10,11,12,13,14}}
    \proofstepwidestar[1,2,3]{%
      \begin{aligned}[t]
        z\in B\cup\{b\}\rightarrow{}&
        \exists x\in A\cup\{a\}\,\\[-2pt]
        &\ExtMap{\iota_{B,B\cup\{b\}}\circ F}{a}{b}(x)=z
      \end{aligned}}{\rRI{9,15}}
    \proofstepwidestar[1,2,3]{%
      \begin{aligned}[t]
        &\forall z\in B\cup\{b\}\,\exists x\in A\cup\{a\}\,\\[-2pt]
        &\qquad\ExtMap{\iota_{B,B\cup\{b\}}\circ F}{a}{b}(x)=z
      \end{aligned}}{\rUI{16}}
    \proofstepwidestar[1,2,3]{%
      \ExtMap{\iota_{B,B\cup\{b\}}\circ F}{a}{b}
      \colon A\cup\{a\}\sur B\cup\{b\}}{%
      \FormulaRefAuto{Surjektive Funktion}{8,17}}
  \closeproofpartwideindR

  \proofpartwide{Schluss}
    \proofstepwidestar[1]{F\colon A\bij B}{\rA}
    \proofstepwidestar[2]{a\notin A}{\rA}
    \proofstepwidestar[3]{b\notin B}{\rA}
    \proofstepwidestar[1,2,3]{%
      \ExtMap{\iota_{B,B\cup\{b\}}\circ F}{a}{b}
      \colon A\cup\{a\}\inj B\cup\{b\}}{%
      \FormulaRefAuto{FreshExtMapInjective}{1,2,3}}
    \proofstepwidestar[1,2,3]{%
      \ExtMap{\iota_{B,B\cup\{b\}}\circ F}{a}{b}
      \colon A\cup\{a\}\sur B\cup\{b\}}{%
      \FormulaRefAuto{FreshExtMapSurjective}{1,2,3}}
    \proofstepwidestar[1,2,3]{%
      \ExtMap{\iota_{B,B\cup\{b\}}\circ F}{a}{b}
      \colon A\cup\{a\}\bij B\cup\{b\}}{%
      \FormulaRefAuto{F\colon A\inj B\dsep F\colon A\sur B \vdash F\colon A\bij B}{4,5}}
  \closeproofpartwide
\end{tabproofsplitwide}


\section{Bijektionen mit vorgeschriebenem Wert}

\FormulaThmDeltaK[Vertauschen zweier Elemente]{%
  b\in B\dsep c\in B\vdash
  \exists\tau\,\bigl(\tau\colon B\bij B\land\tau(c)=b\bigr)%
}{TranspositionBijectionExists}{%
  \DeltaRow{Mengen}{B\dsep b\dsep c\dsep y\dsep z}
  \DeltaRow{Funktionen}{\tau}
}
Im Beweis ist \(t(y)\) die Abkürzung des bereits eingeführten Fallterms
\[
  t(y)\coloneqq\IfThenElse{y=c}{b}{\IfThenElse{y=b}{c}{y}}.
\]
\begin{tabproofsplitwide}
  \proofpartwideindR[Konstruktion der Fallabbildung]{%
    b\in B\dsep c\in B\vdash\exists\tau\,\bigl(\tau\colon B\to B\land\forall y\in B\,\tau(y)=t(y)\bigr)
  }{%
    b\in B\dsep c\in B\vdash\exists\tau\,\bigl(\tau\colon B\to B\land\forall y\in B\,\tau(y)=t(y)\bigr)
  }[TranspositionTermExists]
    \proofstepwidestar[1]{b\in B}{%
      \rA}
    \proofstepwidestar[2]{c\in B}{%
      \rA}
    \proofstepwidestar[3]{y\in B}{%
      \rA}
    \proofstepwidestar[]{y=c\lor y\neq c}{%
      \FormulaRefAuto{P\lor\neg P}[thm]}
    \proofstepwidestar[5]{y=c}{%
      \rA}
    \proofstepwidestar[5]{t(y)=b}{%
      \FormulaRefAuto{P\vdash\IfThenElse{P}{a}{b}=a}[thm]{5}}
    \proofstepwidestar[1,5]{t(y)\in B}{%
      \rIE{\rIE{6,\rII},1}}
    \proofstepwidestar[8]{y\neq c}{%
      \rA}
    \proofstepwidestar[8]{t(y)=\IfThenElse{y=b}{c}{y}}{%
      \FormulaRefAuto{\neg P\vdash\IfThenElse{P}{a}{b}=b}[thm]{8}}
    \proofstepwidestar[]{y=b\lor y\neq b}{%
      \FormulaRefAuto{P\lor\neg P}[thm]}
    \proofstepwidestar[11]{y=b}{%
      \rA}
    \proofstepwidestar[11]{\IfThenElse{y=b}{c}{y}=c}{%
      \FormulaRefAuto{P\vdash\IfThenElse{P}{a}{b}=a}[thm]{11}}
    \proofstepwidestar[2,8,11]{t(y)\in B}{%
      \rIE{\rIE{\rIE{12,9},\rII},2}}
    \proofstepwidestar[14]{y\neq b}{%
      \rA}
    \proofstepwidestar[14]{\IfThenElse{y=b}{c}{y}=y}{%
      \FormulaRefAuto{\neg P\vdash\IfThenElse{P}{a}{b}=b}[thm]{14}}
    \proofstepwidestar[3,8,14]{t(y)\in B}{%
      \rIE{\rIE{\rIE{15,9},\rII},3}}
    \proofstepwidestar[2,3,8]{t(y)\in B}{%
      \rOE{10,11,13,14,16}}
    \proofstepwidestar[1,2,3]{t(y)\in B}{%
      \rOE{4,5,7,8,17}}
    \proofstepwidestar[1,2]{\forall y\in B\,t(y)\in B}{%
      \rUI{\rRI{3,18}}}
    \proofstepwidestar[1,2]{\exists\tau\,\bigl(\tau\colon B\to B\land\forall y\in B\,\tau(y)=t(y)\bigr)}{%
      \FormulaRefAuto{TypedFunctionDefinitionPrinciple}[thm]{19}}
  \closeproofpartwideindR

  \proofpartwideindR[Werte der Fallabbildung]{%
    \begin{aligned}[t]
      &b\in B\dsep c\in B\dsep\forall y\in B\,\tau(y)=t(y)\vdash{}\\[-2pt]
      &\quad\tau(c)=b\land\tau(b)=c\land{}\\[-2pt]
      &\quad\forall y\in B\,(y\neq c\land y\neq b\rightarrow\tau(y)=y)
    \end{aligned}
  }{%
    b\in B\dsep c\in B\dsep\forall y\in B\,\tau(y)=t(y)\vdash\tau(c)=b\land\tau(b)=c\land\forall y\in B\,(y\neq c\land y\neq b\rightarrow\tau(y)=y)
  }[TranspositionTermValues]
    \proofstepwidestar[1]{b\in B}{%
      \rA}
    \proofstepwidestar[2]{c\in B}{%
      \rA}
    \proofstepwidestar[3]{\forall y\in B\,\tau(y)=t(y)}{%
      \rA}
    \proofstepwidestar[]{t(c)=b}{%
      \FormulaRefAuto{P\vdash\IfThenElse{P}{a}{b}=a}[thm]{\rII}}
    \proofstepwidestar[2,3]{\tau(c)=t(c)}{%
      \rRE{\rUE{3},2}}
    \proofstepwidestar[2,3]{\tau(c)=b}{%
      \rIE{4,5}}
    \proofstepwidestar[1,3]{\tau(b)=t(b)}{%
      \rRE{\rUE{3},1}}
    \proofstepwidestar[]{b=c\lor b\neq c}{%
      \FormulaRefAuto{P\lor\neg P}[thm]}
    \proofstepwidestar[9]{b=c}{%
      \rA}
    \proofstepwidestar[9]{t(b)=b}{%
      \FormulaRefAuto{P\vdash\IfThenElse{P}{a}{b}=a}[thm]{9}}
    \proofstepwidestar[1,3,9]{\tau(b)=c}{%
      \rIE{10,9,7}}
    \proofstepwidestar[12]{b\neq c}{%
      \rA}
    \proofstepwidestar[12]{t(b)=\IfThenElse{b=b}{c}{b}}{%
      \FormulaRefAuto{\neg P\vdash\IfThenElse{P}{a}{b}=b}[thm]{12}}
    \proofstepwidestar[]{\IfThenElse{b=b}{c}{b}=c}{%
      \FormulaRefAuto{P\vdash\IfThenElse{P}{a}{b}=a}[thm]{\rII}}
    \proofstepwidestar[1,3,12]{\tau(b)=c}{%
      \rIE{13,14,7}}
    \proofstepwidestar[1,3]{\tau(b)=c}{%
      \rOE{8,9,11,12,15}}
    \proofstepwidestar[17]{y\in B}{%
      \rA}
    \proofstepwidestar[18]{y\neq c\land y\neq b}{%
      \rA}
    \proofstepwidestar[18]{t(y)=\IfThenElse{y=b}{c}{y}}{%
      \FormulaRefAuto{\neg P\vdash\IfThenElse{P}{a}{b}=b}[thm]{\rAEa{18}}}
    \proofstepwidestar[18]{\IfThenElse{y=b}{c}{y}=y}{%
      \FormulaRefAuto{\neg P\vdash\IfThenElse{P}{a}{b}=b}[thm]{\rAEb{18}}}
    \proofstepwidestar[3,17]{\tau(y)=t(y)}{%
      \rRE{\rUE{3},17}}
    \proofstepwidestar[3,17,18]{\tau(y)=y}{%
      \rIE{19,20,21}}
    \proofstepwidestar[3]{\forall y\in B\,(y\neq c\land y\neq b\rightarrow\tau(y)=y)}{%
      \rUI{\rRI{17,\rRI{18,22}}}}
    \proofstepwidestar[1,2,3]{\tau(c)=b\land\tau(b)=c\land\forall y\in B\,(y\neq c\land y\neq b\rightarrow\tau(y)=y)}{%
      \rAI{6,\rAI{16,23}}}
  \closeproofpartwideindR

  \proofpartwideindR[Die Fallabbildung ist bijektiv]{%
    b\in B\dsep c\in B\dsep\tau\colon B\to B\dsep\forall y\in B\,\tau(y)=t(y)\vdash\tau\colon B\bij B\land\tau(c)=b
  }{%
    b\in B\dsep c\in B\dsep\tau\colon B\to B\dsep\forall y\in B\,\tau(y)=t(y)\vdash\tau\colon B\bij B\land\tau(c)=b
  }[TranspositionTermBijective]
    \proofstepwidestar[1]{b\in B}{%
      \rA}
    \proofstepwidestar[2]{c\in B}{%
      \rA}
    \proofstepwidestar[3]{\tau\colon B\to B}{%
      \rA}
    \proofstepwidestar[4]{\forall y\in B\,\tau(y)=t(y)}{%
      \rA}
    \proofstepwidestar[1,2,4]{\tau(c)=b\land\tau(b)=c\land\forall y\in B\,(y\neq c\land y\neq b\rightarrow\tau(y)=y)}{%
      \FormulaRefAuto{TranspositionTermValues}[thm]{1,2,4}}
    \proofstepwidestar[1,2,4]{\tau(c)=b}{%
      \rAEa{5}}
    \proofstepwidestar[1,2,4]{\tau(b)=c}{%
      \rAEa{\rAEb{5}}}
    \proofstepwidestar[1,2,4]{\forall y\in B\,(y\neq c\land y\neq b\rightarrow\tau(y)=y)}{%
      \rAEb{\rAEb{5}}}
    \proofstepwidestar[9]{y\in B}{%
      \rA}
    \proofstepwidestar[]{y=c\lor y\neq c}{%
      \FormulaRefAuto{P\lor\neg P}[thm]}
    \proofstepwidestar[11]{y=c}{%
      \rA}
    \proofstepwidestar[1,2,4,11]{\tau(\tau(y))=y}{%
      \rIE{\rIE{6,\rII},\rIE{11,\rII},7}}
    \proofstepwidestar[13]{y\neq c}{%
      \rA}
    \proofstepwidestar[]{y=b\lor y\neq b}{%
      \FormulaRefAuto{P\lor\neg P}[thm]}
    \proofstepwidestar[15]{y=b}{%
      \rA}
    \proofstepwidestar[1,2,4,15]{\tau(\tau(y))=y}{%
      \rIE{\rIE{7,\rII},\rIE{15,\rII},6}}
    \proofstepwidestar[17]{y\neq b}{%
      \rA}
    \proofstepwidestar[1,2,4,9,13,17]{\tau(y)=y}{%
      \rRE{\rRE{\rUE{8},9},\rAI{13,17}}}
    \proofstepwidestar[1,2,4,9,13,17]{\tau(\tau(y))=y}{%
      \rIE{\rIE{18,\rII},18}}
    \proofstepwidestar[1,2,4,9,13]{\tau(\tau(y))=y}{%
      \rOE{14,15,16,17,19}}
    \proofstepwidestar[1,2,4,9]{\tau(\tau(y))=y}{%
      \rOE{10,11,12,13,20}}
    \proofstepwidestar[1,2,4]{\forall y\in B\,\tau(\tau(y))=y}{%
      \rUI{\rRI{9,21}}}
    \proofstepwidestar[23]{z\in B}{%
      \rA}
    \proofstepwidestar[24]{\tau(y)=\tau(z)}{%
      \rA}
    \proofstepwidestar[1,2,4,23]{\tau(\tau(z))=z}{%
      \rRE{\rUE{22},23}}
    \proofstepwidestar[1,2,4,23,24]{\tau(\tau(y))=z}{%
      \rIE{\rIE{24,\rII},25}}
    \proofstepwidestar[1,2,4,9,23,24]{y=z}{%
      \rIE{21,26}}
    \proofstepwidestar[1,2,4]{\forall y\in B\,\forall z\in B\,(\tau(y)=\tau(z)\rightarrow y=z)}{%
      \rUI{\rRI{9,\rUI{\rRI{23,\rRI{24,27}}}}}}
    \proofstepwidestar[1,2,3,4]{\tau\colon B\inj B}{%
      \FormulaRefAuto{Injektive Funktion}[def]{3,28}}
    \proofstepwidestar[3,9]{\tau(y)\in B}{%
      \FormulaRefAuto{F\colon A\to B\dsep x\in A\vdash F(x)\in B}[thm]{3,9}}
    \proofstepwidestar[1,2,3,4,9]{\exists z\in B\,\tau(z)=y}{%
      \rEI{\rAI{30,21}}}
    \proofstepwidestar[1,2,3,4]{\forall y\in B\,\exists z\in B\,\tau(z)=y}{%
      \rUI{\rRI{9,31}}}
    \proofstepwidestar[1,2,3,4]{\tau\colon B\sur B}{%
      \FormulaRefAuto{Surjektive Funktion}[def]{3,32}}
    \proofstepwidestar[1,2,3,4]{\tau\colon B\bij B}{%
      \FormulaRefAuto{F\colon A\inj B\dsep F\colon A\sur B\vdash F\colon A\bij B}[thm]{29,33}}
    \proofstepwidestar[1,2,3,4]{\tau\colon B\bij B\land\tau(c)=b}{%
      \rAI{34,6}}
  \closeproofpartwideindR

  \proofpartwide{Existenzschluss}
    \proofstepwidestar[1]{b\in B}{%
      \rA}
    \proofstepwidestar[2]{c\in B}{%
      \rA}
    \proofstepwidestar[1,2]{\exists\tau\,\bigl(\tau\colon B\to B\land\forall y\in B\,\tau(y)=t(y)\bigr)}{%
      \FormulaRefAuto{TranspositionTermExists}[thm]{1,2}}
    \proofstepwidestar[4]{\tau\colon B\to B\land\forall y\in B\,\tau(y)=t(y)}{%
      \rA}
    \proofstepwidestar[1,2,4]{\tau\colon B\bij B\land\tau(c)=b}{%
      \FormulaRefAuto{TranspositionTermBijective}[thm]{1,2,\rAEa{4},\rAEb{4}}}
    \proofstepwidestar[1,2,4]{\exists\tau\,\bigl(\tau\colon B\bij B\land\tau(c)=b\bigr)}{%
      \rEI{5}}
    \proofstepwidestar[1,2]{\exists\tau\,\bigl(\tau\colon B\bij B\land\tau(c)=b\bigr)}{%
      \rEE{3,4,6}}
  \closeproofpartwide
\end{tabproofsplitwide}

\FormulaThmDeltaK[Bijektion mit vorgeschriebenem Wert]{%
  \exists u\,(u\colon A\bij B)\dsep a\in A\dsep b\in B\vdash
  \exists f\,\bigl(f\colon A\bij B\land f(a)=b\bigr)%
}{PointedBijectionExists}{%
  \DeltaRow{Mengen}{A\dsep B\dsep a\dsep b}
  \DeltaRow{Funktionen}{u\dsep\tau\dsep f}
}
\begin{tabproofwide}
  \proofstepwidestar[1]{\exists u\,(u\colon A\bij B)}{\rA}
  \proofstepwidestar[2]{a\in A}{\rA}
  \proofstepwidestar[3]{b\in B}{\rA}
  \proofstepwidestar[4]{u\colon A\bij B}{\rA}
  \proofstepwidestar[2,4]{u(a)\in B}{%
    \FormulaRefAuto{F\colon A\to B\dsep x\in A\vdash F(x)\in B}[thm]{4,2}}
  \proofstepwidestar[2,3,4]{%
    \exists\tau\,\bigl(\tau\colon B\bij B\land\tau(u(a))=b\bigr)}{%
    \FormulaRefAuto{TranspositionBijectionExists}[thm]{3,5}}
  \proofstepwidestar[7]{\tau\colon B\bij B\land\tau(u(a))=b}{\rA}
  \proofstepwidestar[4,7]{\tau\circ u\colon A\bij B}{%
    \FormulaRefAuto{F\colon A\bij B\dsep G\colon B\bij C\vdash
      G\circ F\colon A\bij C}[thm]{4,\rAEa{7}}}
  \proofstepwidestar[2,4,7]{(\tau\circ u)(a)=\tau(u(a))=b}{%
    \FormulaRefAuto{CompositionDef}[def]{4,7,2}; \rAEb{7}}
  \proofstepwidestar[2,4,7]{%
    \exists f\,\bigl(f\colon A\bij B\land f(a)=b\bigr)}{\rEI{\rAI{8,9}}}
  \proofstepwidestar[2,3,4]{%
    \exists f\,\bigl(f\colon A\bij B\land f(a)=b\bigr)}{\rEE{6,7,10}}
  \proofstepwidestar[1,2,3]{%
    \exists f\,\bigl(f\colon A\bij B\land f(a)=b\bigr)}{\rEE{1,4,11}}
\end{tabproofwide}

\section{Leere Bereiche}

\FormulaThmDelta[Leere Relation ist bijektiv auf $\varnothing$]{%
  \varnothing\colon \varnothing\bij \varnothing%
}{%
  \DeltaDecl{Mengen}{}
}
\begin{tabproof}
  \proofstep{}{ \varnothing\colon \varnothing\inj \varnothing }{%
    \FormulaRefAuto{\varnothing\colon \varnothing\inj M}}

  \proofstep{}{ \varnothing\colon \varnothing\sur \varnothing }{%
    \FormulaRefAuto{\varnothing\colon \varnothing\sur \varnothing}}

  \proofstep{}{ \varnothing\colon \varnothing\bij \varnothing }{%
    \FormulaRefAuto{F\colon A\inj B\dsep F\colon A\sur B \vdash F\colon A\bij B}{1,2}}
\end{tabproof}

\end{document}
